\documentclass[11pt]{article}
\usepackage[letterpaper,margin=0.7in,headheight=15pt]{geometry}


\usepackage{microtype,amsmath,amssymb,mathtools}
\usepackage{graphicx,xcolor,booktabs,array,enumitem,fancyhdr}
\usepackage[hidelinks]{hyperref}
\usepackage{bookmark}
\definecolor{Ink}{HTML}{20303A}
\definecolor{Accent}{HTML}{226B66}
\definecolor{Muted}{HTML}{637680}
\definecolor{Pale}{HTML}{EAF2F1}
\color{Ink}
\setlength{\parindent}{0pt}
\setlength{\parskip}{0.32em}
\setlength{\emergencystretch}{2em}
\raggedbottom
\pagestyle{fancy}
\fancyhf{}
\fancyhead[L]{\small\color{Muted}Realization Calculus}
\fancyhead[R]{\small\color{Muted}Special Functions / Visual guide}
\fancyfoot[C]{\small\thepage}
\renewcommand{\headrulewidth}{0.3pt}
\newcounter{lesson}
\newcommand{\pendingvolume}{}
\newcommand{\startvolume}[2]{\gdef\pendingvolume{#1: #2}}
\newcommand{\lesson}[8]{%
  \clearpage
  \stepcounter{lesson}%
  \phantomsection
  \ifx\pendingvolume\empty\else
    \addcontentsline{toc}{section}{\pendingvolume}%
    \global\let\pendingvolume\empty
  \fi
  \phantomsection
  \addcontentsline{toc}{subsection}{\protect\numberline{\thelesson}#2}%
  {\small\sffamily\bfseries\color{Accent}#1\hfill LESSON \thelesson\par}%
  \vspace{0.16in}
  {\LARGE\rmfamily\color{RCPlum}#2\par}%
  \vspace{0.08in}
  \begin{center}
    \includegraphics[width=\linewidth,height=3.65in,keepaspectratio]{figures/#3.pdf}
  \end{center}
  \vspace{-0.03in}
  {\color{Accent}\[#4\]}%
  {\bfseries Read the figure.} #5\par
  \vspace{0.05in}
  {\bfseries Work an example.} #6\par
  \vspace{0.05in}
  \noindent\colorbox{Pale}{\parbox{\dimexpr\linewidth-2\fboxsep\relax}{\small\textbf{Carry forward.} #7}}\par
  \vfill
  {\footnotesize\color{Muted}\textbf{Source and scope.} #8\par}%
}

% One bookmark per lesson and one per thematic group.
\makeatletter
\def\toclevel@section{0}
\def\toclevel@subsection{1}
\makeatother

\hypersetup{pdftitle={Special Functions: Expanded Visual Guide},pdfauthor={Komorebi P.},pdfsubject={171 illustrated lessons across Projector Generated Special Functions I--VIII}}

\newcommand{\AtlasRunningTitle}{SPECIAL FUNCTIONS / VISUAL GUIDE}
% Realization Calculus companion design. Fonts are system dependencies.
\usepackage{fontspec,mathpazo,tikz,titlesec}
\usetikzlibrary{calc}
\IfFileExists{/usr/share/fonts/opentype/urw-base35/P052-Roman.otf}{
 \setmainfont{P052-Roman.otf}[Path=/usr/share/fonts/opentype/urw-base35/,BoldFont=P052-Bold.otf,ItalicFont=P052-Italic.otf,BoldItalicFont=P052-BoldItalic.otf,Ligatures=TeX]
}{\setmainfont{TeX Gyre Pagella}[Ligatures=TeX]}
\IfFileExists{/usr/share/fonts/opentype/urw-base35/NimbusSans-Regular.otf}{
 \setsansfont{NimbusSans-Regular.otf}[Path=/usr/share/fonts/opentype/urw-base35/,BoldFont=NimbusSans-Bold.otf,ItalicFont=NimbusSans-Italic.otf,BoldItalicFont=NimbusSans-BoldItalic.otf,Scale=.94]
}{\setsansfont{TeX Gyre Heros}[Scale=.94]}
\IfFontExistsTF{DejaVu Sans Mono}{\setmonofont{DejaVu Sans Mono}[Scale=.82]}{\setmonofont{Latin Modern Mono}[Scale=.86]}
\definecolor{RCPlum}{HTML}{39263F}
\definecolor{RCGold}{HTML}{A78A50}
\definecolor{RCMuted}{HTML}{68616C}
\definecolor{RCPale}{HTML}{DDD6DF}
\colorlet{Ink}{black!85}
\colorlet{Accent}{RCPlum}
\colorlet{Muted}{RCMuted}
\colorlet{Pale}{RCPlum!5}
\colorlet{ink}{Ink}\colorlet{accent}{RCPlum}\colorlet{warm}{RCGold}
\color{Ink}
\pagestyle{fancy}\fancyhf{}
\fancyhead[L]{\sffamily\addfontfeatures{LetterSpace=5}\fontsize{7.1}{10}\selectfont\color{Muted}REALIZATION CALCULUS}
\fancyhead[R]{\sffamily\fontsize{7.1}{10}\selectfont\color{Muted}\AtlasRunningTitle}
\fancyfoot[R]{\rmfamily\fontsize{10}{12}\selectfont\thepage}
\renewcommand{\headrulewidth}{0pt}
\fancypagestyle{plain}{\fancyhf{}\fancyfoot[R]{\rmfamily\fontsize{10}{12}\selectfont\thepage}\renewcommand{\headrulewidth}{0pt}}
\titleformat{\section}{\normalfont\Large\color{RCPlum}}{\thesection}{.7em}{}
\titleformat{\subsection}{\normalfont\large\color{RCPlum}}{\thesubsection}{.7em}{}
\hypersetup{colorlinks=true,linkcolor=RCPlum,citecolor=RCPlum,urlcolor=RCPlum,pdfauthor={Komorebi P.},pdfdisplaydoctitle=true,bookmarksopen=true,bookmarksopenlevel=0}
\newcommand{\AtlasCover}[3]{%
 \hypersetup{pageanchor=false}\thispagestyle{empty}\begingroup
 \begin{tikzpicture}[remember picture,overlay]
 \fill[RCPlum] (current page.south west) rectangle (current page.north east);
 \node[anchor=north west,text=white,font=\sffamily\addfontfeatures{LetterSpace=16}\fontsize{9}{12}\selectfont] at ([xshift=.9in,yshift=-1in]current page.north west) {REALIZATION CALCULUS};
 \draw[RCGold,line width=.8pt] ([xshift=.92in,yshift=-1.52in]current page.north west) -- ++(.56in,0);
 \node[anchor=north west,text=white,align=left,inner sep=0pt,font=\rmfamily\fontsize{52}{60}\selectfont] at ([xshift=.88in,yshift=-1.92in]current page.north west) {#1};
 \node[anchor=north west,text=white,align=left,inner sep=0pt,font=\rmfamily\fontsize{23}{29}\selectfont] at ([xshift=.92in,yshift=-4.32in]current page.north west) {#2};
 \node[anchor=north west,text=RCPale,align=left,inner sep=0pt,font=\rmfamily\fontsize{13}{19}\selectfont] at ([xshift=.92in,yshift=-5.53in]current page.north west) {#3};
 \draw[white!25!RCPlum,line width=.45pt] ([xshift=.92in,yshift=1.29in]current page.south west) -- ([xshift=-.88in,yshift=1.29in]current page.south east);
 \node[anchor=south west,text=white,inner sep=0pt,font=\rmfamily\fontsize{14}{18}\selectfont] at ([xshift=.92in,yshift=.82in]current page.south west) {Komorebi P.};
 \end{tikzpicture}\null\endgroup\clearpage\hypersetup{pageanchor=true}
}

\hypersetup{pdftitle={Special Functions: Expanded Visual Guide},pdfsubject={171 worked lessons across Projector Generated Special Functions I--VIII}}
\begin{document}
\AtlasCover{Special\\Functions}{Expanded visual guide}{Packets, sequences, polyphase channels\\and exact mathematical examples}

\hypersetup{pageanchor=false}

\pagenumbering{roman}
\hypersetup{pageanchor=true}
\pdfbookmark[0]{How to use this guide}{reading-guide}
{\LARGE\bfseries How to use this guide\par}
\vspace{0.2in}
This guide develops the pictures into a sequence of worked mathematical lessons. It accompanies the consolidated manuscript and the eight original papers. Each lesson has one central figure, a formula, directions for reading the figure, a worked example, and a source-and-scope note.

The figures are produced by Python and Matplotlib from explicit formulas, finite matrices, declared quadrature, or bounded arithmetic sums. Vector diagrams describe exact relationships; numerical curves are labeled with their parameter ranges and approximation limits. Original source statements are interpreted through the correction register in the consolidated manuscript.

\vspace{0.1in}
\begin{tabular}{>{\raggedright\arraybackslash}p{1.25in}>{\raggedright\arraybackslash}p{5.1in}}
\toprule
Route & What to follow\\
\midrule
Foundations & Coefficient addresses, refinement, carry, detail banks, truncation, interference, moments, and parameter domains.\\[0.09in]
Papers I--III & Analytic packet reconstruction, Bernoulli--Hurwitz refinement, arithmetic depth, positive measures, Gamma limits, and spectral reduction.\\[0.09in]
Papers IV--VI & Time-average Gram lifts, signed rational models, log-Gamma windows, Fourier geometry, and continuous shape refinement.\\[0.09in]
Papers VII--VIII & Finite lowest-weight spaces, Casimir atlases, signed recoupling amplitudes, measurement channels, tree consistency, and q-deformation.\\[0.09in]
Synthesis & The precise bridges between the three kinds of packets, with examples that retain the hypotheses.\\[0.09in]
Sequence and signal atlases & OEIS families and transforms, exact polyphase operations, finite taps, filter banks, and streaming execution.\\[0.09in]
Bridges and capstone & Seeds, pole cancellation, operator contracts, complete residual calculations, information loss, model-based repair, and the streaming phase clock.\\
\bottomrule
\end{tabular}\par

\vspace{0.18in}
\textbf{A first pass.} Read the eight foundations, then the first three lessons of each paper. Return to the later lessons for the asymptotic, spectral, and representation-theoretic developments. The contents list every lesson so the guide also works as a reference.

\textbf{Evidence scope.} A plot illustrates a statement within its declared range. It does not prove an unbounded claim. The signed arithmetic obstruction remains open. The scalar-to-matrix q-Racah comparison remains a separate normalization task. Full analytic Lean formalization is not established by the available source evidence. These distinctions are taught at the points where they matter.

\clearpage
\pdfbookmark[0]{Lesson contents}{guide-contents}
{\small\tableofcontents}
\clearpage
\pagenumbering{arabic}
\startvolume{Foundations}{Read the packet language}
\lesson{FOUNDATIONS / 01}{An exponent becomes an address}{F01_coefficient_addresses}
{n=Mm+r,\qquad F_{M,r}(u)=\sum_{m\ge0}a_{Mm+r}u^m}
{Read each row as a residue class and each column as its internal index. The figure lists the first eighteen exponents at modulus three. Every number occurs once. The row selects a packet; the column tells that packet where the coefficient belongs.}
{For exponent $14$, divide by three: $14=3\cdot4+2$. Its coefficient is the coefficient of $u^4$ in $F_{3,2}$. Multiplying this term by $z^2$ and substituting $u=z^3$ restores $a_{14}z^{14}$. No coefficient has been averaged or discarded.}
{A packet address is a reversible coordinate for a coefficient. It is useful before any special-function identification is made.}
{Paper I, ``Residue-section functions'' and ``Exact polyphase reconstruction.'' The displayed address table is exact integer arithmetic.}

\lesson{FOUNDATIONS / 02}{Refinement preserves every address}{F02_refinement_addresses}
{F_{M,r}(u)=\sum_{\ell=0}^{L-1}u^\ell F_{LM,r+M\ell}(u^L)}
{Start at either parent on the left and follow its three arrows. Refining modulus two by three produces modulus six. The children of the even parent are addresses $0,2,4$; those of the odd parent are $1,3,5$. The factor $u^\ell$ keeps the child offset.}
{Take $M=2$, $r=1$, and parent index $m=7$. Split $7=3\cdot2+1$. The original exponent is $2\cdot7+1=15$, and its refined address is $6\cdot2+3=15$. The coefficient passes to child $r+M\ell=3$ at internal index two.}
{Refinement reorganizes the same coefficients. Parent and child indices must be tracked together.}
{Paper I, ``Parent-child refinement'' and the coefficient refinement identity in the consolidated manuscript. This diagram uses $M=2$, $L=3$.}

\lesson{FOUNDATIONS / 03}{Multiplication carries between packets}{F03_carry_matrix}
{C_3(u)=\begin{pmatrix}0&0&u\\1&0&0\\0&1&0\end{pmatrix},\qquad \mathcal S_3(C_3 f)=z\mathcal S_3(f)}
{Read matrix columns as input addresses and rows as output addresses. Multiplication by $z$ advances addresses zero to one and one to two. Address two wraps to zero, bringing one factor of $u=z^3$. The right panel evaluates this carry at $u=0.4$.}
{For the constant packet vector $(1,2,3)^T$, synthesis gives $1+2z+3z^2$. The carried vector is $(3u,1,2)^T$, whose synthesis is $3z^3+z+2z^2$. This is exactly $z(1+2z+3z^2)$. At $u=2/5$, its entries are $(6/5,1,2)$.}
{A scalar operation can become a small matrix operation, with the wrap recorded explicitly.}
{Paper I, ``Intertwining laws.'' The matrix and polynomial identity are exact; the bars evaluate one rational parameter.}

\lesson{FOUNDATIONS / 04}{A detail bank separates average and variation}{F04_unitary_detail_bank}
{W_j=L^{-1/2}\sum_{\ell=0}^{L-1}e^{-2\pi i j\ell/L}b_\ell,\qquad \sum_j|W_j|^2=\sum_\ell|b_\ell|^2}
{The left panel has four child amplitudes. The right shows their unitary discrete Fourier channel energies. Channel zero is the aggregate; the other channels encode variation among children. Because the transform is unitary, energy is preserved, even though the original real amplitudes can become complex.}
{For $b=(1,2,4,1)$, the channels are $4$, $(-3-i)/2$, $1$, and $(-3+i)/2$. Their energies are $16$, $5/2$, $1$, and $5/2$, totaling $22$. The inverse transform restores all four children. Keeping only channel zero would discard the details.}
{A reversible detail bank needs all channels. Aggregation alone is a projection.}
{Paper I, the unitary child bank. This four-child example uses exact integer inputs and algebraic DFT values; numerical Parseval is also checked.}

\lesson{FOUNDATIONS / 05}{Separate routing error from truncation error}{F05_identity_and_truncation}
{P_{12}(z)=\sum_{n=0}^{12}\frac{z^n}{n!},\qquad |e^z-P_{12}(z)|\le e^2\frac{2^{13}}{13!}\quad(0\le z\le2)}
{The first panel overlays the exponential and its degree-twelve truncation. The second enlarges the small remainder and its Taylor bound. Reconstructing the finite polynomial from three packets introduces no mathematical routing error. Approximating the infinite exponential by that polynomial does introduce a remainder.}
{The three finite packets use exactly the coefficients with $3m+r\le12$. Their synthesized sum is $P_{12}$. At real $z=2$, Taylor's theorem bounds the omitted tail by $e^2 2^{13}/13!$, about $9.72\times10^{-6}$. The figure plots the actual remainder in millionths.}
{An exact coordinate identity can coexist with a numerical approximation. Report the approximation at the stage where it enters.}
{Paper I reconstruction, illustrated with $F=e^z$. Degree twelve, real domain $[0,2]$; the remainder bound is the elementary Taylor bound.}

\lesson{FOUNDATIONS / 06}{Positive energy does not retain interference}{F06_phase_and_power}
{\left|\frac45+\frac35e^{i\phi}\right|^2=1+\frac{24}{25}\cos\phi}
{The solid curve is coherent power as the relative phase changes. The dashed line keeps only the sum of the two component energies. The latter is constant, while the coherent value can be almost canceled or almost doubled. Removing cross terms changes the information being represented.}
{Both component energies sum to $(4/5)^2+(3/5)^2=1$. At phase zero the coherent power is $49/25$. At phase $\pi$ it is $1/25$. Averaging uniformly over a full phase period gives one because the cosine term has zero mean.}
{A positive diagonal lift is valuable, but a bound on it alone does not recover a signed or coherent cancellation estimate.}
{An exact two-component teaching model for the distinction used in Papers I, IV, VII, and VIII. It is not an arithmetic cancellation theorem.}

\lesson{FOUNDATIONS / 07}{One measure creates a sequence of moments}{F07_measure_and_moments}
{E_k=\sum_jw_jx_j^k,\qquad c^THc=\sum_jw_j\left(\sum_rc_rx_j^r\right)^2\ge0}
{The left panel shows three atoms and their positive weights. The right plots nine moments of this same measure. Higher powers favor atoms nearer one. The zero atom contributes to $E_0$ and disappears from positive moments, so discarding it would change the starting value.}
{Use nodes $(0,2/5,17/20)$ and weights $(1/5,3/10,1/2)$. Then $E_0=1$, $E_1=109/200$, and $E_2=1637/4000$. The three-by-three Hankel matrix is positive definite because a nonzero quadratic cannot vanish at all three distinct atoms. The four-by-four matrix is singular because this toy support has only three points.}
{Strict positivity at every size requires enough support. The arithmetic source measure has infinitely many distinct atoms.}
{Papers II--III positive-moment construction. This explicitly finite teaching measure has rank three; its moments and determinant are computed with rational arithmetic.}

\lesson{FOUNDATIONS / 08}{Parameters have different mathematical jobs}{F08_parameter_domains}
{\alpha=\frac{2a-1}{\beta}>0,\qquad \kappa=\frac\alpha2}
{The curves convert arithmetic decay $a$ and scale exponent $\beta$ into a positive Gamma shape. The shaded strip marks $1/2<a\le1$: positive energy is defined there, while the absolutely convergent signed transfer uses $a>1$. The table continues each shape to its ordinary ladder label.}
{For $a=2$ and $\beta=2$, the shape is $\alpha=3/2$ and the ladder parameter is $\kappa=3/4$. This parameter choice differs from the half-shape recoupling example $\kappa=1/2$. Also distinguish an integer character modulus from the real deformation parameter $0<q\le1$ used in Paper VIII.}
{Follow the declared parameter domain whenever a formula crosses from arithmetic packets to continuous shapes or finite recoupling.}
{Consolidated Chapters 1--3; Papers II, V, and VII. The table is exact rational arithmetic, and the curves show the stated parameter map.}

\startvolume{Paper I}{Exact lifts, Bernoulli refinement, and signed transfer}

\lesson{PAPER I / 01}{A scalar germ has exact coefficient addresses}{I01_coefficient_addresses}
{F(z)=\sum_{r=0}^{M-1}z^rF_{M,r}(z^M),\qquad F_{M,r}(u)=\sum_{m\ge0}a_{Mm+r}u^m.}
{Follow each colored row horizontally. The label on a dot is the original Taylor exponent, while its row is the residue address. The three rows cover every exponent once. Moving from the scalar function to packets changes coordinates and preserves all coefficients; synthesis restores each packet's residue factor. This diagram shows addresses, not the magnitudes of the coefficients.}
{For $e^z$ and $M=3$, the first packets begin $E_{3,0}=1+u/6+\cdots$, $E_{3,1}=1+u/24+\cdots$, and $E_{3,2}=1/2+u/120+\cdots$. The term $u/120$ in address 2 reconstructs to $z^2z^3/120=z^5/5!$.}
{An exact packet is a coefficient address together with its synthesis rule.}
{Volume I, Chapter 1, Definition 1.1 and Theorem 1.2, ``Exact polyphase reconstruction.'' Analytic germs at zero; the series identities hold in their convergence domains.}

\lesson{PAPER I / 02}{Multiplication creates a carry}{I02_carry_matrix}
{C_3(u)=\begin{pmatrix}0&0&u\\1&0&0\\0&1&0\end{pmatrix},\qquad \mathcal S_3(C_3\boldsymbol f)=z\mathcal S_3\boldsymbol f.}
{Read a matrix column as an input address and a row as an output address. Multiplication advances addresses 0 and 1 by one position. Address 2 wraps to address 0 and acquires a factor $u$. The left panel evaluates that matrix at $u=1/2$ only to make its entries visible; the algebraic matrix retains the variable $u$.}
{The scalar term $z^2u^m$ represents exponent $3m+2$. After multiplication it becomes $z^3u^m=u^{m+1}$ in address 0. The factor $u$ at the matrix's upper-right corner records exactly this carry. Without it, the exponent and coefficient reconstruction would both fail.}
{Carry is part of multiplication, so it must be transported with the packets.}
{Volume I, Chapter 2, Proposition 2.1, ``Intertwining laws.'' The matrix acts on analytic packet germs. The displayed numerical heatmap is an illustrative evaluation, not a fitted operator.}

\lesson{PAPER I / 03}{The Euler operator remembers the full exponent}{I03_euler_addresses}
{\Theta_M=Mu\frac{d}{du}I_M+\operatorname{diag}(0,\ldots,M-1),\qquad \Theta_M(u^me_r)=(Mm+r)u^me_r.}
{Choose a column to fix the packet degree $m$, then move between rows to change address $r$. Each cell returns the original scalar exponent $3m+r$. Differentiating only in $u$ would recover $3m$ and lose the address contribution. The diagonal term restores it. The table is exact integer arithmetic.}
{For $m=2$ in address $r=1$, synthesis gives $z^1(z^3)^2=z^7$. The operator $3u\,d/du+1$ multiplies that packet monomial by $3\cdot2+1=7$, just as $z\,d/dz$ multiplies $z^7$ by 7. Apply multiplication once more and the eigenvalue rises to 8.}
{The transported Euler and carry operators satisfy $[\Theta_M,C_M]=C_M$.}
{Volume I, Chapter 2, Proposition 2.1 and Corollary 2.2, ``Weyl commutator.'' Addresses and packet degrees are nonnegative integers; operator ordering is retained.}

\lesson{PAPER I / 04}{A differential equation survives the lift}{I04_exponential_system}
{\Theta_3\boldsymbol E_3=C_3\boldsymbol E_3,\qquad 3uE_{3,0}'(u)=uE_{3,2}(u).}
{The left panel plots 22-term samples of the three exponential sections. The right panel compares the two sides of the first lifted row. Their agreement expresses a coefficient identity inherited from $\theta e^z=ze^z$. The scalar differential equation becomes a coupled matrix system, with the same ordered operator products.}
{For $m\ge1$, a term $u^m/(3m)!$ in $E_{3,0}$ contributes $3m\,u^m/(3m)!=u^m/(3m-1)!$. This is the term obtained from $uE_{3,2}$ by using packet degree $m-1$. At $u=0$ both sides of this row vanish.}
{Transport an equation by intertwining its ordered operators, rather than approximating its solutions.}
{Volume I, Chapter 2, Theorem 2.3, ``Differential lift,'' and Chapter 3, exponential example. On $0\le u\le8$, each 22-term packet tail is below $10^{-60}$; recorded bounds also include the unmatched last finite row term. Floating-point rounding is separate.}

\lesson{PAPER I / 05}{Integer Bernoulli values sit inside an entire family}{I05_bernoulli_interpolation}
{\mathfrak B(z;a)=-z\zeta(1-z,a),\quad \mathfrak B(0;a)=1,\quad \mathfrak B(n;a)=B_n(a)\ (n\ge1).}
{The curve varies the order continuously while keeping $a=1/2$. The dots mark positive-integer Bernoulli-polynomial samples and the removable value at zero. A family indexed by complex order carries more local information than its integer samples. The order-zero value is defined by removal of the zeta pole, not by evaluating an undefined product directly.}
{At order four, $B_4(1/2)=7/240$. At every odd integer $n\ge3$, reflection symmetry gives $B_n(1/2)=0$. These isolated samples fit inside the continuous curve. The order-one value is also zero at this translation parameter.}
{Keep the translation parameter and order variable distinct when taking derivatives or refinements.}
{Volume I, Chapter 4, ``Complex-order Bernoulli interpolation.'' The complex-order family requires $a>0$; the plot uses 100-digit function evaluations at sampled real orders. Classical Hurwitz special values supply the integer interpolation.}

\lesson{PAPER I / 06}{A parent is exactly the sum of its children}{I06_parent_child_refinement}
{\mathfrak B_{M,r}(z;a)=\sum_{\ell=0}^{L-1}\mathfrak B_{LM,r+\ell M}(z;a).}
{The parent curve and the dashed sum coincide, although the two child curves differ substantially. Here the parent has address $(M,r)=(2,1)$ and is refined to addresses $(4,1)$ and $(4,3)$ at integer order four. Refinement does not approximate the parent or reduce its size; it exposes an exact decomposition that later supports detail channels.}
{At $a=1/3$, the parent is $8B_4(2/3)=52/405$. Its children are $64B_4(1/3)=416/405$ and $64B_4(5/6)=-364/405$. Their exact rational sum is $52/405$. The larger child magnitudes show why complete reconstruction matters.}
{Information in separate children can disappear when only their parent sum is retained.}
{Volume I, Chapter 5, Theorem 5.1, ``Parent-child refinement,'' with Theorem 5.2, ``Integer samples.'' The general complex-order identity uses $a>0$ and the classical Hurwitz multiplication formula.}

\lesson{PAPER I / 07}{Detail phases preserve complete information}{I07_unitary_details}
{W_j=L^{-1/2}\sum_{\ell=0}^{L-1}\omega_L^{-j\ell}b_\ell,\qquad \sum_j|W_j|^2=\sum_\ell|b_\ell|^2.}
{The left panel shows four real child values. The right panel places their transformed channels in the complex plane. A nonreal channel is a phase coordinate, not a lost child. The normalized Fourier transform has an inverse and preserves the total squared magnitude. The DC channel is the parent divided by $\sqrt L$; the other channels record differences between children.}
{Here $L=4$, $n=4$, and $a=1/3$. Applying the inverse formula $b_\ell=L^{-1/2}\sum_j\omega_L^{j\ell}W_j$ reconstructs each individual child. Real input also gives conjugate channels $W_3=\overline{W_1}$, while the DC and Nyquist channels are real.}
{A detail bank should come with inverse synthesis and a normalization certificate.}
{Volume I, Chapter 6, Definition 6.1, ``Cyclotomic detail,'' and the subsequent inverse-transform and Parseval identities. Finite integer-order polynomial data; the figure uses normalized channels, not unnormalized resultant evaluations.}

\lesson{PAPER I / 08}{Bernoulli and Euler share one channel bank}{I08_bernoulli_euler_bank}
{W_{2,0}(n;a)=\frac{B_n(a)}{\sqrt2},\qquad W_{2,1}(n;a)=-\frac{n}{2\sqrt2}E_{n-1}(a).}
{Compare the two panels at the same translation parameter. The left pair consists of Bernoulli children; the right pair consists of their sum and difference channels. The low-pass output is a Bernoulli polynomial and the high-pass output is an Euler polynomial with the displayed scale factor. Both pairs carry the same information. The plotted curves have order four.}
{At order three and $a=1/2$, the children are $3/16$ and $-3/16$. Their DC output vanishes and their detail output is $3/(8\sqrt2)$. This agrees with $E_2(1/2)=-1/4$ in the high-pass formula.}
{Classical special functions can occur as complementary outputs of one exact refinement.}
{Volume I, Chapter 8, Theorem 8.1, ``Two-channel identity.'' Positive integer orders; endpoint $a=0$ in this figure uses the Bernoulli-polynomial extension, not a zero-argument complex-order Hurwitz evaluation.}

\lesson{PAPER I / 09}{A Wronskian can reject a broad conjecture}{I09_wronskian_negative_control}
{\mathcal W_4(a)=\frac{a(a-1)}5\bigl(5a^4-10a^3+5a-2\bigr).}
{The real graph is positive between zero and one, but its complete zero set tells a more precise story. The right panel includes two nonreal roots. A favorable sign on one real interval therefore does not imply universal real-rootedness. The shaded interval highlights a valid local feature, while the root plot supplies the negative control against a stronger claim.}
{The nonreal pair is $\frac12\pm i\sqrt{(2\sqrt{65}-15)/20}$. The other four roots are $0$, $1$, and $\frac12\pm\sqrt{(15+2\sqrt{65})/20}$. At $a=1/2$, the value is $7/320>0$. These exact formulas explain both panels.}
{A zero diagnostic is useful when it records counterexamples as carefully as positive cases.}
{Volume I, Chapter 9, Worked example 9.2, ``Orders two through four,'' and Boundary 9.3. The plot uses the exact degree-six polynomial and algebraic roots; no general interlacing theorem is inferred.}

\lesson{PAPER I / 10}{Third-root norms remember parity}{I10_cyclotomic_norm_parity}
{N_n=\begin{cases}[3^{n-1}(B_n-B_n(1/3))]^2,&n\ge2\text{ even},\\3[3^{n-1}B_n(1/3)]^2,&n\ge3\text{ odd}.\end{cases}}
{The equal second and third children in the left panel collapse the even-order primitive detail to a rational number. The right panel divides each norm by its corresponding rational square. The ratio is one at even orders and three at odd orders. The rational square factor has different definitions in the two cases, supplied by Bernoulli reflection.}
{At $n=8$, the children are $-729/10$, $1093/30$, and $1093/30$. Because $\omega_3+\omega_3^2=-1$, the unnormalized primitive detail is $-328/3$. Its field norm is $107584/9$. At odd order three, the primitive detail is purely imaginary and its norm is $1/3$.}
{A cyclotomic norm combines exact arithmetic with the symmetry of its channel orbit.}
{Volume I, Chapter 10, Theorem 10.1, ``Parity law for the third-root detail norm,'' and Chapter 21, order-eight example. Integer Bernoulli-polynomial samples at $a=0$; unnormalized details are used.}

\lesson{PAPER I / 11}{Scale filters select an arithmetic window}{I11_kernel_localization}
{f_k(x)=x^{-2}(1-x^{-2})^k,\qquad x_*(k)=\sqrt{k+1}.}
{Each curve divides the kernel by its own maximum and the horizontal coordinate by its exact maximizing location. The maxima align at one. The curves become similar as the packet order rises, but their finite-order forms remain distinct. Small integers are suppressed by the filter and large integers by the power weight, producing an active window.}
{For $k=160$, the center is $\sqrt{161}\approx12.689$. More generally, the kernel $x^{-a}(1-x^{-\beta})^k$ has center $(1+\beta k/a)^{1/\beta}$. Differentiating its logarithm identifies the center directly. Here all plotted points satisfy $x\ge1$ and the normalized coordinate ranges from 0.18 to 4.}
{A scale center is a localization statement, not yet a cancellation estimate for signs.}
{Volume I, Chapter 14, ``Scale localization and the Mertens wavelet.'' Continuous kernels at $k=40,160,640$; this is an exact elementary kernel calculation. Zero-mean wavelet statements in that chapter require $k\ge1$.}

\lesson{PAPER I / 12}{Coherence has a proved scale and an open estimate}{I12_signed_coherence_boundary}
{\Gamma_k=\frac{|c_k|^2}{E_k},\qquad \mathrm{RH}\iff\Gamma_k=O_\varepsilon(k^\varepsilon)\ \text{for every }\varepsilon>0.}
{The left panel compares the signed amplitude, the square root of positive energy, and the amplitude with all arithmetic signs made positive. These are different quantities. The right panel plots finite principal coherence and an absolute truncation-tail band. Finite small values do not establish the required asymptotic bound. The denominator scale is proved; the remaining assertion concerns Möbius cancellation.}
{At $k=1200$ and cutoff $N=200000$, the finite calculation gives $|c_k|\approx1.1153\times10^{-5}$ and $E_k\approx2.3091\times10^{-6}$. Their ratio is about $5.3868\times10^{-5}$. The omitted amplitude is bounded absolutely by $1/N$, and the omitted energy by $1/(3N^3)$.}
{Positive energy controls normalization; the signed estimate remains RH-equivalent.}
{Volume I, Theorems 15.1 and 16.1; Chapter 22's finite example. The RH equivalence imports the classical Báez-Duarte theorem. Bands control truncation only, not interval-certified floating-point rounding.}

\startvolume{Paper II}{Local germs, positive moments, and character geometry}

\lesson{PAPER II / 01}{Stieltjes constants encode a detail's local slope}{II01_stieltjes_tangent}
{G_{1,0}(1)=\sqrt2\log2,\qquad W_1'(0;1)=-\frac{\log2}{\sqrt2}.}
{The first panel compares the two-channel detail near its removable zero with its exact tangent. The second removes the tangent and displays the remaining curvature. Generalized Stieltjes constants are Taylor data for the Hurwitz-zeta Laurent expansion. Their finite Fourier transform therefore supplies the local expansion of a Bernoulli detail, rather than just its integer values.}
{For $L=2,M=1,a=1$, the child Hurwitz arguments are $1/2$ and 1. Their zeroth Stieltjes difference is $\gamma_0(1/2)-\gamma_0(1)=2\log2$. Dividing by $\sqrt2$ gives $G_{1,0}$. The derivative formula adds the factor $-1/2$, giving the slope about $-0.490129$.}
{Local derivatives and integer samples probe different information in the same entire function.}
{Volume II, Chapter 2, Theorem 2.2, ``Complete local packet expansion,'' and Example 2.3. Fixed $a=1>0$; sampled real orders $-0.2\le z\le0.3$ evaluated at 100-digit precision.}

\lesson{PAPER II / 02}{A character packet returns an L-value}{II02_character_stieltjes}
{\sum_{r=1}^{q}\chi(r)\gamma_0(r/q)=qL(1,\chi)\qquad(\chi\text{ nonprincipal}).}
{The first two bars are signed contributions and the third is their sum. The character removes the pole term of the Hurwitz-zeta representation. What remains is a finite transform of local constants that equals the value of the nonprincipal L-function at one. Cancellation of the common pole is essential to this formula.}
{For the character modulo four, $\chi_4(1)=1$, $\chi_4(3)=-1$, and the even residues contribute zero. Since $\gamma_0(x)=-\psi(x)$, digamma reflection gives $\gamma_0(1/4)-\gamma_0(3/4)=\pi$. Thus $4L(1,\chi_4)=\pi$, agreeing with the classical Dirichlet-beta value.}
{The higher Stieltjes packets recover the whole derivative sequence $L^{(m)}(1,\chi)$.}
{Volume II, Chapter 3, Theorem 3.2, ``Character transform at $s=1$.'' This example fixes actual modulus $q=4$ and a nonprincipal character; the principal character needs separate pole treatment.}

\lesson{PAPER II / 03}{Detail variance measures child disagreement}{II03_detail_variance}
{\mathcal V_L=\sum_\ell(b_\ell-\bar b)^2=\frac1L\sum_{\ell<m}(b_\ell-b_m)^2,\quad \bar b=L^{-1}\sum_\ell b_\ell.}
{The first panel compares three real child values with their common mean. The second shows their squared deviations. Replacing the complete child vector by its mean removes exactly these deviations. For real children the variance also equals the summed squared magnitudes of non-DC unitary channels. For complex children, the holomorphic formula uses squares and paired channels, not absolute values.}
{At $L=3,n=4,a=0$, the integer-polynomial children are $-9/10,13/30,13/30$, with mean $-1/90$. The pairwise form has two nonzero differences of magnitude $4/3$. Hence $\mathcal V_3=(1/3)\cdot2(4/3)^2=32/27$.}
{A scalar invariant can quantify lost detail while complete synthesis retains every child.}
{Volume II, Chapter 4, Theorem 4.2, ``Three exact forms,'' and Chapter 5, Example 5.2. The displayed children are real integer-order polynomial samples; complex holomorphic variance is not a positive norm.}

\lesson{PAPER II / 04}{A resultant compresses primitive detail channels}{II04_primitive_resultants}
{\mathcal N_3=\operatorname{Res}(X^2+X+1,P(X))=P(\omega_3)P(\omega_3^2)=3W_1W_2.}
{Each curve evaluates the child polynomial on the entire unit circle. Highlighted points evaluate it at primitive third roots. Their product is the resultant. The order-three pair is imaginary and conjugate; the order-four pair coincides on the real axis. The construction compresses the primitive orbit into one scalar, while the plotted curve retains more detail.}
{At order three, $P(X)=(X-X^2)/3$ and the primitive product is $1/3$. At order four, $P(X)=-9/10+(13/30)(X+X^2)$; each primitive value is $-4/3$, so the resultant is $16/9$. The factor 3 appears because $W_j$ is normalized by $1/\sqrt3$.}
{Separate field-resultant normalization from unitary-channel normalization.}
{Volume II, Chapter 6, Definition 6.1 and Theorem 6.2, ``Primitive collapse criterion.'' Here $a=0$ is used only for positive-integer Bernoulli-polynomial samples. No complete global zero classification is claimed.}

\lesson{PAPER II / 05}{The energy is a positive moment measure}{II05_positive_measure}
{\nu=\sum_nw_n\delta_{x_n},\quad w_n=\mu(n)^2n^{-4},\quad x_n=(1-n^{-2})^2,\quad E_k=\int x^k\,d\nu.}
{The left panel retains the large atom at $x_1=0$. The right magnifies the remaining weights on a logarithmic scale. The weights are positive at squarefree integers and vanish elsewhere. Atom locations approach one while their masses decay. This produces a finite measure with infinite support, the structural source of the later positivity and spectral results.}
{At $n=1$, the atom has weight 1 and contributes to $E_0$ under the moment convention $0^0=1$. At $n=2$, $x_2=9/16$ and $w_2=1/16$. The total mass is exactly $E_0=\zeta(4)/\zeta(8)=105/\pi^4$.}
{Keep the zero atom when forming moments and the resolvent.}
{Volume II, Chapter 7, ``The squarefree packet energy.'' Principal parameters $a=\beta=2,q=1$. The picture shows positive atoms through $n=60$; the mathematical measure has infinitely many atoms.}

\lesson{PAPER II / 06}{Every alternating difference layer is positive}{II06_complete_monotonicity}
{(-1)^m\Delta^mE_k=\sum_nw_nx_n^k(1-x_n)^m\ge0.}
{The horizontal direction changes packet depth and the vertical direction changes difference order. Color records the logarithm of a direct positive sum. The figure does not form large differences of nearly equal numerical energies. The moment representation supplies the positive expression first, making the sign transparent at every displayed depth and difference order.}
{For $m=1$, $E_k-E_{k+1}=\sum_nw_nx_n^k(1-x_n)$. For $m=2$, $E_k-2E_{k+1}+E_{k+2}=\sum_nw_nx_n^k(1-x_n)^2$. Both are positive because the measure has positive mass away from the endpoints. The same mechanism continues to every integer $m\ge0$.}
{Complete monotonicity is an exact representation, not a pattern inferred from a heatmap.}
{Volume II, Chapter 7, Theorem 7.1, ``Complete monotonicity.'' Plot: $1\le k\le30$, $0\le m\le6$, $N=200000$. Omitted layer $m$ is at most $2^m/[(3+2m)N^{3+2m}]$; rounding is separate.}

\lesson{PAPER II / 07}{Hankel positivity is a Gram-space statement}{II07_hankel_gram}
{c^*H_dc=\int\left|\sum_{j=0}^{d}c_jx^j\right|^2d\nu(x)>0\qquad(c\ne0).}
{The first panel normalizes the entries of a three-by-three Hankel matrix to show correlations among monomials. The second plots its positive eigenvalues using high-precision low-order moments. Very unequal eigenvalues do not contradict positivity. The proof uses infinitely many distinct atoms: a nonzero polynomial cannot vanish at all of them.}
{For the two-by-two matrix, strict positivity reads $E_0E_2-E_1^2>0$. With principal parameters it is about $0.03026887920$. This follows from the nonconstant function $x$ having positive variance under the normalized measure, and independently agrees with the exact zeta-ratio evaluations.}
{Finite numerical eigenvalues illustrate an exact Gram argument and must respect conditioning.}
{Volume II, Chapter 8, Theorem 8.1, ``Strict finite Hankel positivity.'' Displayed moments and eigenvalues use 100-digit arithmetic. The claim is strict positivity for every finite degree, not just the displayed matrix.}

\lesson{PAPER II / 08}{Positive mass gives a zero-free slit resolvent}{II08_zero_free_resolvent}
{\operatorname{Im}\mathcal E(z)=\operatorname{Im}z\int\frac{x\,d\nu(x)}{|1-zx|^2},\qquad z\notin[1,\infty).}
{The upper half-plane has positive imaginary part and the lower half-plane negative imaginary part. The dark segment marks the excluded physical cut, not a line of proved zeros. On the real axis below one, the resolvent is positive. These three regions exclude zero everywhere in the slit domain. The proof uses positive mass on $x>0$, not signed Möbius cancellation.}
{At $z=0$, $\mathcal E(0)=E_0>0$. At $z=i$, every positive atom contributes a positive imaginary part $w_nx_n/(1+x_n^2)$. The zero atom adds a real constant and therefore cannot cancel that imaginary part.}
{Zero-freeness here belongs to the positive resolvent family.}
{Volume II, Chapter 8, Theorem 8.3, ``Zero-free slit plane.'' Grid: $-2\le\Re z\le2$, $|\Im z|\le1.5$, $N=2000$. The audit bounds omitted terms using the minimum denominator distance on each retained sampled point; the cut is masked.}

\lesson{PAPER II / 09}{A zeta-ratio stencil exactly recovers the energy}{II09_zeta_ratio_stencil}
{E_k=\sum_{j=0}^{2k}(-1)^j\binom{2k}{j}\frac{\zeta(4+2j)}{\zeta(8+4j)}.}
{The first panel overlays the exact finite stencil with a direct truncated positive sum. The second resolves the small difference between them. This difference is the omitted positive tail, not disagreement between two mathematical definitions. The stencil comes from expanding the packet filter and interchanging a finite sum with an absolutely convergent Dirichlet series.}
{At $k=1$, the stencil has three terms: $E_1=D(4)-2D(6)+D(8)$, where $D(s)=\zeta(s)/\zeta(2s)$. It gives approximately $0.04780461034$. Direct summation through $N=2000$ undershoots the infinite value by at most $1/(3N^3)$.}
{Exact coordinate changes can require extra numerical precision when their coefficients alternate.}
{Volume II, Chapter 9, Theorem 9.1, ``Exact zeta-ratio finite difference.'' Principal parameters; $0\le k\le6$. Stencil values use 100-digit arithmetic, and the finite sums use the explicit absolute tail bound.}

\lesson{PAPER II / 10}{Squarefree counting fixes the energy scale}{II10_energy_asymptotic}
{E_k\sim A k^{-3/2},\qquad A=\frac{3}{2^{5/2}\pi^{3/2}}.}
{Multiplying by $k^{3/2}$ removes the leading power decay. The finite arithmetic curve approaches the horizontal constant. The dashed curve is the density contribution expressed by a beta integral. The remaining difference includes arithmetic counting error. The asymptotic theorem uses the classical squarefree counting estimate; the graph displays its finite behavior.}
{For the principal parameters, the density term is $(3/\pi^2)\mathrm B(3/2,2k+1)$. The beta-function asymptotic yields $A\approx0.09524045390$. At general $a>1/2,\beta>0$, the exponent is $\alpha=(2a-1)/\beta$ and the leading constant also contains the coprime squarefree density for the actual modulus.}
{Record which counting theorem supplies an asymptotic result and which data merely illustrate it.}
{Volume II, Chapter 10, Theorem 10.1, ``Packet-energy asymptotic.'' Plot: $10\le k\le100000$, $N=200000$; absolute energy tail $\le1/(3N^3)$. Classical counting and beta asymptotics are source inputs.}

\lesson{PAPER II / 11}{The active coordinate approaches a Gamma law}{II11_gamma_packet_law}
{U_k=\frac{2k}{n^2}\Longrightarrow\operatorname{Gamma}(3/2,1),\qquad \mathbb E[U^m]=\frac{\Gamma(3/2+m)}{\Gamma(3/2)}.}
{The distribution curves describe where the normalized positive energy lives in inverse-scale coordinates. Different packet depths approach the same Gamma cumulative distribution. The second panel checks two moments at a large finite depth. The law is weighted by packet energy; it is not the unweighted histogram of integers or squarefree labels.}
{The limiting mean is $3/2$ and the limiting second moment is $15/4$. For each integer $m\ge0$, the finite-depth identity shifts the energy parameter: $\mathbb E_k[U_k^m]=(2k)^mE_{2+m,2,1}(k)/E_{2,2,1}(k)$. Applying the energy law gives those Gamma moments.}
{The moment identity supplies more information than weak convergence alone.}
{Volume II, Chapter 11, Theorem 11.1, ``Gamma scaling limit.'' $k=100,1000,10000$, $N=200000$; CDFs sampled on $0\le U\le9$. Omitted normalized probability is bounded by $1/(3N^3E_k^{(N)})$.}

\lesson{PAPER II / 12}{Distinct characters decohere at a fixed modulus}{II12_character_decoherence}
{\frac{G_{\chi,\psi}(k)}{E_{a,\beta,q}(k)}=O_{a,\beta,q}(k^{-1/(2\beta)})\qquad(\chi\ne\psi\bmod q).}
{The diagonal of the normalized character Gram matrix stays one. The off-diagonal entries measure overlap between distinct multiplicative phases. The second panel follows the largest such overlap as depth increases. Its reference line shows the theorem's rate shape, anchored at the first point; it is not a proved numerical upper envelope. All four characters are represented with the same actual modulus five.}
{Here $a=\beta=2$, so the normalized source rate is $k^{-1/4}$. The product $\chi\overline\psi$ is nonprincipal when the two characters differ. Its squarefree weighted partial sums are $O_q(\sqrt X)$, which supplies the off-diagonal estimate after partial summation.}
{Cross-character cancellation differs from cancellation inside the signed principal Möbius amplitude.}
{Volume II, Chapter 12, Lemma 12.1 and Theorem 12.2, ``Character decoherence.'' $N=200000$, $20\le k\le100000$, fixed modulus $q=5$. The audit records a normalized-entry truncation bound. No character-uniform GRH estimate is inferred.}

\startvolume{Paper III}{Jacobi models, endpoint limits, and the spectral edge}

\lesson{PAPER III / 01}{The zero packet belongs to the positive contraction}{III01_zero_atom_contraction}
{T|n\rangle=x_n|n\rangle,\qquad E_k=\langle v,T^kv\rangle,\qquad \|T\|=1.}
{The left panel distinguishes zero depth from positive depth. The atom at $x_1=0$ contributes its unit mass to $E_0$ and vanishes from $E_k$ for $k\ge1$. The right panel shows the spectral gap of each integer node approaching zero. Individual nodes lie below one, but their supremum is one; there is no uniform gap for the infinite contraction.}
{At $n=2$, $x_2=9/16$ and $1-x_2=7/16$. Along squarefree primes tending to infinity, $1-x_n=2/n^2-1/n^4\to0$. The logarithmic scale operator $-\log(1-n^{-2})$ is defined on $n\ge2$; the zero packet is retained separately.}
{A zero-depth identity must preserve the endpoint atom and state its convention.}
{Volume III, Chapters 2 and 13, ``Positive contraction realization'' and ``Number-basis realization.'' Figure uses principal parameters. The gap panel shows all integer nodes; squarefree nodes select an infinite subsequence with the same limiting norm.}

\lesson{PAPER III / 02}{Generalized Hankel minors are strictly positive}{III02_generalized_minors}
{\det[E_{r_i+s_j}]_{i,j=0}^{d}>0\qquad(r_0<\cdots<r_d,\ s_0<\cdots<s_d).}
{Choose a row label and a column label to specify two strictly increasing exponent sets. Each colored cell is a two-by-two generalized minor, which need not be a principal minor of a symmetric matrix. High-precision values show several such witnesses. The proof expands the determinant into products of generalized Vandermonde determinants at ordered positive support nodes.}
{For row exponents $(0,1)$ and column exponents $(0,2)$, the minor is $E_0E_3-E_1E_2$. For rows and columns both $(1,3)$, it is $E_2E_6-E_4^2$. The infinite support contains enough positive nodes to make the relevant Vandermonde contribution strictly positive.}
{Total positivity is stronger than positivity of principal Hankel quadratic forms.}
{Volume III, Chapter 3, Theorem 3.1, ``Strict total positivity.'' Nonnegative, strictly increasing integer exponent sets; low-order values here use 100-digit exact zeta-ratio moments. The theorem covers all finite sizes, not only this grid.}

\lesson{PAPER III / 03}{One measure determines a Jacobi recurrence}{III03_global_jacobi}
{\begin{gathered}xp_m(x)=b_{m+1}p_{m+1}(x)+a_mp_m(x)+b_mp_{m-1}(x),\\ b_m>0\ (m\ge1),\qquad b_0=0,\qquad p_{-1}=0.\end{gathered}}
{The matrix panel shows only nearest-neighbor couplings in the orthogonal-polynomial basis. The coefficient panel separates diagonal positions from off-diagonal couplings. This Jacobi system represents the same positive measure that was diagonal in the integer basis. Its coefficients are determined by moments rather than selected as free fitting parameters. The displayed global system is distinct from the later tilted, scaled system.}
{The first diagonal coefficient is $a_0=\int x\,d\nu/\int d\nu=E_1/E_0$. The first off-diagonal coefficient is the square root of the normalized variance of $x$. Both therefore follow directly from the first three moments. Higher coefficients continue the orthogonalization.}
{Changing basis makes a diagonal arithmetic operator into a canonical tridiagonal transfer.}
{Volume III, Chapter 4, ``The zero-free Stieltjes-Jacobi function,'' recurrence and Jacobi matrix after Theorem 4.1. The leading six-by-six block is computed from 100-digit moment data.}

\lesson{PAPER III / 04}{A Gaussian node can lie between arithmetic atoms}{III04_gaussian_support_gap}
{\xi_1=\frac{E_1}{E_0},\qquad 0<\xi_1<\frac5{48}<\frac9{16}.}
{The colored atoms lie on the horizontal support line. The diamond is the one-node Gaussian location. It is a weighted mean inside the interval spanned by the measure and lies in a gap of the discrete arithmetic spectrum. Gaussian compression therefore does not preserve the individual atom addresses. The shaded region isolates the simplest exact counterexample to that stronger claim.}
{For principal parameters, $\xi_1\approx0.04434860610$, while the first positive arithmetic atom is $9/16=0.5625$. Since $E_0\ge1$ and $E_1<\sum_{n\ge2}n^{-4}<5/48$, the gap statement follows without numerical computation. The one-node rule matches moments zero and one.}
{Use the support interval for Gaussian nodes, rather than claiming membership in the discrete spectrum.}
{Volume III, Chapter 5, Theorem 5.1, ``Moment-exact packet reduction.'' This is the consolidation correction to the original abstract's spectrum wording. The displayed mean uses 100-digit exact moments.}

\lesson{PAPER III / 05}{Four nodes preserve exactly eight moments}{III05_exact_moment_contract}
{E_r=E_0\sum_{j=1}^{m}\lambda_j\xi_j^r=E_0e_1^TJ_m^re_1,\qquad 0\le r\le2m-1.}
{Bars before the vertical divider lie below the display floor; they witness the exact moment contract numerically at high precision. Bars after the divider show the first moments beyond that contract. The divider marks degree seven for a four-node rule. Matching early moments is a precise finite statement, not equality of the complete infinite measure and reduced model.}
{With $m=4$, degrees $0,1,\ldots,7$ are exact by Gaussian quadrature. Degree eight is outside that guarantee and displays a nonzero residual. The degree-zero equality states that positive normalized weights sum to one, while the degree-one equality preserves the measure's mean.}
{State which moments a reduction preserves and where its guarantee ends.}
{Volume III, Chapter 5, Theorem 5.1, ``Moment-exact packet reduction.'' Gaussian nodes and weights use 100-digit arithmetic; residuals are plotted with floor $10^{-60}$. This is a numerical witness of the source theorem, not a new universal proof.}

\lesson{PAPER III / 06}{Positive spectral weights have negative z-residues}{III06_poles_and_residue_signs}
{\mathcal E_m(z)=E_0\sum_j\frac{\lambda_j}{1-z\xi_j},\qquad \operatorname{Res}_{z=1/\xi_j}\mathcal E_m=-\frac{E_0\lambda_j}{\xi_j}<0.}
{The first panel places positive weights at their physical-cut poles. The second records the residues in the variable $z$. Their signs are negative because differentiating the denominator contributes $-\xi_j$. The weights remain positive. This distinction depends on the coordinate and should not be hidden under the word positivity. The logarithmic pole axis keeps distant poles visible.}
{For a one-node model, the pole is about $22.54862301$ and the residue about $-24.30579519$. For the plotted four-node model, each node lies strictly in $(0,1)$, so each pole lies strictly above one. No exact Gaussian node equals the zero atom.}
{Identify whether positivity refers to weights, moments, or coefficients in a chosen coordinate.}
{Volume III, Chapter 5, Corollary 5.2, corrected in consolidation. Four-node values use 100-digit moments. The formula repairs the original ``positive residues'' wording while preserving its positive-measure content.}

\lesson{PAPER III / 07}{The two Gamma coordinates share all moment limits}{III07_all_moment_bridge}
{u_k=U_k-\frac{U_k^2}{4k},\qquad 0\le U_k^m-u_k^m\le\frac{m}{4k}U_k^{m+1}\quad(m\ge1).}
{Solid and dashed curves use $U_k=2k/n^2$ and $u_k=k(1-x_n)$, with $x_n=(1-n^{-2})^2$, and compare their first three moments with Gamma targets. They approach the same targets without being equal at finite depth. The displayed inequality supplies the bridge: convergence of the next $U$ moment bounds the discrepancy in the current moment. Weak convergence against bounded test functions alone would not provide this conclusion.}
{For $m=1$, the exact difference is $\mathbb E[U_k]-\mathbb E[u_k]=\mathbb E[U_k^2]/(4k)$. Its limiting numerator is $15/4$, so the discrepancy is asymptotic to $15/(16k)$. Both means still tend to $3/2$.}
{All-moment control is the needed input for polynomial and determinant convergence.}
{Volume II, Theorem 11.1's shifted-energy identity; Volume III, Theorems 6.1 and 7.1; explicit consolidation bridge. $N=200000$, $20\le k\le10000$; the audit records absolute normalized-moment truncation bounds.}

\lesson{PAPER III / 08}{Fixed-degree packet polynomials become Laguerre}{III08_laguerre_polynomials}
{P_m^{(k)}(u)\longrightarrow\Lambda_m(u)=(-1)^m m!L_m^{(1/2)}(u)\quad(m\text{ fixed}).}
{Each panel fixes polynomial degree and changes packet depth. The dashed curve is the classical monic Laguerre limit. The finite packet polynomials are monic and come from the tilted arithmetic measure. They are generally different at finite depth. All-moment convergence gives convergence of their coefficient systems; fixing degree is essential to that implication. The plot does not claim convergence uniform in degree.}
{At degree two the limit is $u^2-5u+15/4$, with roots approximately $0.91886117$ and $4.08113883$. At degree three it is $u^3-(21/2)u^2+(105/4)u-105/8$. These coefficients follow from orthogonality under the normalized Gamma-$3/2$ law.}
{Polynomial universality follows from a moment bridge and a nonsingular limiting Gram system.}
{Volume III, Chapter 7, Theorem 7.1, ``Fixed-degree Laguerre universality,'' with the consolidated all-moment justification. Degrees two and three, $k=100,10000$, $N=200000$, displayed on $0\le u\le8$.}

\lesson{PAPER III / 09}{The scaled Jacobi coefficients have a fixed point}{III09_jacobi_fixed_point}
{a_m(k)\to2m+\alpha,\qquad b_{m+1}(k)\to\sqrt{(m+1)(m+\alpha)},\qquad \alpha=3/2.}
{The left panel follows diagonal recurrence coefficients and the right follows nearest-neighbor couplings. At larger depth the arithmetic Jacobi block approaches the Gamma--Laguerre block. These coefficients belong to the endpoint coordinate $u$, so they should not be compared directly with the unscaled global Jacobi coefficients on $x\in[0,1]$. Only each fixed leading block is asserted to converge.}
{For $m=0,1,2$, the limiting diagonal is $3/2,7/2,11/2$. The first three off-diagonal entries are $\sqrt{3/2},\sqrt5,\sqrt{21/2}$. They determine the first leading block and its Gaussian rule.}
{Specify both the coordinate and the fixed-degree range of a recurrence limit.}
{Volume III, Chapter 7, Corollary 7.2, ``Jacobi fixed point.'' Principal shape $\alpha=3/2$; first five coefficients at $k=100,10000$, $N=200000$. Finite Lanczos computations use reorthogonalization.}

\lesson{PAPER III / 10}{A shifted determinant has an explicit leading constant}{III10_shifted_hankel_asymptotic}
{\Delta_d(k)\sim A^{d+1}\!\prod_{j=0}^d j!(3/2)_j\;k^{-3(d+1)/2-d(d+1)}\quad(d\text{ fixed}).}
{Each curve divides a shifted Hankel determinant by its predicted leading term. The ratios approach one. The computation changes basis to the endpoint coordinate before evaluating a determinant, avoiding subtraction of nearly equal shifted energies. This numerical method follows the exact affine basis-change argument used in the proof. The plot illustrates fixed determinant orders one, two, and three.}
{For $d=1$, the determinant is $E_kE_{k+2}-E_{k+1}^2$. Its leading term is $(3/2)A^2k^{-5}$. The factor $k^{-2}$ beyond the two energy factors comes from the affine transformation of the two monomial bases.}
{Small determinants need a stable coordinate representation as well as an asymptotic formula.}
{Volume III, Chapter 8, Theorem 8.1, ``Complete leading Hankel asymptotic.'' $30\le k\le10000$, $N=200000$. Underlying moment tails are recorded; this plot is not a certified determinant-error interval.}

\lesson{PAPER III / 11}{The kernel and Gaussian rule share the same limit}{III11_kernel_quadrature_limit}
{K_d^{(k)}\to K_d^{(3/2)},\qquad (\xi_j^{(k)},\lambda_j^{(k)})\to(\xi_j^{\mathrm{GL}},\lambda_j^{\mathrm{GL}})\quad(d,m\text{ fixed}).}
{The left panel sums three orthonormal polynomial squares on the diagonal of a fixed-degree kernel. The right compares the associated three-node Gaussian rule with its normalized Gauss--Laguerre limit. Their common input is the converging finite Jacobi block. Kernel convergence and quadrature convergence are separate outputs of the same moment and recurrence structure.}
{The limiting three-node locations are approximately $0.66632591$, $2.80077505$, and $7.03289904$. Their positive weights integrate the normalized Gamma density, so they sum to one. This differs from unnormalized generalized-Laguerre weights, whose sum is $\Gamma(3/2)$.}
{State the measure normalization when comparing kernels or quadrature weights.}
{Volume III, Chapters 9--10, Theorems 9.1 and 10.1, ``Fixed-degree kernel convergence'' and ``Quadrature fixed point.'' $d=2,m=3,k=10000,N=200000$. Sampled diagonal kernels on $0\le u\le8$ illustrate the stronger local-uniform source result.}

\lesson{PAPER III / 12}{The edge cusp follows conditionally from energy decay}{III12_principal_edge_cusp}
{\mathcal E(1)-\mathcal E(z)\sim2\sqrt\pi A\sqrt{1-z}=\frac3{2^{3/2}\pi}\sqrt{1-z}\qquad(z\uparrow1).}
{Read the horizontal axis toward the right: $\varepsilon=1-z$ decreases. Dividing the positive resolvent difference by $\sqrt\varepsilon$ exposes the cusp constant. The shaded band adds a positive omitted-tail bound. The asymptotic implication follows from the source energy law through the consolidation's Abelian argument; it does not require an informal continuous density for a discrete measure.}
{Write the difference as $\sum_{k\ge0}E_k(1-z^k)$. Conditional on $E_k\sim Ak^{-3/2}$, set $t=-\log z$ and use the integrable profile $u^{-3/2}(1-e^{-u})$. Its integral is $2\sqrt\pi$, giving the displayed constant. Terms away from the endpoint contribute lower-order corrections.}
{A repaired analytic passage should name its source input and preserve its exact conclusion.}
{Volume III, Chapter 11, Theorem 11.1, ``Principal edge cusp,'' with the consolidation Abelian repair. $N=200000$, $10^{-5}\le\varepsilon\le0.1$; omitted difference $\le1/N$. The band controls truncation, not interval-certified rounding; no RH bound follows.}

% Consolidated entry point; each component owns exactly twelve lesson pages.
\startvolume{Paper IV}{Arithmetic depth, time, and signed coordinates}
\lesson{PAPER IV / 01}{Keep the signs while changing depth}{IV01_signed_seed}
{C_{a,\beta}(k,t)=\sum_{n\ge1}\frac{\mu(n)}{n^{a+it}}y_n^k,\qquad y_n=1-n^{-\beta},\quad a>1,\ \beta>0}
{The seed panel shows signed amplitudes. Squareful integers vanish, while squarefree integers can contribute either sign. The depth panel displays a separate positive filter. It favors larger integers as depth grows. Filtering does not remove the M\"obius signs or replace the signed amplitude with its squared magnitude.}
{Set $a=\beta=2$. The $n=1$ amplitude is one at depth zero. Its filter is zero at every positive depth, so this atom disappears immediately. At $n=2$, the amplitude is $-1/4$ and the depth-five filter is $(3/4)^5$. The complete coefficients begin $C(0,0)\approx0.607927$ and $C(1,0)\approx-0.316011$.}
{Arithmetic decay $a$, scale exponent $\beta$, discrete depth $k$, and logarithmic time $t$ have different jobs.}
{IV Chapter 1, ``One signed seed and two exact evolutions.'' Seed integers through 30; filters through 80. Subsequent principal calculations use $a=\beta=2$ and cutoff $N=200000$.}

\lesson{PAPER IV / 02}{Compute one coefficient in two coordinates}{IV02_finite_difference}
{C_{2,2}(k,0)=\sum_{j=0}^k(-1)^j\binom kj\frac1{\zeta(2+2j)},\qquad |C(k,0)-C_N(k,0)|\le\frac1N}
{The coefficient panel overlays a direct number sum and a finite difference of reciprocal zeta values. Both retain the signed cancellation. The error panel compares their difference with the number-cutoff bound. This bound concerns the omitted infinite tail. It does not describe rounding in a large alternating finite difference, which needs its own precision choice.}
{At depth one, the filter expansion gives $C(1,0)=\zeta(2)^{-1}-\zeta(4)^{-1}$, about $-0.316011$. At depth two it gives $\zeta(2)^{-1}-2\zeta(4)^{-1}+\zeta(6)^{-1}$, about $-0.256997$. Truncating the number sum at $N=200000$ introduces at most $5\times10^{-6}$ absolute error, uniformly in every displayed nonnegative integer depth.}
{Integer depth gives a finite algebraic identity. Infinite number support introduces a separately controlled approximation.}
{IV Chapter 1 finite-difference identity and Worked example 1.1, ``Principal Bernoulli--zeta orbit.'' Depths $0$--$10$; zeta coordinates use 90-digit arithmetic.}

\lesson{PAPER IV / 03}{Watch finite time extract the diagonal}{IV03_finite_time_average}
{\frac1{2T}\int_{-T}^{T}C_N(k,t+\tau)\overline{C_N(\ell,t)}\,dt=\sum_{n,m\le N}a_n\overline{b_m}\operatorname{sinc}\!\bigl(T\log(n/m)\bigr)}
{Each panel follows one component of an exactly integrated finite correlation. Off-diagonal pairs oscillate through their sinc factors. Diagonal pairs retain the lag phase and form the dashed target. Increasing the observation window need not make the error decrease monotonically. Different logarithmic frequency gaps have different cancellation scales.}
{For the pair $n=2$, $m=3$, the off-diagonal factor is $\sin(T\log(2/3))/(T\log(2/3))$. Its first positive zero occurs at $T=\pi/\log(3/2)$, about $7.748$. Other pairs generally remain at that time. The plots keep all $40^2$ pairs, use $k=\ell=3$, and preserve the fixed lag $\tau=0.5$.}
{Finite sinc extraction makes the diagonal mechanism explicit. The infinite time-average theorem has an additional absolute-convergence argument.}
{IV Theorem 2.1, ``Exact time-average identity.'' Here $a_n=\mu(n)n^{-2}y_n^3n^{-i/2}$ and $b_m=\mu(m)m^{-2}y_m^3$; sinc means $\sin x/x$. No infinite-model finite-time bound is inferred.}

\lesson{PAPER IV / 04}{The positive kernel is a complex Gram matrix}{IV04_positive_gram}
{v_{k,t}(n)=\mu(n)n^{-2}y_n^kn^{-it},\qquad K_N(k,\ell;t-s)=\sum_{n\le N}v_{k,t}(n)\overline{v_{\ell,s}(n)}}
{The heatmap shows magnitudes of normalized complex overlaps. The underlying matrix still retains complex phases and is Hermitian. The eigenvalue panel tests the displayed finite packet set. Positivity follows from the feature representation: every quadratic Gram sum is a squared norm. The numerical eigenvalues illustrate that identity rather than proving positivity for every possible packet set.}
{Take the eight packets with depths $0,10,100,1000$ and times zero or one. Divide each finite feature vector by its norm. Every diagonal overlap becomes one. For any coefficients $c_j$, the quadratic form is the norm squared of $\sum_jc_jv_j$, so it cannot be negative, even when individual off-diagonal entries are complex.}
{Time averaging removes arithmetic interference on the diagonal while leaving a positive depth-time geometry.}
{IV Theorem 2.2, ``Positive-definite packet kernel.'' Cutoff $N=200000$; every unnormalized omitted kernel entry is bounded by $1/(3N^3)$. Normalized finite overlaps are displayed.}

\lesson{PAPER IV / 05}{Use Beta to connect depth with Gamma}{IV05_beta_gamma_depth}
{K(k,\ell;\tau)=\frac{\delta}{\beta}B(\alpha+i\tau/\beta,k+\ell+1)+O\!\left((k+\ell)^{-(2a-1/2)/\beta}\right),\quad\alpha=\frac{2a-1}{\beta}}
{At zero lag, multiplying out the leading depth power makes the Gamma limit visible. The Beta term retains a finite-depth correction. The second panel restores a compact interval of nonzero lags and their complex exponent. A single real asymptotic at zero lag would not determine this complete complex law.}
{For $a=\beta=2$, shape is $\alpha=3/2$ and principal squarefree density is $\delta=6/\pi^2$. Writing $m=k+\ell$, the leading zero-lag kernel is $(3/\pi^2)\Gamma(3/2)m^{-3/2}$. At nonzero lag, the power becomes $m^{-3/2-i\tau/2}$. Removing both parts leaves the Gamma factor, whose real part is shown.}
{The continuous shape arises from arithmetic decay and scale. Compact lag intervals are part of the uniform theorem.}
{IV Theorem 4.1, ``Uniform complex packet asymptotic.'' Principal modulus one; $|\tau|\le2$. Finite-cut data illustrate convergence; the theorem's implicit error constant is not fitted from this plot.}

\lesson{PAPER IV / 06}{Remove the carrier before comparing limits}{IV06_carrier_dephasing}
{A_k(\tau)=\frac{K(k,k;\tau)}{K(k,k;0)},\qquad (2k)^{i\tau/\beta}A_k(\tau)\longrightarrow\frac{\Gamma(\alpha+i\tau/\beta)}{\Gamma(\alpha)}}
{The two panels compare real and imaginary components after carrier removal. Without that factor, the logarithmic frequency scale keeps drifting with depth. Normalization removes the changing positive energy; the carrier removes the changing phase reference. Both operations are needed to see the same shape-dependent limit across increasing depths.}
{Equal depth means the combined depth is $2k$. With $\beta=2$, the carrier is $\exp(i\tau\log(2k)/2)$, not $\exp(i\tau\log k/2)$. At zero lag the normalized finite kernel and its limit both equal one. For $k=100,1000,10000$, the figures use this same convention and compare against the shape-$3/2$ Gamma ratio.}
{Changing the phase reference does not alter modulus. It makes complex correlations at different arithmetic scales comparable.}
{IV Corollary 5.1, ``Equal-depth dephasing.'' Compact interval $|\tau|\le4$; finite energies normalize the displayed kernels. Number-cut errors and normalized bounds are recorded separately.}

\lesson{PAPER IV / 07}{An exact envelope can have non-Gaussian tails}{IV07_principal_envelope}
{\left|\frac{\Gamma(3/2+i\tau/2)}{\Gamma(3/2)}\right|^2=\frac{1+\tau^2}{\cosh(\pi\tau/2)}}
{The first panel overlays the Gamma expression with its elementary closed envelope. The second uses a logarithmic vertical scale to compare its tails with a Gaussian having the same limiting variance. Agreement at the initial curvature does not force agreement farther away. The Gamma recurrence and reflection identity determine the full envelope.}
{At lag one the squared modulus is $2/\cosh(\pi/2)$. At zero it is one. The limiting frequency variance is $\psi_1(3/2)/4\approx0.233701$, so the variance-matched Gaussian has squared characteristic modulus $e^{-0.233701\tau^2}$. The exact principal expression instead has an exponential tail multiplied by a quadratic factor.}
{A few moments determine local behavior. A closed characteristic-function formula determines every displayed lag.}
{IV Worked example 5.2, ``Closed principal envelope.'' Exact formulas on $0\le\tau\le8$. The Gaussian is an explicitly variance-matched teaching comparison, not an asserted approximation theorem.}

\lesson{PAPER IV / 08}{Center the logarithmic frequency law}{IV08_centered_log_frequency}
{Z_m=\log n-\beta^{-1}\log m\Longrightarrow-\beta^{-1}\log U,\qquad U\sim\operatorname{Gamma}(\alpha,1)}
{The histograms use positive squarefree packet weights after subtracting the drifting log scale. The limiting density remains continuous even though each finite spectrum consists of integer frequencies. The moment panel compares finite means and variances with their polygamma targets. Histogram convergence alone does not justify convergence of every unbounded moment.}
{For $\alpha=3/2$ and $\beta=2$, the limiting density is $2e^{-3z-e^{-2z}}/\Gamma(3/2)$. Its mean is $-\psi(3/2)/2\approx-0.018244987$ and its variance is $\psi_1(3/2)/4\approx0.23370055$. The finite weights are proportional to $\mu(n)^2n^{-4}(1-n^{-2})^m$, with their full finite mass used for normalization.}
{Log-frequency centering connects the arithmetic active scale with a fixed log-Gamma probability law.}
{IV Corollary 6.1, ``Polygamma cumulants.'' Histograms at $m=1000,10000$; finite moment records through $m=100000$. Preserve the source's moment-convergence proviso when claiming cumulant limits.}

\lesson{PAPER IV / 09}{Compare depth ratios before absolute depths}{IV09_cross_depth_overlap}
{\mathcal O(k,\ell)=\frac{K(k,\ell;0)}{\sqrt{K(k,k;0)K(\ell,\ell;0)}}\longrightarrow\left(\frac{2\sqrt r}{1+r}\right)^\alpha,\qquad r=k/\ell}
{The first panel compares several absolute depth scales against one ratio profile. Equal depths give overlap one. Exchanging the depths leaves the overlap unchanged, so reciprocal ratios share a value. The second follows a single four-to-one ratio while both depths increase. This separates ratio dependence from finite-depth arithmetic fluctuations.}
{Take $\alpha=3/2$ and $k=4\ell$. The limiting overlap is $(4/5)^{3/2}\approx0.715541753$. The same value occurs at ratio $1/4$. At ratio one the limit is exactly one. The displayed comparisons keep ratios between $1/16$ and $16$, and round each sampled larger depth to an integer before evaluating the finite kernel.}
{Comparable depths retain overlap even at large scale. Multiplicative depth separation is measured by a log-depth gap.}
{IV Theorem 7.1, ``Universal scale overlap,'' and Worked example 7.2, ``A four-to-one depth ratio.'' Compact positive ratio domain; finite number-cut sums are not asymptotic proofs.}

\lesson{PAPER IV / 10}{Use the actual character modulus}{IV10_character_sectors}
{\chi_j(2^r)=i^{jr}\quad(\bmod\ 5),\qquad\chi_j(n)=0\ \text{if }5\mid n,\qquad K_{\chi,\psi}=\sum_n\mu(n)^2\chi(n)\overline{\psi(n)}n^{-4-i\tau}y_n^{k+\ell}}
{Each matrix compares all four characters modulo five at one equal depth. Diagonal sectors have unit normalized magnitude. Off-diagonal sectors use nonprincipal character products and weaken with depth. The coprimality restriction is determined by the actual modulus. Replacing it with conductor can change even the principal sector's support and normalization.}
{The units are the cycle $1,2,4,3$ generated by two. For $j=1$, their values are $1,i,-1,-i$, while multiples of five have value zero. The principal character modulo five has conductor one but still vanishes at five. Its squarefree density is $(5/6)/\zeta(2)$, not $1/\zeta(2)$.}
{Distinct characters must share a fixed actual modulus; their product is then nonprincipal and the stated decoherence bound applies.}
{IV Theorem 8.1, ``Character depth-time decoherence,'' with the consolidation's actual-modulus correction. Depths $10,10000$, zero lag; comparable-depth and compact-lag hypotheses govern the asymptotic theorem.}

\lesson{PAPER IV / 11}{Change generating coordinates without losing signs}{IV11_resolvent_borel}
{\begin{gathered}(1-z)\mathcal C_a(z,t)+z\mathcal C_{a+\beta}(z,t)=L(a+it,\chi)^{-1},\\ e^{-x}\sum_{k\ge0}C_a(k,t)\frac{x^k}{k!}=\mathcal R_a(x,t)/x.\end{gathered}}
{The first panel combines two signed resolvents using the parameter-shift identity. The result returns the initial signed coefficient. The second changes ordinary discrete-depth coordinates to normalized exponential coordinates. Removing $e^x$ reveals the Riesz companion. These identities retain the signs throughout; positivity or zero-freeness does not follow from a coordinate change.}
{For the principal choice $a=\beta=2$, $t=0$, the shifted combination returns $1/\zeta(2)$. The normalized Borel value is $\sum_n\mu(n)n^{-2}e^{-x/n^2}$. At $x=0$, define $\mathcal R(x)/x$ by this continuous expression, which returns the same initial coefficient. The finite number cutoff and depth-40 EGF truncation introduce separate errors.}
{Resolvent shifts and Borel--Riesz normalization are exact transformations of one signed family.}
{IV Theorem 9.1, ``Parameter-shift equation,'' and Chapter 10, ``The Borel--Riesz bridge.'' Real $z\in[-0.6,0.8]$, $x\in[0,8]$; bounds for number support and EGF degree are recorded.}

\lesson{PAPER IV / 12}{State the signed Pad\'e contract precisely}{IV12_signed_positive_pade}
{\begin{gathered}\Delta_m\ne0\ \Longrightarrow\ \mathcal C-P/Q=O(z^{2m}),\quad \deg P\le m-1,\ \deg Q\le m,\\ \deg Q=m\ \text{and simple zeros}\ \Longrightarrow\ m\text{ simple packets}.\end{gathered}}
{The node panel contrasts positive and signed effective models. Positive nodes remain in the physical interval; signs can produce a node outside it. The moment panel shades exactly the contracted indices. Later moments are predictions and can differ. Neither numerical nonbreakdown at one order nor this finite contract establishes universal stability.}
{The first signed node is $c_1/c_0\approx-0.519817755$. It already leaves $[0,1)$. Separately, the order-four models use the first eight principal moments evaluated with 90 digits. Their first eight reconstructed moments agree at that precision. Moments eight through eleven have nonzero discrepancies. Full denominator degree and separated simple roots are checked for these displayed models.}
{Hankel nonbreakdown gives rational matching. A simple-packet expansion also needs full denominator degree and simple zeros.}
{IV Theorem 11.1, ``Finite signed Pad\'e reduction,'' with the consolidation correction, and Worked example 11.2, ``The first signed packet leaves the physical interval.'' Numerical examples; plotted errors are floored at $10^{-90}$.}

\startvolume{Paper V}{The log-Gamma phase-space atlas}

\lesson{PAPER V / V01}{Normalize the log-Gamma window}{V01_log_gamma}
{g_\alpha(y)=\Gamma(\alpha)^{-1/2}e^{\alpha y/2-e^y/2},\qquad
 \int_{\mathbb R}|g_\alpha(y)|^2dy=1}
{Read amplitude on the left and probability density on the right. A taller curve
need not have more total probability. Changing the positive shape moves the
active log position and changes its width, while the full squared integral
stays one. The curves are unrenormalized crops: the shape $0.5$ window has a
visible left tail beyond the page.}
{For $\alpha=1.5$, substitute $u=e^y$. The norm becomes
$\Gamma(1.5)^{-1}\int_0^\infty u^{0.5}e^{-u}du=1$.
The mode is $y=\log(1.5)$, whereas the mean is
$\psi(1.5)\approx0.036490$. Those two positions differ because the
log-Gamma law is asymmetric. Its variance is approximately $0.934802$.}
{Use the normalized window as the common input to every later transform.}
{V Definition 2.1 and Proposition 2.2, source lines 339 and 369. Shapes
$0.5,1.5,4$, $y\in[-7,3]$. Exact Gamma-CDF crop probabilities are recorded in
the evidence ledger; the displayed crop is not the full norm.}

\lesson{PAPER V / V02}{Build an ambiguity integral}{V02_ambiguity_integral}
{A_\alpha(h,\eta)=\int g_\alpha(y+h/2)g_\alpha(y-h/2)e^{i\eta y}dy
 =\frac{\Gamma(\alpha+i\eta)}{\Gamma(\alpha)}
 \operatorname{sech}^{\alpha+i\eta}(h/2)}
{First multiply the shifted windows. Their product localizes the integral
before modulation is applied. The second panel separates the cosine and sine
contributions to its complex value. Negative real or imaginary contributions
are phases of the integrand, not negative probability weights. Integrating both
components yields one centered ambiguity value.}
{Take $\alpha=1.5$, $h=1.4$, and $\eta=0.8$. The closed formula gives
approximately $0.536116-0.048326i$. Direct quadrature of the displayed
definition agrees. At zero modulation the sine contribution vanishes; at zero
displacement the two windows coincide and the remaining integral is the
log-Gamma characteristic function. Neither operation changes the Fourier
convention.}
{Translation and modulation are joined by one precisely centered integral.}
{V Theorem 3.1, ``Exact ambiguity formula,'' source line 404. Quadrature uses
$y\in[-30,6]$, tolerance $2\times10^{-13}$; analytic omitted-amplitude bounds
and numerical error estimates are recorded separately.}

\lesson{PAPER V / V03}{Separate magnitude from wrapped phase}{V03_phase_space}
{A_\alpha(h,\eta)=\frac{\Gamma(\alpha+i\eta)}{\Gamma(\alpha)}
 e^{-(\alpha+i\eta)\log\cosh(h/2)}\ne0
 \quad(h,\eta\in\mathbb R)}
{The left panel measures overlap strength. The right panel records the
argument modulo $2\pi$, so its sharp color transitions may be phase wraps.
They do not indicate a zero or discontinuity of the complex function.
Increasing displacement suppresses magnitude through the hyperbolic factor.
Increasing modulation introduces the separate Gamma envelope and its phase.}
{At the origin the value is exactly one. At $h=\log4$ and $\eta=0$, its
value is $(4/5)^{3/2}\approx0.715542$. At nonzero modulation, use the same
formula with a complex exponent. The real logarithm of $\cosh(h/2)$ fixes
that exponent without a branch ambiguity on the real plane. The Gamma
function has no zeros, which supplies the nonvanishing proof.}
{A phase wrap is a plotting convention; a zero is a mathematical event.}
{V Theorem 3.1 and Corollary 3.2, source lines 404 and 447. Shape $1.5$;
101-by-101 samples on $[-5,5]\times[-4,4]$. A positive sampled minimum is
finite evidence, not the universal nonvanishing argument.}

\lesson{PAPER V / V04}{Raise shape and read its differential equation}{V04_shift_differential}
{\begin{aligned}
 A_{\alpha+1}&=\frac{\alpha+i\eta}{\alpha}
 \operatorname{sech}(h/2)A_\alpha,\qquad \nu=\alpha+i\eta,\\
 [-4\partial_h^2-\nu(\nu+1)\operatorname{sech}^2(h/2)]A_\alpha
 &=-\nu^2A_\alpha.
 \end{aligned}}
{The first panel compares two constructions of the raised shape: direct
evaluation and the Gamma recurrence. The second shows the potential coefficient
in the differential equation. Both its real and imaginary parts matter when
modulation is nonzero. This complex equation should not be silently read as a
self-adjoint Hamiltonian with a real energy.}
{For $\alpha=1.5$ and $\eta=0.8$, raising shape multiplies the kernel by
$(1.5+0.8i)\operatorname{sech}(h/2)/1.5$. At $h=1.4$, the result is
approximately $0.447661+0.189299i$. Differentiate
$\operatorname{sech}^{\nu}(h/2)$ twice to obtain the displayed equation.
The Gamma prefactor is independent of $h$, so it survives unchanged in
this derivative calculation.}
{One recurrence and one differential equation describe the same kernel.}
{V Proposition 3.3 and Theorem 3.4, source lines 466 and 491.
$h\in[-5,5]$; derivatives are evaluated analytically. The residual checks
floating-point evaluation of the identity, not a discretized spectral problem.}

\lesson{PAPER V / V05}{Transform the window into a Gamma envelope}{V05_fourier_gamma}
{\widehat g_\alpha(\omega)=
 \frac{2^{\alpha/2-i\omega}\Gamma(\alpha/2-i\omega)}
 {\sqrt{2\pi\Gamma(\alpha)}},\qquad
 \rho_\alpha(\omega)=|\widehat g_\alpha(\omega)|^2}
{Keep the unitary convention
$\widehat f=(2\pi)^{-1/2}\int f(y)e^{-i\omega y}dy$ beside the plot.
The Fourier amplitude has real and imaginary parts, even though the starting
window is real. Squaring its magnitude removes the carrier phase and gives
the positive Meixner density. The second panel compares these equivalent
density constructions.}
{Set $u=e^y$. The transform becomes the Gamma integral with exponent
$\alpha/2-i\omega$ and scale $1/2$. For $\alpha=1.5$, evaluating this
integral at frequencies $-2,0,2$ agrees with the closed envelope. Conjugating
the value at $2$ gives the value at $-2$, as required for a real position
window. The positive density is consequently symmetric in frequency.}
{Preserve the Fourier sign before comparing phases or polynomial bases.}
{V Theorem 4.1, source line 540. Frequency crop $[-5,5]$. Independent
quadrature uses $y\in[-30,6]$ with explicit amplitude-tail bounds. The
ledger separates the quadrature estimate from the omitted integral.}

\lesson{PAPER V / V06}{Read depth overlap as a characteristic function}{V06_meixner_overlap}
{\rho_\alpha(\omega)=\frac{2^{\alpha-1}}{\pi\Gamma(\alpha)}
 |\Gamma(\alpha/2+i\omega)|^2,\qquad
 \int\rho_\alpha(\omega)e^{ih\omega}d\omega
 =\operatorname{sech}^{\alpha}(h/2)}
{The two panels express one probability law in frequency and characteristic
coordinates. Larger shape broadens the frequency density because its variance
is $\alpha/4$. Its characteristic function then narrows in displacement.
The right curve is also the zero-modulation ambiguity slice, so it has the
same depth-overlap interpretation as the arithmetic limit.}
{Take $\alpha=1.5$ and a four-to-one depth ratio. Then $h=\log4$ and
$\cosh(h/2)=5/4$. The exact overlap is
$(4/5)^{3/2}\approx0.715542$. This is a transform identity for the
continuous window. A finite squarefree arithmetic packet approaches it under
the compact ratio and time hypotheses of V Theorem 1.1; the identity alone
does not give a finite-depth error bound.}
{The arithmetic scale overlap becomes an exact Meixner characteristic law.}
{V Corollary 4.2, source line 582. Shapes $0.5,1.5,4$. Density crop
$[-5,5]$ is not renormalized. Full-line adaptive cosine quadrature checks
$h=0,\log4$ with declared numerical tolerance.}

\lesson{PAPER V / V07}{Resolve the negative Wigner lobes}{V07_wigner_negative}
{W_\alpha(y,\omega)=\frac1{2\pi}\int
 g_\alpha(y+h/2)g_\alpha(y-h/2)e^{-i\omega h}dh
 =\frac{2e^{\alpha y}}{\pi\Gamma(\alpha)}K_{2i\omega}(e^y)}
{The left panel shows three exact position slices. The right enlarges the
small lobes around zero. A negative value is meaningful here: the Wigner
function records phase-space interference and is not a pointwise probability
density. Its frequency symmetry does not remove those negative regions.
Changing position changes where the oscillatory Bessel lobes appear.}
{For $\alpha=1.5$, $y=0$, and $\omega=2$, the exact expression is
$2K_{4i}(1)/(\pi\Gamma(1.5))\approx-0.00155215$.
The integral remains real because its sine part is odd in the displacement
variable. Negative phase-space values therefore coexist with real-valued
Wigner slices. Integrating over an entire coordinate produces positive
marginals, which is the next lesson's separate statement.}
{Wigner negativity and positive probability marginals are compatible.}
{V Theorem 5.1, source line 615. Shape $1.5$; $y=-1,0,1$ and 241
frequency samples on $[0,4]$. Modified Bessel values use 60-digit working
precision; smooth lines connect finite samples.}

\lesson{PAPER V / V08}{Project Wigner onto its probability marginals}{V08_wigner_marginals}
{\int_{\mathbb R}W_\alpha(y,\omega)d\omega=|g_\alpha(y)|^2,
 \qquad
 \int_{\mathbb R}W_\alpha(y,\omega)dy=\rho_\alpha(\omega)}
{Integrating over frequency leaves the position density; integrating over
position leaves the frequency density. The figure draws these proved exact
marginals directly. It does not integrate a cropped Wigner image, which would
discard part of an oscillatory integral. The shaded regions show only the
probability mass inside the visible windows.}
{At shape $1.5$, the position mean is
$\psi(1.5)\approx0.036490$ and the frequency mean is zero. The position
variance is $\psi_1(1.5)\approx0.934802$, while the frequency variance is
$1.5/4=0.375$. Thus the two marginal widths differ. The negative value from
the previous lesson disappears only after complete marginal integration,
not because the Wigner function has become pointwise positive.}
{A projection can be positive even when the joint phase-space representation is not.}
{V Corollary 5.2, source line 649. Position crop $[-7,3]$, frequency crop
$[-5,5]$. Gamma-CDF position mass and adaptively integrated frequency mass
are recorded; no truncated-image marginal identity is asserted.}

\lesson{PAPER V / V09}{Compare uncertainty with orbit geometry}{V09_uncertainty_metric}
{\Delta Y\Delta P=\tfrac12\sqrt{\alpha\psi_1(\alpha)}>\tfrac12,
 \qquad g_{uu}=\alpha/4,\qquad
 g_{\eta\eta}=\psi_1(\alpha)}
{Follow the uncertainty curve toward the dashed lower bound. It approaches
that bound with increasing shape but remains strictly above it at every
finite shape. The second panel reads the same position and momentum variances
as metric coefficients of a fixed-shape Weyl orbit. Translation sensitivity
increases while modulation sensitivity decreases.}
{At $\alpha=1.5$, the momentum variance is $0.375$ and the position
variance is approximately $0.934802$. Their square-root product is
$0.592073$, greater than $1/2$. The quantum Fisher coefficients in these
two orbit directions are four times the metric, namely $1.5$ and
approximately $3.739209$. These are local parameter sensitivities, not a
claim that the packet is Gaussian at finite shape.}
{Uncertainty and geometric sensitivity use the same exact packet variances.}
{V Theorem 6.1 and Proposition 6.2, source lines 678 and 748.
Shapes $0.2\leq\alpha\leq30$. This is the translation/modulation metric
at fixed shape; mixed coefficients involving changing shape are not inferred.}

\lesson{PAPER V / V10}{Create orthogonal Laguerre packets}{V10_laguerre_gram}
{\Lambda_{n,\alpha}(y)=\sqrt{\frac{n!}{\Gamma(n+\alpha)}}
 e^{\alpha y/2-e^y/2}L_n^{(\alpha-1)}(e^y),\qquad
 \int\Lambda_{n,\alpha}\Lambda_{m,\alpha}dy=\delta_{nm}}
{Each degree adds oscillations around the same envelope. A point where two
curves cross does not determine their inner product. Orthogonality is the
weighted integral over the entire position line. The second panel displays
the degree-zero through degree-four Gram matrix, computed in the transformed
Gamma coordinate rather than by integrating the cropped curves.}
{Use $u=e^y$. The inner product becomes a generalized Laguerre integral
with weight $u^{\alpha-1}e^{-u}$. At $\alpha=1.5$, products of the shown
polynomials have degree at most eight. A 48-node generalized Gauss--Laguerre
rule integrates such polynomial products exactly in mathematical arithmetic.
Its small recorded residual measures floating-point implementation, whereas
the full orthogonality theorem covers all degrees.}
{A common envelope becomes an orthogonal basis through its natural weight.}
{V Theorem 7.1, source line 803. Shape $1.5$, degrees $0$ through $4$;
position curves cropped to $[-5,3]$. The Gram uses 48 weighted nodes,
whose exact polynomial range extends through degree 95.}

\lesson{PAPER V / V11}{Keep Fourier phases in the common ladder}{V11_fourier_ladder}
{\begin{aligned}
 \widehat\Lambda_{n,\alpha}(\omega)&=i^n\widehat g_\alpha(\omega)
 p_{n,\alpha}(\omega),\\
 J_L&=2K_0-K_+-K_-,\qquad J_M=(K_++K_-)/2.
 \end{aligned}}
{The first panel retains the degree phase $i^n$. Dropping it would replace
the actual Fourier-transformed position basis with a different frequency
basis. The second compares neighbor weights of the two real polynomial
recurrences. Their coefficients come from one abstract ladder, although
their signs and scaling differ between coordinate realizations.}
{At shape $1.5$, the first raising coefficient is
$\sqrt{1.5}$. The Laguerre neighbor has its negative, while the real
Meixner polynomial recurrence uses half that magnitude. Test
$[K_-,K_+]=2K_0$ on a finite-support basis vector below the truncation
boundary. A six-dimensional cut satisfies the interior calculation but has
a top-row defect because its last raising operation was removed.}
{Coordinate changes retain degree phases and infinite-basis boundary conditions.}
{V Theorems 8.1 and 9.1, source lines 855 and 917. Shape $1.5$; four
Fourier packet curves and a six-dimensional ladder cut. The audit records
both the small interior residual and the nonzero top-boundary defect.}

\lesson{PAPER V / V12}{Extend the arithmetic depth continuously}{V12_fractional_depth}
{C_{a,\beta}(r,0)=\sum_{j\geq0}(-1)^j\binom rj
 \zeta(a+\beta j)^{-1}
 =\sum_{j\geq0}(-1)^j\binom rj
 [\zeta(a+\beta j)^{-1}-1],\quad r>0}
{The left curve embeds the integer signed orbit in a smooth positive-depth
parameter. The right measures differences from a finite $J=80$ reference for
two equivalent infinite binomial coordinates.
The original finite partial sum retains an artificial remainder from the
unit packet. Subtracting one removes that packet and accelerates the displayed
finite reconstruction. These are signed amplitudes, not positive densities.}
{Choose arithmetic parameters $a=\beta=2$ and depth $r=2.5$. The direct
sieve sums $\mu(n)n^{-2}(1-n^{-2})^{2.5}$ through $N=200000$.
For $r>0$, the full generalized binomial coefficient sum is zero, which
justifies the centered reciprocal-zeta series. At integer depths that series
terminates. Increasing its last index improves the plotted comparison but
does not constitute a universal infinite-series error certificate.}
{Continuous arithmetic depth preserves the unit-packet boundary and signed character.}
{V Theorems 10.1 and 10.2, source lines 995 and 1020. $a>1$ is required
for the stated absolute exchanges. Direct cutoff tail is at most $1/N$;
the reference is an 80-index binomial sum at 60-digit precision. The log plot
floors differences at $10^{-18}$.}

\startvolume{Paper VI}{Shape composition and transfer}
\lesson{PAPER VI / 01}{Read shape through the logarithmic coordinate}{VI01_normalized_shape}
{g_\alpha(y)=\Gamma(\alpha)^{-1/2}e^{(\alpha y-e^y)/2},\qquad |g_\alpha(y)|^2\,dy=\frac{u^{\alpha-1}e^{-u}}{\Gamma(\alpha)}\,du}
{The left panel shows probability densities, the squared packet amplitudes. The right shows the same laws after $u=e^y$. Each curve has total mass one on its full domain. The logarithmic Jacobian changes the curve and its mode. Shape therefore does more than move a fixed profile along an axis.}
{For shape two, the log density is $e^{2y-e^y}$. Its derivative vanishes at $y=\log2$. The corresponding Gamma density is $u e^{-u}$, whose mode is $u=1$. The mean of the log coordinate is $1-\gamma_{\mathrm E}$, about $0.4228$, where $\gamma_{\mathrm E}$ is Euler's constant; its variance is $\pi^2/6-1$.}
{A change of coordinate preserves probability after including its Jacobian. Packet normalization uses the full real line.}
{VI Proposition 1.1, ``Normalization.'' Shapes $1,2,4,8$; finite plot windows. Omitted mass is recorded using regularized Gamma functions; no plotted curve is renormalized.}

\lesson{PAPER VI / 02}{Separate shape affinity from local curvature}{VI02_shape_affinity}
{F(\alpha,\gamma)=\frac{\Gamma((\alpha+\gamma)/2)}{\sqrt{\Gamma(\alpha)\Gamma(\gamma)}},\qquad F(\alpha,\alpha+d)=1-\frac{\psi_1(\alpha)}8d^2+O(d^3)}
{The first panel compares every displayed second shape with the fixed shape two. Only the equal-shape point reaches one. The second enlarges a neighborhood of that point and compares the exact affinity with its quadratic expansion. Their separation away from zero shows why a local expansion should not be used as a global identity.}
{For shapes two and four, $F=2/\sqrt6=\sqrt{2/3}$. Thus the squared overlap at zero translation and modulation is exactly $2/3$. At increment $d=0.2$, the exact affinity is about $0.996969$; the quadratic approximation gives $0.996775$. Its curvature comes from $\psi_1(2)=\pi^2/6-1$.}
{Origin affinity compares shapes at one common coordinate. A later translation can increase their overlap.}
{VI Proposition 2.2, ``Strict shape affinity,'' and its following local expansion. The local quadratic is displayed as an approximation; the shape metric coefficient is $\psi_1(\alpha)/4$.}

\lesson{PAPER VI / 03}{Keep magnitude and phase in cross ambiguity}{VI03_cross_ambiguity}
{\mathcal A_{\alpha,\gamma}(h,\eta)=e^{(\alpha-\gamma)h/4}\frac{\Gamma(\lambda+i\eta)}{\sqrt{\Gamma(\alpha)\Gamma(\gamma)}}\operatorname{sech}^{\lambda+i\eta}(h/2),\quad\lambda=(\alpha+\gamma)/2}
{The two panels represent one complex overlap. Magnitude measures localization; the principal phase records information that a squared magnitude removes. A phase-color jump is a branch wrap, not a zero. Unequal shapes shift the magnitude peak away from zero displacement, while the frequency envelope remains symmetric. The white cross marks the zero-frequency displacement maximum.}
{Take $\alpha=2$, $\gamma=4$. At zero frequency, the magnitude is $\sqrt{2/3}\,e^{-h/2}\operatorname{sech}^3(h/2)$. Differentiating its logarithm gives the maximum at $h=-\log2$, with value about $0.967700$. The origin value is only $\sqrt{2/3}$. At $(h,\eta)=(0,1)$ the squared magnitude is $5\pi/(3\sinh\pi)$.}
{Translation, modulation, and shape are distinct coordinates. The complex Gamma factor retains their phase relation.}
{VI Theorem 2.1, ``Cross-shape Gamma--sech formula.'' Shapes $2,4$ on $[-3,3]^2$; exact-formula samples. Independent overlap quadrature is checked separately.}

\lesson{PAPER VI / 04}{Factor energy into two independent laws}{VI04_factorization}
{\frac{|\mathcal A_{\alpha,\gamma}(h,\eta)|^2}{2\pi}=p_{\alpha,\gamma}(h)\rho_{\alpha+\gamma}(\eta),\qquad \rho_s(\eta)=\frac{2^{s-1}|\Gamma(s/2+i\eta)|^2}{\pi\Gamma(s)}}
{The displacement panel reflects when the two shapes are exchanged. The frequency panel uses only their sum. These are probability marginals of the normalized cross-spectrogram, whose product is its full energy density. Independence here belongs to that product probability law. It does not say that the complex ambiguity factors into two real densities or loses its phase.}
{With shapes two and four, displacement is the logit of a Beta$(2,4)$ variable. Its mean is $\psi(2)-\psi(4)=-5/6$, while the symmetric shape-six frequency has mean zero and variance $3/2$. At the origin, $p_{2,4}(0)=5/16$ and $\rho_6(0)=16/(15\pi)$, so the joint density is exactly $1/(3\pi)$.}
{A positive two-dimensional energy law can have independent marginals even when its underlying amplitude has nontrivial phase.}
{VI Theorem 3.1, ``Exact cross-spectrogram factorization.'' Position tails use Beta CDFs; frequency tails use a stated Meixner Chernoff bound. Finite-window mass is not reset to one.}

\lesson{PAPER VI / 05}{Add shapes by convolution}{VI05_shape_convolution}
{\rho_{s+t}=\rho_s*\rho_t,\qquad \rho_{s+2}(x)=\frac{s^2+4x^2}{s(s+1)}\rho_s(x)}
{The first panel overlays a closed density with an independently integrated convolution. The second uses a different operation: pointwise quadratic multiplication raises shape by two. Both identities change shape, but their inputs and computational costs differ. The convolution adds independent random variables. The ladder changes one density locally at each displayed frequency.}
{Shape one has density $\operatorname{sech}(\pi x)$. Convolving two such laws gives shape two, $2x/\sinh(\pi x)$, with the removable origin value $2/\pi$. Its variance is $1/4+1/4=1/2$. Separately, multiplying shape one by $(1+4x^2)/2$ produces shape three. At $x=1$, shape two has density about $0.173179$.}
{The characteristic function $\operatorname{sech}^s(h/2)$ turns shape addition into multiplication. The Gamma recurrence gives the ladder.}
{VI Theorem 4.1, ``Meixner shape semigroup,'' and Proposition 4.2, ``Two-step density ladder.'' Convolution integrates on $[-8,8]$; its omitted integral is uniformly bounded by $1.55\times10^{-11}$.}

\lesson{PAPER VI / 06}{Continuous shape produces jumps and cumulants}{VI06_levy_bernoulli}
{-\log\cosh(\xi/2)=\int_0^\infty\frac{\cos(\xi x)-1}{x\sinh(\pi x)}\,dx,\qquad \kappa_{2m}=s\frac{(2^{2m}-1)|B_{2m}|}{2m}}
{The generator panel compares its exact Fourier symbol with the L\'evy integral. The numerator cancels the apparent origin singularity, so the combined integrand is finite there. The second panel tracks even cumulants through shape. The second and sixth cumulants overlap because both equal $s/4$; their coincidence does not make the distribution Gaussian.}
{At shape four, variance is one, the fourth cumulant is $1/2$, and the sixth is one. A Gaussian with variance one would have both higher cumulants zero. The nonzero values measure departures from Gaussianity. For dual frequency $\xi=2$, the unit-shape generator symbol is $-\log\cosh1$, about $-0.433781$.}
{The L\'evy measure describes infinitesimal shape increments; Bernoulli numbers encode the complete even-cumulant sequence.}
{VI Theorems 5.1 and 5.2, ``L\'evy--Khintchine representation'' and ``Bernoulli cumulant law.'' Quadrature uses $[0,10]$, stable sine-squared cancellation, and a tail bound below $2.90\times10^{-15}$.}

\lesson{PAPER VI / 07}{Split one polynomial without losing energy}{VI07_two_child_refinement}
{p_{n,\alpha+\gamma}(x+y)=\sum_{r=0}^n C_{n,r}^{\alpha,\gamma}p_{r,\alpha}(x)p_{n-r,\gamma}(y),\qquad (C_{n,r}^{\alpha,\gamma})^2=\binom nr\frac{(\alpha)_r(\gamma)_{n-r}}{(\alpha+\gamma)_n}}
{Every bar is an allowed degree allocation between two children. Their heights sum to one, so the transfer preserves coefficient energy. The polynomial panel overlays the parent with the sum of all five products. Individual products can be large or signed; the probability bars describe squared transfer coefficients, not pointwise polynomial contributions.}
{Use degree four and child shapes two and three. The energy weights, in order $r=0,\ldots,4$, are $(15,20,18,12,5)/70$. Their sum is one and their mean allocation to the first child is $8/5$. At $x=0.6$, $y=0.85$, the normalized parent is approximately $-0.716400$, reproduced by the complete five-term child sum.}
{An isometric refinement is a complete change of polynomial coordinates. Omitting an allocation changes the reconstruction.}
{VI Theorem 6.1, ``Two-child packet refinement.'' Degree four, shapes $2,3$, fixed $y=0.85$; all five terms are retained. Numerical pointwise reconstruction and coefficient normalization are checked.}

\lesson{PAPER VI / 08}{Choose a refinement tree freely}{VI08_coassociative_refinement}
{C_{n,\mathbf r}^{\boldsymbol\alpha}=\sqrt{\frac{n!\prod_j(\alpha_j)_{r_j}}{(\sum_j\alpha_j)_n\prod_jr_j!}},\qquad C_{n,r_1+r_2}^{\alpha_1+\alpha_2,\alpha_3}C_{r_1+r_2,r_1}^{\alpha_1,\alpha_2}=C_{n,\mathbf r}^{\boldsymbol\alpha}}
{The triangle contains every degree-four allocation among three children. Color encodes energy; marker area increases with it. The second panel lists all fifteen terminal amplitudes and overlays two binary refinement orders. Coincident markers express coassociativity. The geometric arrangement is a degree-allocation simplex, not the phase space used by the ambiguity function.}
{Take shapes $(2,3,1)$ and allocation $(1,2,1)$. The direct squared coefficient is $2/21$. Splitting the first two children first gives squared factors $5/18$ and $12/35$. Splitting the last two first gives $20/63$ and $3/10$. Both products equal $2/21$, and both amplitudes are the positive square root, about $0.308607$.}
{Intermediate shape factors cancel exactly. A deeper binary tree gives the same final many-child coefficients.}
{VI Theorem 7.1, ``Multinomial packet refinement,'' and Corollary 7.2, ``Coassociativity.'' All fifteen compositions of degree four are shown; both trees and total energy are checked.}

\lesson{PAPER VI / 09}{Read conditional transfer from its singular values}{VI09_martingale_spectrum}
{T_{s,t}p_{n,s+t}=\sigma_n(s,t)p_{n,s},\qquad \sigma_n(s,t)=\sqrt{\frac{(s)_n}{(s+t)_n}},\quad \sigma_n(s,t)\sim\sqrt{\frac{\Gamma(s+t)}{\Gamma(s)}}n^{-t/2}}
{Each spectral curve measures contraction of one polynomial degree. Larger increments suppress high degrees more strongly. Degree zero, excluded from the logarithmic horizontal axis, remains exactly one. The polynomial panel shows why orthonormal packets contract even though the standard, unnormalized Meixner--Pollaczek polynomials satisfy the martingale identity without a scalar loss.}
{For $s=2$, $t=3$, the first singular values are $1$, $\sqrt{2/5}$, $1/\sqrt5$, and $2/\sqrt{35}$. Since a shape-three increment has variance $3/4$, conditioning $(2(X+Y)^2-5/2)/\sqrt{15}$ gives $(2X^2-1)/\sqrt{15}$. This equals $p_{2,2}(X)/\sqrt5$, and at $X=1$ it is $1/\sqrt{15}$.}
{The exact cocycle permits successive increments. Compactness holds for every positive increment; Schatten membership requires $pt>2$.}
{VI Theorem 8.1, ``Polynomial martingale law,'' and Proposition 8.2, ``Cocycle and compactness.'' Degrees through 60; a rank-61 truncation at increment three has exact omitted norm $\sigma_{61}\approx0.00957$.}

\lesson{PAPER VI / 10}{A continuous coherent orbit resolves identity}{VI10_coherent_frame}
{\psi_{\alpha,u,\eta}(y)=e^{i\eta(y-u/2)}g_\alpha(y-u),\qquad \frac1{2\pi}\int_{\mathbb R^2}|\langle\psi_{\alpha,u,\eta},f\rangle|^2\,du\,d\eta=\|f\|_2^2}
{Translation moves the window; modulation oscillates its complex amplitude. Only real parts are drawn on the left. The right integrates the coherent-transform energy of $f=g_2$ over a position strip, including every modulation frequency. Its approach to one illustrates the continuous frame normalization. Three displayed states alone do not form the identity resolution.}
{For $f=g_2$, the frame energy density is $p_{2,2}(u)\rho_4(\eta)$. Integrating frequency removes its normalized factor. In the strip $|u|\le\log3$, the logistic endpoints are $1/4$ and $3/4$. The Beta$(2,2)$ CDF is $3v^2-2v^3$, so the captured energy is exactly $11/16$.}
{An overlap kernel becomes a coherent transform and a reproducing kernel when the full normalized orbit is used.}
{VI Theorems 9.1 and 9.2, ``Cross-shape coherent overlap'' and ``Continuous tight frame.'' Shape two; strip energy uses an exact Beta CDF with frequency integrated on all $\mathbb R$.}

\lesson{PAPER VI / 11}{Positive smoothing can have unstable inversion}{VI11_berezin_inverse}
{\widehat{\mathcal B}_\alpha(\xi,x)=\frac{|\Gamma(\alpha+i\xi)|^2}{\Gamma(\alpha)^2}\operatorname{sech}^{2\alpha}(x/2)>0,\qquad \widetilde F=\mathcal B_\alpha*F}
{The multiplier stays positive while shrinking quickly along either dual coordinate. The bars show its action on four cosine modes: convolution preserves each mode and multiplies its amplitude by the displayed number. Therefore exact information survives algebraically, but recovering a small attenuated mode requires amplifying any error attached to it.}
{At shape two, $\widehat{\mathcal B}_2(1,0)=2\pi/\sinh\pi\approx0.544058$. The mode $(3,0)$ retains only $30\pi/\sinh(3\pi)\approx0.015212$, requiring inverse amplification about $65.7$. For mode $(1,2)$, the additional factor is $\operatorname{sech}^4(1)$, reducing its amplitude to about $0.095960$. The constant mode retains amplitude one.}
{Injectivity is an exact uniqueness statement. Stable numerical recovery requires an additional regularization or restricted function class.}
{VI Theorem 10.1, ``Exact Berezin factorization,'' and Corollary 10.2, ``Injectivity and ill-posed inversion.'' Bounded cosine symbols illustrate attenuation; the $L^2$ inverse statement is the source theorem.}

\lesson{PAPER VI / 12}{Rescale shape toward a Gaussian heat kernel}{VI12_gaussian_heat}
{Z_\alpha=\sqrt{\alpha/2}\,H_\alpha,\quad W_\alpha=\sqrt{2/\alpha}\,\Omega_\alpha\ \Longrightarrow\ N(0,1),\qquad S_\alpha F\longrightarrow e^{\Delta/2}F}
{The two panels use different rescalings because unscaled displacement narrows while frequency broadens. Both normalized laws approach the same standard normal curve. At finite shape, the position variance is still slightly above one; frequency variance is exactly one. Matching a limiting curve does not remove the higher finite-shape cumulants or establish a function-space convergence theorem.}
{At shape two, $\operatorname{Var}(Z_2)=2\psi_1(2)\approx1.289868$, while $\operatorname{Var}(W_2)=1$. At shape 32, the first variance is about $1.015788$. The rescaled two-dimensional multiplier at dual mode $(1,1)$ is then about $0.365956$, approaching the heat multiplier $e^{-1}$. For a Gaussian test symbol $e^{-|z|^2/2}$, the limiting smoothed value is $\tfrac12e^{-|z|^2/4}$.}
{The natural rescaling gives independent standard-normal marginals and the two-dimensional heat operator as the large-shape limit.}
{VI Theorem 11.1, ``Gaussian marginal limits,'' and Corollary 11.2, ``Heat-operator limit.'' Exact finite-shape densities; tail bounds recorded. All-Schwartz-seminorm convergence belongs to the source proof, not the plotted comparison.}


\startvolume{Paper VII}{Classical recoupling packets}

\lesson{Paper VII / 01}
{A ladder begins with a positive shape}
{vii_01_ladder}
{K_0|n\rangle=(n+\kappa)|n\rangle,\qquad K_+|n\rangle=\sqrt{(n+1)(n+2\kappa)}\,|n+1\rangle}
{Read the left panel vertically: an occupation label selects a weight eigenvalue. Changing the positive shape shifts that entire ladder. The right panel reads the same labels as raising amplitudes. Those coefficients determine how neighboring degrees are joined, rather than merely naming them. The curves are computed from the defining representation, with no fitted model. These two pieces of data are the input to every tensor construction in this volume.}
{Set $\kappa=1/2$ and start at degree $n=2$. Its weight is $2.5$, while its raising coefficient is $\sqrt{3\cdot3}=3$. Thus $K_+|2\rangle=3|3\rangle$ in the algebraic ladder. This coefficient is an operator amplitude, not a probability. A normalized physical transition would require its own state and dynamics. The example explains why later Casimir matrices retain products of ladder amplitudes.}
{Keep the degree label, weight eigenvalue, and transition amplitude as three distinct quantities.}
{VII, Chapter 1, Definition 1.1, \emph{Positive shape packet}, source lines 219--243. Direct representation formulas; positive $\kappa$, nonnegative integer degree, finite-support vectors.}

\lesson{Paper VII / 02}
{Project a large excitation simplex}
{vii_02_simplex}
{\dim H_N^{(3)}=\binom{N+2}{2},\qquad \dim M_N^{(3)}=N+1}
{Each point on the left is one occupation triple with the same total excitation. The three digits give its exact address, and color records the third occupation. The plot contains fifteen points because weak compositions count the ambient basis. On the right, compare that full count with the dimension of the new lowest-weight kernel. The smaller curve counts an exact subspace. It does not mean that a few occupation states were discarded by sampling.}
{At $N=4$, the full space has $\binom62=15$ basis states. The lowering map targets degree three, whose full dimension is ten. Its full row rank therefore leaves a five-dimensional kernel. A subsequent recoupling matrix needs $5\times5$ entries, even though its defining projector was constructed in fifteen dimensions. At $N=0$, both spaces have dimension one, so the reduction also respects the initial degree.}
{Compression comes from an operator kernel, with its dimension proved before any numerical atlas is chosen.}
{VII, Chapter 2, Proposition 2.1 and Theorem 2.2. The explicit full-row-rank proof used here is VIII, Theorem 3.1, specialized at $q=1$.}

\lesson{Paper VII / 03}
{See the lowest-weight plane}
{vii_03_lowest_plane}
{M_N=\ker(K_-^{[3]}|_{H_N}),\qquad N=1:\quad x_{100}+x_{010}+x_{001}=0}
{The three-dimensional view is coefficient space. Its axes are amplitudes of three occupation basis states, rather than spatial coordinates of particles. The shaded plane is the exact lowering kernel for three half-shapes at one excitation. Two arrows show an orthonormal atlas inside that plane. The right panel advances one degree and shows the rectangular lowering matrix. The next kernel is again selected by an equation, not by the apparent proximity of points in a picture.}
{At degree one, $(1,1,-2)/\sqrt6$ lies in the plane because its coefficients sum to zero. The vector $(1,-1,0)/\sqrt2$ also lies there, and their inner product vanishes. At degree two, the displayed $3\times6$ lowering map has a three-dimensional nullspace. The generated frame satisfies its lowering equation to floating-point roundoff. This residual checks this frame; the rank argument proves the general dimension.}
{A basis is a choice inside the kernel; the kernel is the mathematical object being represented.}
{VII, Chapter 2, \emph{Fixed excitation and the lowest-weight projector}, lines 259--343; Chapter 5, the one-excitation construction. Newly computed plane and matrix.}

\lesson{Paper VII / 04}
{Two observables select two atlases}
{vii_04_pair_atlases}
{\lambda_j^{ab}=(\kappa_a+\kappa_b+j)(\kappa_a+\kappa_b+j-1),\qquad \mathcal R=L^{\mathsf T}R}
{The left curves are the spectra of two intermediate Casimirs. Their labels record how two neighboring shape factors are coupled. The right panel is the overlap of their complete eigenbases in the same lowest-weight space. A spectrum gives the addresses of an atlas, while the overlap tells how coordinates move between atlases. Eigenvalues alone do not determine that movement. The signed entries retain the phase convention of the selected real eigenvectors.}
{Use shapes $(0.7,1.1,0.9)$ and $N=4$. The left pair begins at $1.8(0.8)=1.44$, and the right pair begins at $2(1)=2$. Both spectra contain five entries because the multiplicity dimension is five. The program restricts each Casimir, orders its eigenvalues, and computes the displayed overlap. It separately compares those eigenvalues with the formula and checks orthogonality. These are complementary numerical checks of the chosen example.}
{The shared space supports reconstruction; the different intermediate observables support different coordinate descriptions.}
{VII, Chapter 3, Proposition 3.1; Chapter 4, Definition 4.1 and Theorem 4.2. Pair-spectrum identification is a classical representation-theoretic input.}

\lesson{Paper VII / 05}
{The middle shape controls a beam splitter}
{vii_05_shape_beam}
{\mathcal R_1=\begin{pmatrix}c&s\\s&-c\end{pmatrix},\qquad c^2=\frac{\kappa_1\kappa_3}{(\kappa_1+\kappa_2)(\kappa_2+\kappa_3)},\quad s^2=1-c^2}
{Hold the outer shapes fixed and move along the middle-shape axis. The first curve is the retained-channel amplitude, and the second is the transferred-channel amplitude in the declared real gauge. Their squares always sum to one. The right panel freezes the curve at one shape triple and displays the full signed transform. Its negative entry is needed for orthogonality. Reading the same numbers as unsigned probabilities would erase the operation that makes exact synthesis possible.}
{For $(\kappa_1,\kappa_2,\kappa_3)=(0.7,1.1,0.9)$, the denominator is $1.8\cdot2=3.6$. Hence $c^2=0.175$ and $s^2=0.825$, giving $c\approx0.41833$ and $s\approx0.90830$. The independent Casimir construction reproduces this matrix. Applying the overlap and then its transpose recovers every two-channel amplitude vector. Squaring entries instead produces a measured channel with a different role.}
{The exact low-degree formula links a continuous shape choice to a finite reversible coordinate transform.}
{VII, Chapter 5, \emph{One excitation as a shape beam splitter}, source lines 452--510. Exact formula checked against the newly computed pair-Casimir overlap.}

\lesson{Paper VII / 06}
{An exact three-channel transform}
{vii_06_exact_three}
{\mathcal R_2=\begin{pmatrix}
1/3&1/\sqrt3&\sqrt5/3\\
1/\sqrt3&1/2&-\sqrt{15}/6\\
\sqrt5/3&-\sqrt{15}/6&1/6
\end{pmatrix},\qquad P_{jm}=|\mathcal R_2(j,m)|^2}
{Read the left matrix as amplitudes and the right matrix as probabilities after measurement. The two negative amplitudes participate in cancellations, while the probability matrix contains only nonnegative entries. Rows and columns of the amplitude matrix are orthonormal. Rows and columns of the probability matrix sum to one. Both panels describe the same exact half-shape example, but they support different operations. The displayed decimals are visual labels for the exact radicals in the central formula.}
{The first amplitude row has squared norm $1/9+1/3+5/9=1$. Its probability row is therefore $(1/9,1/3,5/9)$. For the first two amplitude rows, positive and negative products cancel, giving inner product zero. The generated Casimir overlap agrees with the complete radical matrix. Checking just the first probability row would confirm one normalization, while the full orthogonality check confirms reversible reconstruction across all three channels.}
{Keep the exact signed matrix available even when the application eventually requests only measured probabilities.}
{VII, Chapter 6, \emph{An exact three-channel example}, source lines 512--559, equal shapes $1/2$ and $N=2$.}

\lesson{Paper VII / 07}
{Phase changes a measured output}
{vii_07_phase_measurement}
{b_m=\sum_j\mathcal R_{jm}a_j,\qquad |b_m|^2\ \text{and}\ \sum_j|\mathcal R_{jm}|^2|a_j|^2\ \text{need not agree}}
{The left bars are the two signed contributions reaching each target channel. The black points are their coherent sums. The right panel compares measuring those sums with dephasing the input first and then using the squared-amplitude channel. Both output laws are normalized, yet they differ because only one retains the cross terms between routes. This is a concrete reason to separate a reversible coordinate transform from a stochastic measurement kernel in the packet record.}
{Start in the exact three-channel example with $a=(1,-1,0)/\sqrt2$. At target label zero, the coherent amplitude is $(1/3-1/\sqrt3)/\sqrt2$. Its squared magnitude is approximately $0.02977$. If the input is dephased first, the corresponding probability is $(1/9+1/3)/2=2/9$. The probability difference comes from interference. It is not a failure of normalization, and the signed matrix itself still has its exact inverse.}
{Amplitude addition occurs before squaring in coherent synthesis; a dephased channel has already lost those cross terms.}
{VII, Chapters 9 and 11, measurement interpretation and Boundary 11.1. Newly worked superposition using the exact Chapter 6 matrix.}

\lesson{Paper VII / 08}
{A positive recurrence needs its own gauge}
{vii_08_recurrence_gauge}
{J=\mathcal R\,\operatorname{diag}(\lambda_m)\mathcal R^{\mathsf T},\qquad S=\operatorname{diag}(s_j):\ (SJS)_{j,j+1}=|J_{j,j+1}|}
{Both panels represent the same opposing Casimir. In the left panel, the overlap convention makes its first row and column positive. The resulting adjacent Casimir entries can be negative. The right panel changes the signs of left-atlas vectors so every adjacent entry is positive. The zero pattern and eigenvalues remain unchanged. This distinction matters when a scalar three-term recurrence is presented as a Jacobi recurrence with positive neighboring coefficients.}
{For three half-shapes at $N=1$, the printed overlap has $J_{01}=-\sqrt3/2$. It cannot simultaneously have the same signed coordinates and a positive adjacent coefficient. Set $s_0=1$, then choose $s_{j+1}=s_j\operatorname{sign}(J_{j,j+1})$. The rephased overlap is $S\mathcal R$, and its recurrence uses $SJS$. The displayed seven-channel example checks the tridiagonal zero pattern and the rephased matrix eigen-equation separately.}
{A gauge changes coordinates; state it explicitly before assigning signs to recurrence coefficients.}
{VII, Chapter 7, Theorem 7.1, read with the corrected row-sign convention. Irreducible tridiagonality is a classical Racah input; the displayed finite matrices are newly checked.}

\lesson{Paper VII / 09}
{A fixed measured channel spreads a law}
{vii_09_entropy}
{P_{jm}=|\mathcal R_{jm}|^2,\qquad H(pP)\geq H(p),\qquad H(p)=-\sum_jp_j\log p_j}
{The bars compare one input law with its output under one fixed squared-overlap matrix. The right panel sorts both distributions before adding their largest entries. The output curve lies below the input curve until both reach total mass one. That is the finite concentration order behind the entropy statement. The comparison is between input and output at a fixed kernel. It does not compare unrelated kernels or promise that entropy rises when a shape parameter is changed.}
{At $N=5$, use shapes $(0.7,1.1,0.9)$ and the input $(0.7,0.1,0.05,0.05,0.05,0.05)$. The code computes its six-channel measured output directly. Each row and column of the kernel sums to one, and the output entropy exceeds the input entropy. Concavity applied with the column weights proves the inequality for every probability law. No random sampling is needed for that theorem.}
{Bistochasticity gives a universal fixed-channel entropy law, while particular mixing rates remain properties of the chosen kernel.}
{VII, Chapter 9, Theorem 9.1; Chapter 10, Theorem 10.1. Natural logarithms give entropy in nats; zero terms use $0\log0=0$.}

\lesson{Paper VII / 10}
{Tree coordinates have joint addresses}
{vii_10_tree_addresses}
{\dim M_N^{(4)}=\binom{N+2}{2},\qquad (((12)3)4):\quad 0\leq j_{12}\leq j_{123}\leq N}
{Each vertex on the left names a binary bracketing of four factors. The two block labels beneath it name the commuting intermediate Casimirs used by that tree. The right panel shows their joint integer addresses for the first tree. These addresses fill a triangular lattice rather than an unconstrained square. An atlas is complete only when all compatible addresses occur in the same final lowest-weight space. The graph explains which coordinate transforms connect neighboring bracketings.}
{At $N=4$, there are fifteen addresses. With shapes $(0.55,0.8,1.05,1.3)$, the address $(j_{12},j_{123})=(1,3)$ gives intermediate shapes $2.35$ and $5.4$. Their ordinary Casimir eigenvalues are $3.1725$ and $23.76$. The numerical joint eigenbasis recovers the full triangular address set and diagonalizes both operators. Dimension, address completeness, and diagonalization each test a different part of this construction.}
{A joint spectrum supplies a finite address system for a tree, rather than just a list of separate eigenvalues.}
{VII, Chapter 12, \emph{Five complete coupling trees} and Proposition 12.1; plotted addresses are extracted from newly computed joint eigenbases.}

\lesson{Paper VII / 11}
{What a pentagon check actually certifies}
{vii_11_loop_scope}
{U_{r,r+1}=B_r^{\mathsf T}B_{r+1},\qquad U_{01}U_{12}U_{23}U_{34}U_{40}=I}
{The left bars show that three constructions all close their common-space overlap loop. The first uses the intended Casimir atlases. The second changes coordinates consistently at every vertex. The third uses a common subspace that fails the intended lowering equation. The right bars reveal that difference. The loop theorem remains exact, but its telescoping proof cannot independently distinguish the intended projector from every consistently shared alternative. The plotted log scale displays roundoff and a large failure together.}
{For four factors at $N=2$, the intended packet dimension is six inside a ten-dimensional excitation space. Five complete orthonormal bases of a common six-dimensional subspace always telescope. This remains true after consistent orthogonal coordinate changes. Choosing six ambient coordinate vectors also produces a closing loop, yet their lowering residual exceeds one. Thus the loop check passes while the separate kernel check rejects the wrong subspace.}
{Pair overlap closure with kernel, dimension, spectrum, and eigen-equation checks when validating an associator implementation.}
{VII, Chapter 13, Theorem 13.1, with its implementation-test interpretation corrected. Newly constructed deterministic negative control, coordinate-change seed 73011.}

\lesson{Paper VII / 12}
{Compile in the reduced dimension}
{vii_12_compiler}
{\mathcal R=G_1^{\mathsf T}\cdots G_L^{\mathsf T}D,\qquad L\leq d(d-1)/2,\qquad D_{jj}\in\{-1,1\}}
{The left panel is the signed diagonal remaining after row rotations eliminate entries below the overlap's diagonal. Reversing those rotations reconstructs the original transform. The right panel compares generated rotation counts with the triangular worst-case bound as the packet dimension grows. Each rotation touches only two coordinate slots. This compiler uses the multiplicity dimension, so its index set is smaller than the original tensor excitation space that defined the projector.}
{At $N=5$ for three factors, the packet dimension is $d=6$ and the bound is fifteen rotations. The program records each Givens matrix, reduces the overlap to diagonal signs, and reconstructs it in reverse order. The residual is checked against the original overlap. This gives a finite algebraic circuit for the example. It does not assert nearest-neighbor hardware locality, minimum circuit depth, or an optimal gate count for a particular device.}
{Reduction defines the working coordinates; rotation factorization supplies an exact reconstruction procedure in those coordinates.}
{VII, Chapter 15, Theorem 15.1, \emph{Givens packet compilation}; Chapter 16, \emph{Projector-resolved Racah transform}. Direct finite linear algebra and new bounded checks.}

\startvolume{Paper VIII}{q-deformation and scalar bridges}

\lesson{Paper VIII / 01}
{Replace a linear number by a q-number}
{viii_01_q_numbers}
{[x]_q=\frac{q^x-q^{-x}}{q-q^{-1}}=\frac{\sinh(x\theta)}{\sinh\theta},\qquad q=e^{-\theta},\quad [x]_1=x}
{The left panel shows the symmetric q-number at three positive deformation choices. The ordinary linear coordinate is recovered at one. The right panel divides by the positive input, making the relative deformation visible. This convention is symmetric under replacing q by its reciprocal, even though the ordered coproduct later retains directional q-weights. A scalar convention must be fixed before formulas from different quantum-algebra sources are compared.}
{At $q=1/2$, $[2]_q=2.5$ and $[3]_q=5.25$, while $[1]_q=1$. Thus the half-shape ladder's raising amplitude at degree one becomes $\sqrt{[2]_q[2]_q}=2.5$. At $q=1$, that coefficient is two. The stable implementation uses the hyperbolic expression near one and the exact continuous value at one. Positive x always gives a positive q-number in the declared real range.}
{Deforming the scalar ladder coefficients preserves the positive shape construction while changing its degree-dependent geometry.}
{VIII, Chapter 1, Definitions 1.1--1.2 and Proposition 1.3. Real range $0<q\leq1$; q=1 is defined by continuity.}

\lesson{Paper VIII / 02}
{The coproduct routes lowering events}
{viii_02_coproduct}
{\Delta K_-=K_-\otimes q^{K_0}+q^{-K_0}\otimes K_-,\qquad (\Delta\otimes I)\Delta=(I\otimes\Delta)\Delta}
{Each curve on the left corresponds to lowering a different occupied factor at total degree one. The labels follow the declared composition order. Their unequal weights come from the q-powers assigned to factors before and after the lowering event. On the right, an explicit three-term expansion agrees with the iterated matrix implementation. This tests coassociative routing in a small transparent case. Coassociativity concerns bracketings of the same ordered factors, rather than freely permuting those factors.}
{At shapes $(0.7,1.1,0.9)$ with $q=0.73$, the $|001\rangle$ route has coefficient $q^{-(0.7+1.1)}\sqrt{[1.8]_q}$. The route from $|010\rangle$ has coefficient $q^{-0.7+0.9}\sqrt{[2.2]_q}$, and the $|100\rangle$ route has $q^{1.1+0.9}\sqrt{[1.4]_q}$. Both coproduct bracketings produce exactly these three terms. Their values also match the freshly built lowering row.}
{Preserve the tensor order and q-number convention when constructing block observables or comparing formulas.}
{VIII, Chapter 2, Theorem 2.1, \emph{Coassociative packet routing}, source lines 293--327. Direct expansion and new degree-one matrix comparison.}

\lesson{Paper VIII / 03}
{A q-Hahn seed keeps signs and weights}
{viii_03_lowest_seed}
{\frac{c_{n+1}}{c_n}=-q^{-(N-1+\kappa_1+\kappa_2)}
\sqrt{\frac{[N-n]_q[2\kappa_2+N-n-1]_q}{[n+1]_q[2\kappa_1+n]_q}}}
{The left panel retains the alternating phases of the one-dimensional lowest-weight seed. The right panel squares those same coefficients and exposes the q-dependent weight tilt. Both come from a solved lowering equation. Changing q alters the distribution without changing the six available occupation labels. The normalized vectors have total probability one. Identifying general coupled coefficients with scalar q-Hahn polynomials is an additional classical representation-theoretic input.}
{Take $N=5$ and shapes $(0.7,1.1)$. Starting with $c_0=1$, apply the displayed ratio five times, then normalize. At $q=0.73$, the resulting vector is annihilated by the two-factor lowering matrix. It also agrees up to its overall sign with a separately computed nullspace vector. The ratio derives the seed analytically; the two numerical checks confirm its implementation at this input. The plotted squared weights are not inferred from a fitted probability model.}
{The signed seed supports synthesis, and its square supports a weight description; neither representation should replace the other silently.}
{VIII, Chapter 4, Theorem 4.1, \emph{Explicit q-Hahn lowest packet}, source lines 407--474. New recurrence and nullspace comparisons at three q choices.}

\lesson{Paper VIII / 04}
{Deform both a spectrum and an atlas}
{viii_04_deformed_atlases}
{\lambda_j^{ab}(q)=[\kappa_a+\kappa_b+j]_q[\kappa_a+\kappa_b+j-1]_q,\qquad \mathcal R_N(q)=L_N(q)^{\mathsf T}R_N(q)}
{The left curves follow one pair's intermediate eigenvalues as q changes. On the right, the complete overlap at one q value records how the left and right Casimir atlases mix. The packet dimension remains seven throughout this example. The changed colors describe coordinate geometry on a fixed-size object, rather than a change in the number of channels. Spectral agreement and overlap orthogonality are checked separately.}
{For $N=6$, shapes $(0.7,1.1,0.9)$ and $q=0.73$, build the full degree-six lowering map, compute its seven-dimensional kernel, and restrict both pair Casimirs. Their ordered eigenvectors define the overlap. The code compares the left eigenvalues with the displayed q-number formula and checks $\mathcal R^{\mathsf T}\mathcal R=I$. This is a finite matrix construction. The exact scalar q-Racah parameter and weight map is kept as a separate comparison obligation.}
{The matrix module is usable before an explicit scalar normalization formula is supplied, provided its evidence is described at that level.}
{VIII, Chapter 5, Theorem 5.1 and \emph{Two intermediate q-Casimir atlases}. Classical pair-spectrum/q-Racah inputs; the plotted matrix is newly computed.}

\lesson{Paper VIII / 05}
{An exact q-coordinate reaches a swap}
{viii_05_q_beam}
{\mathcal R_q=\begin{pmatrix}c_q&s_q\\s_q&-c_q\end{pmatrix},\qquad
c_q=\frac{q}{1+q^2},\quad s_q=\frac{\sqrt{1+q^2+q^4}}{1+q^2}}
{Read the amplitude curves from small q toward one. The diagonal amplitude rises from zero to one half, while the transferred amplitude approaches $\sqrt3/2$. The right panel displays the complete signed transform at the midpoint q value. Its limit at zero is a routing swap. This is an exact limit for this normalized two-channel model. The unnormalized q-number operators themselves do not have a uniform finite limit at zero.}
{At $q=0.5$, $c_q=0.4$ and $s_q=\sqrt{0.84}\approx0.916515$. Their squares sum to one, and the lower-right sign makes the two columns orthogonal. The numerical pair-Casimir construction matches this matrix at four q values. At $q=1$, it recovers the classical half-shape beam splitter. At a very small positive q, the closed expression approaches the swap matrix directly, avoiding an ill-conditioned large-degree operator construction.}
{State the normalization and tested model before promoting a visually sparse q endpoint into a general routing theorem.}
{VIII, Chapter 6, Theorem 6.1, \emph{Closed two-channel q-Racah transform}, half-shapes and $N=1$. Routing concerns this closed formula only.}

\lesson{Paper VIII / 06}
{Two-channel mixing has a monotone gap}
{viii_06_entropy_gap}
{p(q)=\frac{q^2}{(1+q^2)^2},\qquad \lambda_2=2p-1,\qquad
\operatorname{gap}_{\rm abs}=1-|\lambda_2|=2p}
{The left curve is the binary entropy of one measured row. The right panel separates the absolute mixing gap from the magnitude of the other eigenvalue. Small q is close to a periodic swap, whose second eigenvalue has magnitude one. Moving toward q=1 reduces that magnitude and increases the absolute gap. The word absolute matters: a negative second eigenvalue still produces an oscillating mode, and its distance from one alone is not the mixing rate used here.}
{At $q=0.5$, the measured matrix has rows $(0.16,0.84)$ and $(0.84,0.16)$. Its second eigenvalue is $-0.68$, its absolute gap is $0.32$, and its row entropy is approximately $0.43967$ nats. The smaller row probability increases only to $1/4$ over the declared interval. Binary entropy is increasing there, which proves the model's q-monotonicity without relying on the plotted grid.}
{The exact one-excitation model has a monotone deformation law; higher multiplicity requires its own analysis.}
{VIII, Chapter 7, Theorem 7.2, \emph{Monotone one-excitation mixing}. Equal half-shapes, $N=1$, $0<q\leq1$, absolute spectral gap.}

\lesson{Paper VIII / 07}
{Higher-level row entropy can decrease}
{viii_07_nonmonotone}
{H_1(q)=-\sum_m|\mathcal R_3(q)_{1m}|^2\log|\mathcal R_3(q)_{1m}|^2,\qquad
H_1(q_1)-H_1(q_0)\approx-0.007784}
{The left panel follows the four probabilities of the second measured row, using zero-based row index one. The right panel computes their entropy and highlights two nearby q values. The curve decreases between them. Every individual kernel remains bistochastic, so each still obeys the input-to-output entropy inequality. The decrease compares two different kernels' row laws. It therefore refutes universal entropy monotonicity in q without contradicting the fixed-channel theorem.}
{For $N=3$ and three half-shapes, take $q_0=307/395\approx0.777215$ and $q_1=1557/1975\approx0.788354$. The entropies are approximately $1.244882$ and $1.237098$ nats. Their difference is negative by about $0.007784$. The guide recomputes both matrices in double precision from the q-coproduct. A separate full-space implementation records a decreasing interval from $q=0.78$ to $0.80$. These are bounded numerical witnesses, not interval proof certificates or a classification of all deformation behavior.}
{A negative control defines the theorem boundary as precisely as a successful example defines its mechanism.}
{VIII, Chapter 7, Boundary 7.3, \emph{No universal entropy monotonicity in q}. New endpoint witness in vii\_viii\_audit.json and separately saved independent full-space evidence.}

\lesson{Paper VIII / 08}
{Differentiate a complete coordinate transform}
{viii_08_tangent}
{A(\theta)=\mathcal R(\theta)^{\mathsf T}\partial_\theta\mathcal R(\theta),\qquad
A^{\mathsf T}=-A,\quad \partial_\theta\mathcal R=\mathcal R A}
{The left matrix is the skew tangent generator for the seven-channel transform at $N=6$. Opposite entries have opposite signs, and diagonal entries vanish. The right curves measure squared speed for two excitation levels in the same $\theta=-\log q$ coordinate. This quantity concerns sensitivity of the complete atlas overlap, rather than entropy of its measured rows. Its size depends on parameterization, while constant orthogonal changes of atlas preserve the Frobenius coefficient.}
{At $\theta=0.35$, compute the overlaps at $\theta\pm2\times10^{-5}$, align their real sign conventions, and form a central difference. Before projecting to a skew matrix, the code checks the raw generator's symmetric residual. It also checks the tangent equation afterward. This avoids treating an explicitly antisymmetrized output as independent evidence of skewness. Analytically, differentiating $\mathcal R^{\mathsf T}\mathcal R=I$ proves the exact skew law in a smooth gauge.}
{A deformation tangent is meaningful only after the atlas conventions are made continuous along the chosen coordinate.}
{VIII, Chapter 8, Theorem 8.1, \emph{Skew deformation generator}. New finite-difference diagnostics; derivative step and raw residual recorded in the audit.}

\lesson{Paper VIII / 09}
{Approach the classical atlas carefully}
{viii_09_classical_limit}
{[x]_q=x+O((\log q)^2),\qquad \mathcal R_N(q)\longrightarrow\mathcal R_N(1)
\quad\text{in a matched sign gauge}}
{Both curves on the left tend toward zero error as q approaches one. One compares complete recoupling overlaps, and the other compares normalized two-body seed vectors. Their rates differ here. The right panel retains the signed overlap correction at $q=0.99$. The q-number's even hyperbolic expansion does not by itself establish the same rate for every seed or coordinate choice, because the coproduct also carries directional q-powers.}
{For $N=6$ and shapes $(0.7,1.1,0.9)$, compute the q=1 overlap as the reference. Align row and column signs before taking entry differences at nearby q. The seed comparison uses the two-body degree-five vector for shapes $(0.7,1.1)$. Both finite comparisons converge over the displayed range. They verify these numerical examples and conventions. The continuity argument is supplied by the finite matrix construction and its nondegenerate spectral data.}
{Match atlases before measuring a limit, and distinguish an operator coefficient expansion from an entire reconstructed object's convergence rate.}
{VIII, Chapter 9, \emph{Classical and routing limits}, source lines 751--763. New q-grid comparisons at fixed positive shapes; no universal rate asserted.}

\lesson{Paper VIII / 10}
{q-trees need more than loop closure}
{viii_10_q_trees}
{U_{r,r+1}(q)=B_r(q)^{\mathsf T}B_{r+1}(q),\qquad
U_{01}U_{12}U_{23}U_{34}U_{40}=I\quad\text{(indices modulo 5)}}
{The five vertices still represent the same ordered planar bracketings after deformation. Each now uses q-Casimirs on contiguous blocks. The right panel plots three separate diagnostics: common-space loop closure, annihilation by total lowering, and joint diagonalization. Their residuals remain small in the displayed finite tests. They answer different questions. Combining them gives more information than calling the closing overlap loop alone a certificate of every label, phase, and projector convention.}
{At four factors and $N=3$, the multiplicity dimension is ten. For q values $0.65$, $0.75$, $0.85$, and one, the code constructs every joint atlas and composes its five overlaps. It separately tests commuting intermediate operators and the displayed kernel and diagonalization residuals. Consistent coordinate changes would still leave the loop equal to identity. A future stronger local-associator check would compare independently constructed q-Racah edge blocks with these global overlaps before composing them.}
{Keep structural operator checks alongside loop checks when the same data construction could make closure automatic.}
{VIII, Chapter 10, Theorem 10.1, \emph{Finite q-pentagon certificate}, with corrected interpretation. Four-factor inputs and residuals recorded separately.}

\lesson{Paper VIII / 11}
{A scalar table is a separate comparison}
{viii_11_scalar_bridge}
{\begin{gathered}
R_n(\mu_{\gamma,\delta}(y);\alpha,\beta,\gamma,\delta\mid q)
=R_y(\mu_{\alpha,\beta}(n);\gamma,\delta,\alpha,\beta\mid q)\\
\mu_{u,v}(t)=q^{-t}+uvq^{t+1}
\end{gathered}}
{The left panel is generated from the terminating scalar q-hypergeometric expression, with independent scalar parameters. Exchanging degree and coordinate together with the parameter pairs gives the recorded duality agreement. This table has not been identified entry by entry with the shape-Casimir overlap. The right panel gives a finite-support counterexample: the two lattice points coincide, and two evaluated polynomial rows coincide. A denominator truncation condition alone therefore cannot guarantee a positive orthogonal family.}
{Set $N=1$, $q=1/2$, $\alpha=\beta=1/2$, $\gamma=4$, and $\delta=1$. Although $\gamma q=q^{-N}$, both lattice values are three. The finite formula gives $R_0=1$ and $R_1(\mu)=6-(5/3)\mu$, so both row vectors are $(1,1)$. Their inner product under positive weights is $w_0+w_1>0$. This exact rational witness separates termination, nondegeneracy, and positivity. The complete shape-to-scalar map still needs base, parameters, weights, norms, and phases.}
{Scalar identities can be verified independently while a matrix-to-polynomial normalization comparison remains explicit unfinished work.}
{VIII, Chapter 13, Theorem 13.1; Chapter 19, first open direction. Scalar sums require $0<q<1$ and nonzero finite denominator factors; $q=1$ is a separate limit. The collapsed-lattice witness uses exact rational arithmetic.}

\lesson{Paper VIII / 12}
{Check a shifted Askey--Wilson operator}
{viii_12_askey_wilson}
{\mathcal L\mathcal A_n=(q^{-n}+abcdq^{n-1})\mathcal A_n,\qquad
\mathcal L=D_q+(1+abcd/q)I}
{The left curves are normalized terminating Askey--Wilson polynomials on the continuous unit-circle coordinate. The right panel substitutes six degrees into the q-difference equation at thirteen nonsingular points. The residual scale shows numerical agreement. The eigenvalue uses the shifted operator: $D_q$ is the two-difference part, while the scalar identity term has been added. Omitting that shift would change the eigenvalue convention and make the comparison inconsistent.}
{Take $q=0.7$ and $(a,b,c,d)=(0.2,0.3,-0.25,0.4)$. At degree two, the shifted eigenvalue is $q^{-2}+abcdq\approx2.036616$. Evaluate the finite polynomial at $z$, $qz$, and $z/q$, then apply the declared operator coefficients away from their poles. The program compares this result with the eigenvalue times the unshifted polynomial value. This checks a classical scalar identity. It does not construct a continuous-series unitary packet transform or its spectral measure.}
{A continuous variable in a scalar polynomial is a bridge to further analysis, not an already completed continuous-series transform.}
{VIII, Chapter 14, Theorem 14.1, \emph{Askey-Wilson eigenfunction law}. Scalar sums require $0<q<1$, nonzero finite denominator factors, and nonsingular operator points; $q=1$ is a separate limit. Classical identity, newly evaluated on the declared grid.}

\startvolume{Synthesis}{Connections across the eight papers}

\lesson{SYNTHESIS / 01}{Three reconstruction contracts}{synthesis_01}
{\mathcal S_M\mathcal A_M=I,\qquad \mathcal V_g^*\mathcal V_g=I,\qquad R^{\mathsf T}R=I_{N+1}}
{The series uses three different coordinate spaces. Coefficient packets partition an index set. A normalized coherent window analyzes a function over a continuous phase plane. Casimir atlases analyze a finite multiplicity space. Each has an exact synthesis map, although its coefficients, measure, and admissible range differ. The middle panel plots coherent coefficient intensity $|A|^2$ for a log-Gamma window with shape $3/2$.}
{Split $(1,2,3,4,5,6)$ into $(1,3,5)$ and $(2,4,6)$, then interleave to recover all six entries. At $N=1$, the displayed recoupling matrix sends $(1,0)$ to $(1/2,\sqrt3/2)$; its transpose recovers the input. In the continuous case, the unit-window Moyal identity integrates squared coefficients with measure $dh\,d\eta/(2\pi)$ to recover the input norm.}
{Identify the coordinate space and its synthesis rule before transporting an identity.}
{I: exact polyphase reconstruction; VI: continuous tight frame; VII--VIII: finite reconstruction. Analytic frame identities come from the corrected shape core; the diagram does not identify these three transform ranges.}

\lesson{SYNTHESIS / 02}{Arithmetic parameters become a shape}{synthesis_02}
{\alpha=\frac{2a-1}{\beta}=2\kappa,\qquad K_+|n\rangle=\sqrt{(n+1)(n+\alpha)}\,|n+1\rangle}
{The arithmetic scale exponent fixes a positive Gamma shape. The same shape enters the log-Gamma window and the ordinary positive ladder after setting $\kappa=\alpha/2$. This is a precise parameter dictionary across IV--VII. The curves show Gamma densities in the variable $u=e^y$, so their normalization uses $du$, whereas the window density uses $dy$.}
{Take $a=2$ and $\beta=2$. Then $\alpha=3/2$, $\kappa=3/4$, and the first two raising coefficients are $\sqrt{3/2}$ and $\sqrt5$. Taking $a=3/4$, $\beta=1/2$ instead gives $\alpha=1$: the positive kernel is admissible because $a>1/2$. The original absolutely convergent signed series requires $a>1$, so that second choice belongs to the positive branch.}
{One shape links the formulas; each realization retains its own domain and measure.}
{IV--V: $\alpha=(2a-1)/\beta$; V--VI: normalized log-Gamma windows and ladder; VII: $\alpha=2\kappa$. The bridge concerns the finite-span Lie-algebra formulas, not a new global-group integration theorem.}

\lesson{SYNTHESIS / 03}{A diagonal lift removes arithmetic signs}{synthesis_03}
{E=\sum_n|v_n|^2,\qquad |v_1+v_2|^2=4\cos^2(\theta/2)\quad(v_1=1,\ v_2=e^{i\theta})}
{Time averaging the arithmetic amplitude retains its diagonal packet energy. Its squared M\"obius factor is positive, but the original signed sum uses information that this diagonal contraction removes. The illustration varies a relative phase while holding both component magnitudes fixed. The energy stays at two, while the coherent sum ranges from four to zero. A full Gram kernel may retain other phase relations; the loss here is the specified diagonal operation.}
{For $(v_1,v_2)=(1,1)$, the energy is two and the squared sum is four. For $(1,-1)$, the energy remains two and the squared sum is zero. Therefore identical positive diagonal data cannot determine the signed amplitude. This exact two-component witness explains why the positive Gamma, frame, and semigroup results do not supply the principal M\"obius cancellation estimate.}
{Record which phases an averaging or measurement operation discards.}
{IV: exact time-average identity; I--VI: signed-coherence boundary; VII--VIII: measured versus coherent recoupling. This finite phase illustration is not an arithmetic counterexample or an estimate for M\"obius sums.}

\lesson{SYNTHESIS / 04}{A binary refinement becomes a tree}{synthesis_04}
{\left(C_{n,\mathbf r}^{\mathbf s}\right)^2=\frac{n!\prod_j(s_j)_{r_j}}{(\sum_js_j)_n\prod_jr_j!},\qquad \sum_jr_j=n}
{Shape addition produces normalized two-child coefficients. Multiplying them down a refinement tree cancels every intermediate factorial and Pochhammer factor. The remaining expression depends only on leaf shapes and degrees. Nodes show both numbers: a shape is not an excitation degree. This coefficient cancellation is the shape-level analogue of composing coordinate changes across recoupling trees; the two constructions use different coefficient systems.}
{Choose shapes $(1,2,3)$, total degree three, and leaf degrees $(1,1,1)$. The left route has squared factors $9/28$ and $1/3$; the right route has $15/56$ and $2/5$. Both products are $3/28$, hence both positive amplitude coefficients equal $\sqrt{3/28}$. Enumerating all ten degree-three leaf compositions gives squared coefficients summing exactly to one. This verifies the displayed example independently of a tree drawing.}
{Cancel intermediate normalization factors before interpreting a refinement as tree independent.}
{VI: two-child refinement and Dirichlet-multinomial extension; VII--VIII: common-space tree atlases. Exact rational checks cover the stated three-child example; general cancellation is the generating-function identity in the corrected shape core.}

\lesson{SYNTHESIS / 05}{Positivity needs a specified measure}{synthesis_05}
{G_{ij}=\int x^{i+j}\,d\nu(x),\qquad \gamma q=q^{-N}\ \text{alone does not imply positive orthogonality}}
{A positive measure makes its moment matrix a Gram matrix: every quadratic form is an integral of a squared polynomial. This mechanism underlies III's positive spectral branch. A terminating scalar q-Racah expression is a separate construction. Nonzero denominators make its values defined, but do not ensure a distinct lattice, positive weights, or nonzero orthogonal norms. The right panel shows two labels landing on the same lattice point.}
{For the illustrative measure $\nu=\tfrac12\delta_{1/4}+\tfrac12\delta_{3/4}$, the first moment matrix has entries $1,1/2,5/16$ and determinant $1/16$. In VIII's exact degenerate example, take $N=1$, $q=1/2$, $\alpha=\beta=1/2$, $\gamma=4$, $\delta=1$. Then $\mu(0)=\mu(1)=3$ and $R_0=R_1=(1,1)$. Their inner product equals one for any normalized positive weights.}
{Check the lattice, weight, norms, parameter dictionary, and phases separately.}
{III: positive moment measure; VIII: scalar formulas, corrected by C11--C12. The left measure is a two-atom illustration, not the infinite arithmetic measure. The scalar counterexample is evaluated exactly with nonzero finite denominators.}

\lesson{SYNTHESIS / 06}{Match the evidence to the question}{synthesis_06}
{c_{1/2}^{\,2}=\frac4{25},\qquad s_{1/2}^{\,2}=\frac{21}{25},\qquad c_{1/2}^{\,2}+s_{1/2}^{\,2}=1}
{Three source rows ask different questions. The coefficient example compares finite reconstruction routes. The q-matrix example checks exact normalization and floating orthogonality at specified inputs. The coherence conjecture requires one constant for each positive exponent, uniformly over depth. A finite calculation preserves that distinction: it answers a bounded question without weakening the quantifiers of another claim.}
{For I.theorem.002, evaluate the six-term polynomial at $z=1/3$ directly and through two residue packets: both give $181/81$. For VIII.theorem.011 at $q=1/2$, the squared coefficients sum to one exactly; five floating-point inputs also have orthogonality residual below $10^{-14}$. The ledger marks these rows partial and queued, respectively. V.conjecture.021 remains conjectural. No Lean kernel build is performed by these calculations.}
{Attach each result to its input, arithmetic domain, source row, and verification method.}
{Preserved claim index: I.theorem.002, VIII.theorem.011, V.conjecture.021, with full IDs in the computation JSON. Status labels come from the retrieved ledger snapshot; they are not a claim that all later formalization has been searched exhaustively.}

\startvolume{Sequence atlas}{OEIS families and exact transforms}
\lesson{SEQUENCES / 01}{Two seeds, one Fibonacci recurrence}{S01_fibonacci_lucas}
{F_{n+2}=F_{n+1}+F_n,\quad L_{n+2}=L_{n+1}+L_n,\quad L_n^2-5F_n^2=4(-1)^n}
{The two families share the same recurrence and have different seeds: $(F_0,F_1)=(0,1)$ and $(L_0,L_1)=(2,1)$. The right panel uses the same index in both families. Its alternating defect stays exactly four in magnitude.}
{At $n=7$, $F_7=13$ and $L_7=29$. Thus $29^2-5\cdot13^2=841-845=-4$. Comparing terms at different indices would destroy this identity.}
{A recurrence operator and its initial state are separate parts of the sequence specification.}
{OEIS A000045 and A000032, indexed from zero. Exact recurrence and common-index identity; the plotted prefix ends at $n=12$.}

\lesson{SEQUENCES / 02}{Longer recurrence memory}{S02_tribonacci_tetranacci}
{T_n=T_{n-1}+T_{n-2}+T_{n-3},\qquad Q_n=\sum_{j=1}^{4}Q_{n-j}}
{The seed table fixes the OEIS convention. Tribonacci starts $0,0,1$; tetranacci starts $0,0,0,1$. The curves use $\log_{10}(1+a_n)$ so that zeros and rapidly growing terms fit together. The transform changes the plotted scale, not the integer sequence.}
{Tribonacci gives $T_5=2+1+1=4$. Tetranacci gives $Q_7=4+2+1+1=8$. The recurrence first applies at $n=3$ for $T$ and $n=4$ for $Q$.}
{These finite-memory recurrences have three and four feedback taps. The seed length must match the memory.}
{OEIS A000073 and A000078, both indexed from zero. Exact integer terms through $n=16$; logarithmic ordinate explicitly labeled.}

\lesson{SEQUENCES / 03}{Jacobsthal separates two modes}{S03_jacobsthal}
{J_0=0,\quad J_1=1,\quad J_{n+2}=J_{n+1}+2J_n=\frac{2^{n+2}-(-1)^{n+2}}{3}}
{The rising bars are Jacobsthal numbers. The right panel removes their growing mode: $2^n-3J_n=(-1)^n$. What remains is an alternating mode of constant magnitude. Both modes belong to the same second-order recurrence.}
{At $n=6$, $J_6=21$ and $2^6-3\cdot21=1$. At $n=7$, $J_7=43$ and $2^7-3\cdot43=-1$.}
{A recurrence can combine growth and oscillation. Splitting the modes explains the parity correction.}
{OEIS A001045, indexed from zero. Direct closed form and exact recurrence agree for $0\le n\le18$.}

\lesson{SEQUENCES / 04}{A Pell pair shares one unit}{S04_sqrt3_pell}
{x_n+y_n\sqrt3=(2+\sqrt3)^n,\qquad x_n^2-3y_n^2=1}
{The table keeps both coordinates at the same index. Multiplication by $2+\sqrt3$ advances the pair using the displayed integer update. Its determinant is one, and it preserves the quadratic norm. The zero-index pair is $(1,0)$.}
{Starting from $(x_2,y_2)=(7,4)$ gives $(x_3,y_3)=(2\cdot7+3\cdot4,7+2\cdot4)=(26,15)$. The norm is $26^2-3\cdot15^2=1$.}
{A shared recurrence state can retain an exact invariant while both coordinates grow.}
{OEIS A001075 and A001353, indexed from zero. Integer matrix iteration and the norm are checked through $n=18$.}

\lesson{SEQUENCES / 05}{Normalize central binomial coefficients}{S05_binomial_catalan}
{B_n=\binom{2n}{n},\qquad C_n=\frac{B_n}{n+1}}
{Each table column shows one central binomial coefficient and its Catalan normalization. The curves extend the comparison through index fourteen on a base-ten logarithmic axis. Division by $n+1$ is exact here and gives an integer at every index.}
{For $n=4$, $B_4=\binom84=70$ and $C_4=70/5=14$. For $n=5$, $252/6=42$. The divisor changes with the index.}
{Normalization is an operator on coefficients. It need not be a fixed convolution or a constant tap weight.}
{OEIS A000984 and A000108, indexed from zero. Binomial formulas and integrality checked through $n=18$.}

\lesson{SEQUENCES / 06}{Catalan convolution rebuilds the next term}{S06_catalan_convolution}
{C_{n+1}=\sum_{k=0}^{n}C_kC_{n-k},\qquad C(z)=1+zC(z)^2}
{The five bars are the five products contributing to $C_5$. The input indices add to four, and their sum supplies the next coefficient. Here $C(z)=\sum_{n\ge0}C_nz^n$ is an ordinary generating function.}
{The contributions are $1\cdot14$, $1\cdot5$, $2\cdot2$, $5\cdot1$, and $14\cdot1$. Their sum is $14+5+4+5+14=42=C_5$.}
{A nonlinear generating-function equation becomes an exact convolution identity among coefficients.}
{OEIS A000108. All convolution identities with $0\le n\le17$ are checked against the binomial definition.}

\lesson{SEQUENCES / 07}{Bell numbers are Stirling row sums}{S07_bell_stirling}
{B_n=\sum_{k=0}^{n}S(n,k),\qquad \sum_{n\ge0}B_n\frac{t^n}{n!}=\exp(e^t-1)}
{Each matrix row contains Stirling numbers of the second kind, counting partitions into $k$ nonempty blocks. Summing across the row gives the Bell number at its right. The generating function is exponential, so its coefficient is $B_n/n!$.}
{The fifth row contains $0,1,15,25,10,1$. Its sum is $52=B_5$. Equivalently, apply the Stirling transform to the constant input sequence $1,1,1,\ldots$.}
{A triangular transform uses all relevant input positions. It is not automatically a fixed-length tap filter.}
{OEIS A000110. Stirling recurrence and an independent binomial Bell recurrence agree through $n=18$; the figure shows rows zero through seven.}

\lesson{SEQUENCES / 08}{Factorial weights order the blocks}{S08_fubini_factorial}
{A_n=\sum_{k=0}^{n}k!S(n,k),\qquad \sum_{n\ge0}A_n\frac{t^n}{n!}=\frac{1}{2-e^t}}
{The same Stirling row now receives factorial weights. Ordering $k$ blocks gives $k!$ arrangements. The upper table shows the factors at $n=5$; the bars show their products. The sum is the ordered Bell, or Fubini, number.}
{At $n=5$, the weighted contributions are $1,30,150,240,120$, totaling $541$. The unweighted row sums to only $52$. The weights encode a different combinatorial operation.}
{Changing the transform's input from one to $k!$ changes the output family while retaining the same triangular kernel.}
{OEIS A000670 and A000142, indexed from zero. Integer Stirling sums and an independent ordered-partition recurrence agree through $n=18$.}

\lesson{SEQUENCES / 09}{Derangements show signed cancellation}{S09_derangements}
{D_n=n!\sum_{k=0}^{n}\frac{(-1)^k}{k!},\qquad D_n=(n-1)(D_{n-1}+D_{n-2})\quad(n\ge2)}
{Positive and negative bars are the inclusion-exclusion contributions at $n=6$. Their large magnitudes cancel to a smaller integer. The dashed line records the final sum rather than an individual contribution. Here $D_0=1$ and $D_1=0$.}
{The sum is $720-720+360-120+30-6+1=265$. The recurrence independently gives $D_6=5(44+9)=265$.}
{Signed transforms must retain cancellation. Magnitudes alone do not determine the result.}
{OEIS A000166. Factorial sum, variable-coefficient recurrence, and brute-force permutation counts through $n=8$ agree; exact formulas checked through $n=18$.}

\lesson{SEQUENCES / 10}{Secant and tangent interleave}{S10_euler_updown}
{\sum_{n\ge0}E_n\frac{t^n}{n!}=\sec t+\tan t,\quad E_{2m}=s_m,\quad E_{2m+1}=\tau_m}
{The coefficient table separates even and odd positions of the Euler up/down family. Even positions are positive secant numbers. Odd positions are tangent numbers. The factorial row makes clear that these are exponential-generating-function numbers, not ordinary Taylor coefficients.}
{At degree four, $E_4=5$ contributes $5t^4/4!$ to secant. At degree five, $E_5=16$ contributes $16t^5/5!$ to tangent. The first tangent entry corresponds to degree one.}
{Parity packets can identify familiar functions, provided the factorial normalization and sequence offset are kept.}
{OEIS A000111, A000364, and A000182. The tangent entry has OEIS initial index one; $\tau_m$ denotes its degree-$2m+1$ value. Exact series checked through degree eighteen.}

\lesson{SEQUENCES / 11}{Bernoulli data need both coordinates}{S11_bernoulli_pair}
{\frac{t}{e^t-1}=\sum_{n\ge0}B_n\frac{t^n}{n!},\qquad B_n=\frac{p_n}{q_n},\quad q_n>0}
{The top panel shows signed numerators on a symmetric logarithmic scale. The lower panel shows positive denominators on a logarithmic scale. Fractions are reduced. A zero Bernoulli number is stored as $0/1$, and this guide uses $B_1=-1/2$.}
{At $n=10$, $p_{10}=5$ and $q_{10}=66$, so $B_{10}=5/66$. Its Taylor coefficient is $(5/66)/10!$. At $n=12$, $B_{12}=-691/2730$.}
{An OEIS integer numerator sequence does not by itself specify the rational family or its generating-function coefficient.}
{OEIS A027641 and A027642, indexed from zero. Exact rational recurrence and multiplication by $(e^t-1)/t$ checked through degree eighteen.}

\lesson{SEQUENCES / 12}{Partitions multiply reciprocal factors}{S12_partition_product}
{P(z)=\prod_{j\ge1}\frac{1}{1-z^j}=\sum_{n\ge0}p(n)z^n}
{A factor indexed by $j$ chooses how many parts of size $j$ occur. The left table lists the five partitions of four. The right bars show the first nineteen partition numbers. This is an ordinary generating function with integer coefficients.}
{The partitions of four are $4$, $3+1$, $2+2$, $2+1+1$, and $1+1+1+1$. Thus $p(4)=5$. To find coefficients through degree eighteen, only factors with $j\le18$ are needed.}
{Reciprocal factors turn independent multiplicities into a coefficient-counting product. The finite prefix can be reconstructed exactly.}
{OEIS A000041. Product coefficients computed by integer dynamic programming, independently checked by partition enumeration through $n=10$.}

\lesson{SEQUENCES / 13}{Apéry uses taps that depend on the index}{S13_apery_recurrence}
{(n+1)^3A_{n+1}=(2n+1)(17n^2+17n+5)A_n-n^3A_{n-1}}
{The prefix annotation begins the Apéry sequence defined by $A_n=\sum_{k=0}^{n}\binom nk^2\binom{n+k}{k}^2$. The curves show its growth through index fourteen on a base-ten logarithmic scale. The lower diagram evaluates the recurrence at $n=2$.}
{Using $A_1=5$ and $A_2=73$ gives $27A_3=535\cdot73-8\cdot5=39015$. Therefore $A_3=1445$. The recurrence applies for $n\ge1$, with $A_0=1$.}
{This is a three-term recurrence with polynomial coefficients. Its weights change with $n$, so it is not a constant-tap filter.}
{OEIS A005259. Exact binomial sums and the displayed polynomial-coefficient recurrence agree for $1\le n\le17$.}

\lesson{SEQUENCES / 14}{Divisor convolution changes the address rule}{S14_divisor_convolution}
{\sum_{d\mid n}\mu(d)=\mathbf1_{n=1},\qquad \sum_{d\mid n}\varphi(d)=n}
{The figure places the six divisors of twelve in matching columns. Möbius contributions sum to zero; totient contributions sum to twelve. The address rule is divisibility. This differs from the additive index splitting used by ordinary power-series convolution.}
{For $n=12$, $\mu(1),\mu(2),\mu(3),\mu(4),\mu(6),\mu(12)$ are $1,-1,-1,0,1,0$. The totient values are $1,1,2,2,2,4$. Their sums are zero and twelve.}
{Choose the convolution law before identifying a transform. Divisor packets and residue packets organize different relations.}
{OEIS A008683 and A000010, both indexed from one. Prime-factor definitions and both divisor sums checked for $1\le n\le60$.}

\lesson{SEQUENCES / 15}{One binary statistic creates three families}{S15_binary_families}
{w(n)=\operatorname{popcount}(n),\quad t(n)=w(n)\bmod2,\quad g(n)=2^{w(n)}}
{Each column has a four-bit address, its binary weight, its Thue--Morse parity, and its Gould value. Reading a final zero bit preserves weight; reading a final one adds one. Parity toggles and the Gould value doubles in the latter case.}
{The address $13=1101_2$ has weight three, so $t(13)=1$ and $g(13)=8$. Address $12=1100_2$ has weight two, parity zero, and Gould value four.}
{The same state statistic can feed several exact outputs. Binary refinement gives simple even and odd channel recurrences.}
{OEIS A000120, A010060, and A001316, indexed from zero. Binary formulas and even/odd recurrences checked for $0\le n\le63$; figure shows zero through fifteen.}

\lesson{SEQUENCES / 16}{Sums and differences form a polynomial ladder}{S16_polynomial_ladder}
{T_n=\sum_{j=1}^{n}j,\quad H_n=\sum_{j=1}^{n}T_j,\quad \Delta T_n=n,\quad \Delta H_n=T_n}
{The upper table follows positive integers into triangular and tetrahedral numbers. Squares and cubes occupy the lower rows for comparison. The arrows reverse prefix sums by backward differences, $\Delta f_n=f_n-f_{n-1}$. All plotted values use their displayed common index.}
{At $n=4$, $T_4=10$ and $H_4=20$. The differences are $10-6=4$ and $20-10=10$. Separately, $4^2-3^2=7$ and $4^3-3^3=37$.}
{Prefix summation raises polynomial degree; differencing lowers it. These are exact operators with boundary values at zero; the displayed difference identities apply for $n\ge1$.}
{OEIS A000027, A000217, A000292, A000290, and A000578. Positive integers start at one; the other four families include index zero. Identities checked through $n=18$.}

\startvolume{Polyphase atlas}{Channels, refinement, and transfer}

\lesson{POLYPHASE / 01}{Split Fibonacci into three actual streams}{P01_fibonacci_streams}
{F_{3,r}(u)=\sum_{m\ge0}F_{3m+r}u^m,\qquad r=0,1,2}
{The large numbers are Fibonacci coefficients, not exponent labels. Each column advances one step inside a channel, which advances three steps in the original sequence. The smaller label underneath records the original index. The three colors together account for every index from zero to seventeen.}
{The first channel begins $0,2,8,34,144,610$. The second begins $1,3,13,55,233,987$, and the third begins $1,5,21,89,377,1597$. Thus the coefficient of $u^4$ in channel two is $F_{14}=377$.}
{Declare the starting index before splitting a sequence. Here $F_0=0$, $F_1=1$, and channel indices start at $m=0$.}
{Fibonacci: OEIS A000045. Exact integer coefficients; the displayed prefix ends at index seventeen. The splitting definition is valid for the full formal series.}

\lesson{POLYPHASE / 02}{Weave the three streams back together}{P02_interleave}
{\sum_{n\ge0}F_nz^n=F_{3,0}(z^3)+zF_{3,1}(z^3)+z^2F_{3,2}(z^3)}
{Read the top rows horizontally to see each internal stream. Read the lower row in increasing original index to see the reconstruction. Its colors repeat zero, one, two. Synthesis restores both the coefficient and the place where it belongs.}
{The channel terms $2u$, $3u$, and $5u$ become $2z^3$, $3z^4$, and $5z^5$. Interleaving the four displayed terms from each channel produces precisely the first twelve Fibonacci coefficients, ending with $F_{11}=89$.}
{Splitting followed by synthesis is a change of coordinates. It does not smooth, average, or estimate the sequence.}
{OEIS A000045; Paper I, exact polyphase reconstruction. Equality follows coefficient by coefficient from the unique division $n=3m+r$. The numerical example is finite.}

\lesson{POLYPHASE / 03}{Two Fibonacci channels share one denominator}{P03_fibonacci_block}
{E(u)=\frac{u}{1-3u+u^2},\quad O(u)=\frac{1-u}{1-3u+u^2},\quad B=\begin{pmatrix}1&1\\1&2\end{pmatrix}}
{Each box is the adjacent pair $(F_{2m},F_{2m+1})$. The same block matrix advances every box by two original indices. Both channel series therefore obey the same second-order recurrence, although they have different starting coefficients.}
{The initial pair is $(0,1)$. Applying $B$ gives $(1,2)$, then $(3,5)$, then $(8,13)$. Consequently $E=0+u+3u^2+8u^3+\cdots$ and $O=1+2u+5u^2+13u^3+\cdots$.}
{A shared recurrence does not imply identical channels. Initial values and numerator data retain the residue information.}
{OEIS A000045. Here $B=T^2$ for $T=\left(\begin{smallmatrix}0&1\\1&1\end{smallmatrix}\right)$. The denominator is a formal-series identity, equivalently a rational identity near zero.}

\lesson{POLYPHASE / 04}{Every Fibonacci stride has a jump recurrence}{P04_jump_recurrences}
{g_m=F_{Mm+r},\qquad g_{m+2}=L_Mg_{m+1}-(-1)^M g_m\quad(M\ge1)}
{Each table row chooses a stride $M$. The trace of $T^M$ is the Lucas number $L_M$, while its determinant is $(-1)^M$. The last three columns use the same starting offset $r=1$ and show the two initial values and the next value.}
{For $M=3$, start with $g_0=F_1=1$ and $g_1=F_4=3$. Since $L_3=4$ and the determinant is $-1$, the next value is $4\cdot3+1=13=F_7$. For $M=4$, the sign changes: $7\cdot5-1=34$.}
{The stride determines the recurrence. The starting offset determines $g_0=F_r$ and $g_1=F_{M+r}$. A residue channel uses $0\le r<M$.}
{OEIS A000045 and A000032, with $L_0=2$, $L_1=1$. Here $M$ is a positive integer and $m,r\ge0$. The formula follows from Cayley--Hamilton for $T^M$; the table displays strides one through six at offset one.}

\lesson{POLYPHASE / 05}{Complex phases select one residue}{P05_root_selector}
{z^rF_{M,r}(z^M)=\frac1M\sum_{j=0}^{M-1}\zeta^{-rj}F(\zeta^jz),\qquad \zeta=e^{2\pi i/M}}
{The triangle places the three cube roots of unity in the complex plane. The selection grid shows the resulting coefficient mask: an exponent survives in exactly one row. Individual evaluations can be complex even when the original coefficients and the reconstructed residue series are real.}
{At $M=3$, the coefficient $F_7=13$ survives the selector for $r=1$ because $7-1$ is divisible by three. The sum of its three phase factors is three. For $r=0$ or $r=2$, those phase factors sum to zero.}
{This selector uses a characteristic-zero field containing the required roots of unity, here the complex numbers. Address splitting itself needs no such extension.}
{Paper I, roots-of-unity filters; $M$ is a positive integer and $0\le r<M$. OEIS A000045 supplies the example. Exact cancellation is checked in $\mathbb Q[\zeta]/(\zeta^2+\zeta+1)$, not inferred from floating-point cancellation.}

\lesson{POLYPHASE / 06}{Refine even and odd into six coefficient streams}{P06_refined_coefficients}
{F_{2,r}(u)=\sum_{\ell=0}^2u^\ell F_{6,r+2\ell}(u^3),\qquad r=0,1}
{The parent color identifies even or odd original indices. The child address is $r+2\ell$, so the even parent has children zero, two, four and the odd parent has children one, three, five. Each child row lists three actual Fibonacci coefficients.}
{Inside the even parent, coefficient $55$ occurs at internal index five: $F_{10}=55$. Splitting $5=3\cdot1+2$ sends it to child four at index one. Its reconstructed exponent is $6\cdot1+4=10$.}
{A child address alone is insufficient. Carry its internal index and the factor $u^\ell$ through the refinement.}
{OEIS A000045; Paper I parent--child refinement. The rows use exact values through $F_{17}$; the identity follows from splitting the parent index $m=3j+\ell$.}

\lesson{POLYPHASE / 07}{Different refinement orders reach the same coefficient}{P07_refinement_routes}
{17=2(3\cdot2+2)+1=3(2\cdot2+1)+2=6\cdot2+5}
{Follow the left path through modulus two and the right path through modulus three. The intermediate residues differ. Both paths arrive at the same modulus-six address, with the same internal index and the same Fibonacci coefficient.}
{Left: $17=2\cdot8+1$, then $8=3\cdot2+2$, giving child $1+2\cdot2=5$. Right: $17=3\cdot5+2$, then $5=2\cdot2+1$, giving child $2+3\cdot1=5$. Both retain $F_{17}=1597$.}
{Refinement orders agree after the address permutation is accounted for. Their intermediate labels need not agree.}
{Paper I refinement identity, illustrated with OEIS A000045. The two routes are exact Euclidean division; the full coefficient check uses indices zero through seventeen.}

\lesson{POLYPHASE / 08}{Three carries produce one block delay}{P08_three_carries}
{C_3(u)=\begin{pmatrix}0&0&u\\1&0&0\\0&1&0\end{pmatrix},\qquad C_3(u)^3=uI_3}
{Columns show successive multiplication by $z$ for the finite Fibonacci polynomial $P=z+z^2+2z^3+3z^4+5z^5$. Rows keep the three residue addresses. After three carries, every packet has returned to its row and acquired exactly one block delay.}
{Initially the packets are $(2u,1+3u,1+5u)$. One carry gives $(u+5u^2,2u,1+3u)$. Three carries give $(2u^2,u+3u^2,u+5u^2)$, which synthesize to $z^3P$.}
{Wrapping adds a power of $u$. It does not silently erase the coefficient at the end of a displayed prefix.}
{Paper I intertwining law. The matrix identity is algebraic over a polynomial ring; the displayed six-coefficient example retains its full output support.}

\lesson{POLYPHASE / 09}{Polynomial multiplication has visible residue carries}{P09_product_carries}
{h_t(u)=\sum_{\substack{0\le r,s<M\\r+s\equiv t\!\!\pmod M}}u^{\lfloor(r+s)/M\rfloor}f_r(u)g_s(u)}
{Each cell multiplies one residue coefficient from $f=1+2z+3z^2$ by one from $g=4+5z+6z^2$. Its color records the output residue. The cells beyond the diagonal wrap to a lower residue and carry a factor of $u$.}
{At modulus three, the product packets are $h_0=4+27u$, $h_1=13+18u$, and $h_2=28$. They synthesize to $4+13z+28z^2+27z^3+18z^4$, exactly the ordinary polynomial product.}
{Residue convolution must retain carry factors. Multiplying only equal-address channels omits cross terms.}
{Worked finite-polynomial example over the integers. The general formula follows from $r+s=M\lfloor(r+s)/M\rfloor+(r+s\bmod M)$ in a commutative coefficient ring.}

\lesson{POLYPHASE / 10}{An exponential generating function changes the weights}{P10_ogf_egf}
{\sum_{n\ge0}\frac{a_nz^n}{n!}=\sum_{r=0}^{M-1}z^r\sum_{m\ge0}\frac{a_{Mm+r}}{(Mm+r)!}(z^M)^m}
{The same Fibonacci values appear on both sides of the table. Ordinary generating-function coefficients retain the integers. Exponential generating-function coefficients divide by the factorial of the original index before channel routing. The denominator does not reset inside a channel.}
{The even coefficient $F_6=8$ becomes $8/6!=1/90$ in the exponential series, at channel index three. Dividing by $3!$ instead would give $4/3$, a different series. The odd coefficient $F_7=13$ becomes $13/5040$.}
{State whether a channel holds sequence values or normalized series coefficients. This choice changes every operator using those coefficients.}
{OEIS A000045 supplies the values. Exact rational normalization; the reconstruction formula is a coefficient identity over a characteristic-zero field. It is not a claim that the channel itself is an EGF in $u$.}

\lesson{POLYPHASE / 11}{Parity separates secant and tangent counts}{P11_secant_tangent}
{\sec z=\sum_{m\ge0}\frac{A_{2m}z^{2m}}{(2m)!},\qquad \tan z=\sum_{m\ge0}\frac{A_{2m+1}z^{2m+1}}{(2m+1)!}}
{The combined up/down stream $A_n$ splits into a secant row and a tangent row. The figure displays the integer counts, with the factorial denominator underneath each value. Thus the parity split is applied to the coefficients of $\sec z+\tan z$ with their original-index normalization intact.}
{At internal index two, the secant term is $5z^4/4!$ and the tangent term is $16z^5/5!$. The tangent row starts $1,2,16,272,7936$ at exponents $1,3,5,7,9$. These are positive counts, not signed Euler numbers.}
{The tangent OEIS sequence starts at native index one. Here its first term is explicitly reindexed to channel index $m=0$.}
{OEIS A000111, A000364, A000182. The finite check independently divides the formal sine and cosine series and compares counts through exponent eleven.}

\lesson{POLYPHASE / 12}{Thue--Morse refines by copying and complementing}{P12_thue_morse}
{t(2m)=t(m),\quad t(2m+1)=1-t(m),\quad T(z)=(1-z)T(z^2)+\frac{z}{1-z^2}}
{Each parent bit creates two children. Its even child copies the parent, while its odd child complements it. The two-row child bank reconstructs the first sixteen bits when interleaved. No statistical fitting is involved: the binary digit rule determines every entry.}
{For $m=5$, the binary expansion is $101$ and its digit sum is two, so $t(5)=0$. Its children are $t(10)=0$ and $t(11)=1$. In the generating identity, the complement contributes the constant stream $1/(1-z^2)$.}
{A refinement operator may be affine. Omitting the constant term changes the odd channel.}
{OEIS A010060, offset zero. $t(n)$ is the binary digit sum modulo two. The universal recurrence follows by appending a binary zero or one; exact checks cover $0\le n<256$.}

\lesson{POLYPHASE / 13}{Binary weight keeps the affine increment}{P13_binary_weight}
{s_2(2m)=s_2(m),\quad s_2(2m+1)=s_2(m)+1,\quad W(z)=(1+z)W(z^2)+\frac{z}{1-z^2}}
{The left bars show the binary weights of eight parents. On the right, even children preserve each height and odd children add one. These are integer digit counts. Reducing their values modulo two recovers the previous lesson's Thue--Morse bits.}
{The parent $m=7$ has binary form $111$, so its weight is three. Appending zero gives $14=1110$ with weight three. Appending one gives $15=1111$ with weight four. The odd increment is visible in every pair.}
{An affine update becomes linear on the augmented state $(s_2(m),1)^T$. The odd update matrix is $\left(\begin{smallmatrix}1&1\\0&1\end{smallmatrix}\right)$.}
{OEIS A000120, offset zero. The recurrence and generating identity follow from binary appending. Exact checks cover $0\le n<256$; the bars display only parents zero through seven.}

\lesson{POLYPHASE / 14}{A Haar bank preserves average and detail together}{P14_haar_bank}
{\binom{a}{d}=\frac1{\sqrt2}\begin{pmatrix}1&1\\1&-1\end{pmatrix}\binom{x}{y},\qquad x^2+y^2=a^2+d^2}
{Group the four input amplitudes into adjacent pairs. The right bars show the normalized average and detail channels for each pair. A negative detail is meaningful: it records which member of the pair is larger. Keeping both channels permits the inverse transform.}
{For $(1,3,2,0)$, the averages are $(2\sqrt2,\sqrt2)$ and the details are $(-\sqrt2,\sqrt2)$. Their energies are ten and four, totaling the original energy fourteen. Inverting the first pair gives $x=(a+d)/\sqrt2=1$ and $y=(a-d)/\sqrt2=3$.}
{These are normalized filtering taps on paired amplitudes. They are distinct from recurrence coefficients or an unnormalized arithmetic mean.}
{Exact Haar two-channel teaching example over $\mathbb R$; amplitudes lie in $\mathbb Q(\sqrt2)$. The plotted bars are numerical renderings of the stated exact algebraic values.}

\lesson{POLYPHASE / 15}{One retained channel cannot identify the original}{P15_discarded_channel}
{\mathcal R_{2,0}(1+2z+3z^2+4z^3)=\mathcal R_{2,0}(1+9z+3z^2-7z^3)=1+3u}
{The two input polynomials share their even coefficients but differ at both odd positions. The right panel overlays the identical retained channel. The missing odd channel prevents the two inputs from being distinguished after projection.}
{Their difference is $7z-11z^3$, whose even channel is zero. Full two-channel routing would retain odd channels $2+4u$ and $9-7u$ and would reconstruct each input exactly. Discarding those channels loses this distinction.}
{Exact routing is reversible only while the complete channel state is retained. A projection needs a separate recovery hypothesis.}
{Explicit finite counterexample over the integers. This proves noninjectivity of retaining only the even channel, rather than estimating reconstruction quality on a sample.}

\lesson{POLYPHASE / 16}{A complete operation agrees along both routes}{P16_commuting_operation}
{\mathcal R_2((1+z)P)=\begin{pmatrix}1&u\\1&1\end{pmatrix}\mathcal R_2(P),\qquad \mathcal S_2\mathcal R_2=I}
{Follow either side of the square. The left route operates on the scalar polynomial first and then splits it. The right route splits first and applies the channel matrix. The two bottom states encode exactly the same seven output coefficients.}
{For $P=z+z^2+2z^3+3z^4+5z^5$, the input channels are $E=u+3u^2$, $O=1+2u+5u^2$. The output channels are $E'=2u+5u^2+5u^3$, $O'=1+3u+8u^2$. Their synthesis ends with $5z^6$, retained from the finite input boundary.}
{A complete contract declares the input support, channel convention, operator, and output support. Equality of those states is the check.}
{Paper I intertwining, worked on a finite prefix of OEIS A000045. The universal matrix identity follows from the two-phase carry; exact integer checks retain the whole output polynomial.}

\startvolume{Tap atlas}{Finite weights and exact execution}

\lesson{TAPS / T01}{Read a tap as a delayed contribution}{T01_taps}
{y[n]=\sum_{k=0}^{L-1}h[k]x[n-k],\qquad h=(1,2,1),\quad x[n]=0\ (n<0)}
{The stem height is a weight, and its horizontal position is a delay. The right panel applies the three weights to three different input samples before adding. A tap specifies which delayed value contributes and with what weight.}
{At $n=5$, the three inputs are $x[5],x[4],x[3]$. Their contribution is $x[5]+2x[4]+x[3]$. At $n=0$, the two delayed inputs are zero under our causal convention.}
{Keep the weight, the delay, and the input sample separate.}
{New worked FIR example. Integer indices $n\ge0$; every input is zero at negative indices. The finite tap list has length $L=3$. This convention applies throughout the tap atlas.}

\lesson{TAPS / T02}{Execute the taps on Fibonacci numbers}{T02_taps}
{F_0=0,\ F_1=1,\qquad y[n]=F_n+2F_{n-1}+F_{n-2}}
{Read the two rows at the same fine index. The input follows its recurrence, while the output comes from a separate weighted convolution. The output is larger because two nearby values contribute in addition to the current value.}
{At $n=5$, the result is $5+2\cdot3+2=13$. At $n=11$, it is $89+2\cdot55+34=233$. The first twelve displayed outputs are computed with integers, including the zero-extended initial positions.}
{A sequence recurrence and an applied tap filter are different operations.}
{Fibonacci convention $F_0=0,F_1=1$; OEIS A000045 as verified in the sequence atlas. New tap example. Indices $0$--$11$; exact inputs and outputs are recorded in \texttt{evidence/tap\_checks.json}.}

\lesson{TAPS / T03}{Turn three taps into a two-phase matrix}{T03_taps}
{\begin{gathered}\begin{pmatrix}Y_0\\Y_1\end{pmatrix}
 =\begin{pmatrix}1+u&2u\\2&1+u\end{pmatrix}
 \begin{pmatrix}E\\O\end{pmatrix},\\
 E=\sum_{m\ge0}x[2m]u^m,\quad O=\sum_{m\ge0}x[2m+1]u^m\end{gathered}}
{Each arrow is one matrix entry. The same-phase contributions use $1+u$, and the cross-phase contributions differ by a coarse delay. Here $u$ shifts an entire pair of fine samples, so it must not be read as a single fine-sample delay.}
{The even output is $E_m+2O_{m-1}+E_{m-1}$. The odd output is $O_m+2E_m+O_{m-1}$. Summing each channel as an ordinary generating function gives the displayed matrix exactly.}
{Polyphase representation reorganizes the original convolution.}
{PSC Paper V, Theorem 4.2 and Corollary 4.5, specialized to $M=2$ and $h=(1,2,1)$; new exact derivation. Rows are output phases, columns are input phases. Causal zero extension fixes the initial positions.}

\lesson{TAPS / T04}{Locate the delay when a tap wraps a block}{T04_taps}
{y[2m]=E_m+2O_{m-1}+E_{m-1},\qquad
 y[2m+1]=O_m+2E_m+O_{m-1}}
{Follow the middle tap in both panels. Subtracting one from the even fine index reaches the odd phase of the previous block. Subtracting one from the odd fine index reaches the even phase of the current block. That single index calculation explains the asymmetry.}
{For $m=2$, $y[4]$ reads $x[4],x[3],x[2]$, whereas $y[5]$ reads $x[5],x[4],x[3]$. The first middle contribution has a factor $u$ in its generating function; the second does not.}
{Assign a coarse delay only after solving the fine-index equation.}
{PSC Paper V, Theorem 4.2. New per-tap expansion of the matrix in T03. The labels $E_m,O_m$ denote phase coefficients, while $E(u),O(u)$ denote their OGFs.}

\lesson{TAPS / T05}{Route a finite filter through four phases}{T05_taps}
{P_{rs}(u)=\sum_{\substack{k\ge0\\0\le Mk+r-s<L}}h[Mk+r-s]u^k,
 \qquad 0\le r,s<M}
{Rows are output phases and columns are input phases. For the five taps $(1,2,3,4,5)$, a same-phase contribution combines tap zero with tap four. Warm cells cross the block boundary and therefore carry one coarse delay.}
{At $M=4$, the first row is $(1+5u,4u,3u,2u)$. For example, $y[4m]$ receives $4x[4m-3]=4X_{1,m-1}$ from tap three. That contribution belongs to column one and has coefficient $4u$.}
{The index equation supplies every matrix entry for any finite tap length.}
{PSC Paper V, Theorem 4.2, Corollary 4.5, and Example 4.6. The plotted matrix uses $M=4,L=5$. Exact routing also checked for moduli $2$--$6$ and tap lengths $1$--$9$; these finite tests verify the implementation.}

\lesson{TAPS / T06}{Compose filters before or after phase splitting}{T06_taps}
{h=(1,2,1),\quad g=(1,-1),\quad g*h=(1,1,-1,-1),
 \qquad P_2(g)P_2(h)=P_2(g*h)}
{The combined tap list is the convolution of the two finite lists. The bottom identity describes the same composition after splitting into even and odd channels. Matrix multiplication retains the delays that are easy to lose in a verbal description.}
{The combined fine-delay polynomial is $(1-z)(1+2z+z^2)=1+z-z^2-z^3$. Its two-phase matrix is $\bigl(\begin{smallmatrix}1-u&u-u^2\\1-u&1-u\end{smallmatrix}\bigr)$, obtained independently by routing the four combined taps.}
{Composition is represented faithfully in both coefficient and phase coordinates.}
{PSC Paper V, Proposition 4.4 and Corollary 4.5; new signed-filter example. The rightmost matrix acts first. Exact polynomial convolution and exact matrix multiplication agree in the check ledger.}

\lesson{TAPS / T07}{Use signed taps as finite differences}{T07_taps}
{\Delta x[n]=x[n]-x[n-1],\qquad
 \Delta(n^2)=2n-1\ (n\ge1),\quad \Delta^2(n^2)=2\ (n\ge2)}
{The first signed filter turns a quadratic sequence into an increasing linear sequence. Applying it again gives a constant after the initial boundary. The shaded positions show why a polynomial identity on unrestricted integer arguments is not automatically the zero-extended causal identity at the start.}
{At $n=3$, $9-4=5$. A second difference gives $5-3=2$. At $n=1$, the causal second difference is $1$, because the missing negative-index input was defined to be zero.}
{State the boundary before calling a filter an annihilator.}
{New exact finite-difference example, tap list $(1,-1)$, input $x[n]=n^2$ for $n\ge0$. Indices $0$--$10$ checked exactly. This lesson does not claim a finite constant-coefficient annihilator for arbitrary sequences.}

\lesson{TAPS / T08}{Let recurrence taps expose initial data}{T08_taps}
{(1-z-z^2)F(z)=z,\qquad
 F(z)=\sum_{n\ge0}F_nz^n,\qquad ((1,-1,-1)*F)[n]=\delta_{n,1}}
{The negative taps subtract the two previous Fibonacci values. Every recurrence position cancels, but the initial value $F_1=1$ survives as one impulse. The generating-function numerator records precisely that surviving initial information.}
{At $n=4$, $3-2-1=0$. At $n=1$, $1-0-0=1$. Thus the filter annihilates the recurrence tail while leaving the initial impulse visible. Here $\delta_{n,j}$ is one at $n=j$ and zero otherwise.}
{An initial numerator is part of the exact recurrence representation.}
{Fibonacci recurrence with causal zero extension; new tap presentation. Twenty-four coefficients checked exactly. The polynomial $1-z-z^2$ uses the fine-delay variable $z$, unlike the coarse variable $u$ in T09.}

\lesson{TAPS / T09}{Give the two Fibonacci phases different numerators}{T09_taps}
{E(u)=\frac{u}{1-3u+u^2},\qquad
 O(u)=\frac{1-u}{1-3u+u^2},\qquad E_m=F_{2m},\ O_m=F_{2m+1}}
{The two phase channels share a recurrence denominator, but their initial impulses differ. Applying the taps $(1,-3,1)$ gives one impulse in the even channel and two signed impulses in the odd channel. Those residuals are the two displayed numerators.}
{The first even terms are $0,1,3,8$, giving residuals $0,1,0,0$. The first odd terms are $1,2,5,13$, giving $1,-1,0,0$. The coarse recurrence holds from $m=2$ onward in both channels.}
{Shared denominator does not mean identical phase initialization.}
{Fibonacci phase extraction and exact recurrence calculation. Twenty-four coefficients in each channel checked. Residual notation: $\delta_1$ and $\delta_0-\delta_1$ on the coarse index.}

\lesson{TAPS / T10}{Separate finite taps from an infinite reciprocal response}{T10_taps}
{H(u)=\frac1{1-3u+u^2}=1+3u+8u^2+21u^3+\cdots,
 \qquad E=uH,\quad O=(1-u)H}
{The long row is an infinite impulse response. The two boxes below it are finite numerator tap lists acting on that response. The denominator taps are finite too, but taking their reciprocal creates a response with infinitely many nonzero coefficients.}
{The even numerator $(0,1)$ shifts the response to $0,1,3,8,\ldots$. The odd numerator $(1,-1)$ subtracts successive response coefficients to give $1,2,5,13,\ldots$. Applying $(1,-3,1)$ directly to $H$ leaves $\delta_0$.}
{Finite numerator and denominator polynomials do not make their reciprocal FIR.}
{New formal-OGF reconstruction of T09. Twenty-four coefficients checked exactly. The reciprocal identity is formal around $u=0$; the plate makes no claim about convergence at every complex value of $u$.}

\lesson{TAPS / T11}{Read frequency response from the tap polynomial}{T11_taps}
{H(e^{-i\omega})=1+2e^{-i\omega}+e^{-2i\omega},\qquad
 |H(e^{-i\omega})|=4\cos^2(\omega/2)}
{The tap polynomial becomes a frequency response by substituting a unit complex delay. Its value at zero frequency is the sum of the weights. The endpoint at $\pi$ cancels because adjacent samples alternate signs there.}
{A constant input is multiplied by $1+2+1=4$ away from its causal starting boundary. At $\omega=\pi$, the three contributions are $1,-2,1$ and sum to zero. The plotted curve agrees with the exact magnitude formula on the stated grid.}
{A frequency curve describes these taps; it is not a performance guarantee.}
{New FIR frequency example. NumPy complex128 on 801 samples of $[0,\pi]$. The endpoint and DC identities are algebraic; the sampled error is recorded separately. No ideal anti-aliasing claim is made.}

\lesson{TAPS / T12}{Expose the delay carried by symmetric taps}{T12_taps}
{H(e^{-i\omega})=e^{-i\omega}(2+2\cos\omega),\qquad
 \arg H(e^{-i\omega})=-\omega\quad(-\pi<\omega<\pi)}
{The symmetric taps factor into one sample of delay and a nonnegative real amplitude. Removing the delay leaves the second panel. At the endpoint zeros the complex number vanishes, so its phase is undefined rather than a meaningful continuation of the plotted straight line.}
{At $\omega=\pi/2$, the amplitude is $2$ and the delayed response is $-2i$. More generally, a symmetric list of length $L$ factors a delay $(L-1)/2$; the remaining real factor may change sign, introducing phase jumps.}
{Linear phase needs its sign and zero conventions stated.}
{New symmetric FIR example. The displayed formula is exact. The numerical phase check uses 801 interior samples and excludes $\pm\pi$. A phase value supplied by floating-point arithmetic at a zero is not used as evidence.}

\lesson{TAPS / T13}{Reconstruct exactly with Haar analysis and synthesis}{T13_taps}
{A=\frac1{\sqrt2}\begin{pmatrix}1&1\\1&-1\end{pmatrix},\qquad
 \begin{pmatrix}c_m\\d_m\end{pmatrix}
 =A\begin{pmatrix}x[2m]\\x[2m+1]\end{pmatrix},\qquad A^{\mathsf T}A=I}
{Analysis combines one nonoverlapping pair into a normalized sum and difference. Synthesis reverses that change of coordinates. This plate fixes the pair convention explicitly, avoiding the extra shifts that appear under other convolution-and-decimation conventions.}
{The pair $(3,5)$ gives $(c,d)=(4\sqrt2,-\sqrt2)$. Reconstruction yields $((c+d)/\sqrt2,(c-d)/\sqrt2)=(3,5)$. Both the normalization and the difference sign are necessary.}
{Perfect reconstruction is an operator identity with a fixed convention.}
{Standard Haar pair transform, independently verified by direct matrix algebra. Exact sum/difference witnesses cover all 121 integer pairs in $[-5,5]^2$, with the common square-root factor canceled symbolically.}

\lesson{TAPS / T14}{Reverse signed prediction and update steps}{T14_taps}
{D=O-PE,\quad S=E+UD,\qquad E=S-UD,\quad O=D+PE}
{Forward lifting first predicts the odd channel from the even channel, then updates the even channel using the residual. Reverse lifting undoes those two steps in the opposite order. Neither step needs an inverse filter for $P$ or $U$.}
{Use the signed causal taps $P=1-u$ and $U=-u$. The combined matrix is $\bigl(\begin{smallmatrix}1+u-u^2&-u\\-1+u&1\end{smallmatrix}\bigr)$ and has determinant one. Any finite predict and update polynomials admit the displayed algebraic reversal.}
{Invertibility comes from the triangular steps, not from a special predictor.}
{New exact lifting example over rational coefficients. Negative coarse indices are zero in both directions. Sixteen signed filter-pair witnesses checked with integers and fractions. Streaming use must retain the histories needed by both filters.}

\lesson{TAPS / T15}{Resample at rate three over two by exact routing}{T15_taps}
{\begin{gathered}y[m]=\sum_n h[2m-3n]x[n]
 =\sum_{k\ge0}h_{r_m}[k]x[\ell_m-k],\\
 r_m=2m\bmod3,\quad \ell_m=\lfloor2m/3\rfloor\end{gathered}}
{First insert two zeros between input samples, convolve on the fine grid, and keep every second result. The table expresses those same kept values through three tap phases: $h_r[k]=h[3k+r]$. The phase selector changes with the output index.}
{For $x=(1,2,3,4)$ and $h=(1,2,1)$, the selected phases are $0,2,1,0,2,1$. The six outputs are $(1,1,4,3,3,8)$. Each method includes the full finite convolution tail.}
{Rational rate conversion can be implemented without materializing inserted zeros.}
{New small-filter $3/2$ example. Exact zero insertion, convolution, downsampling, and phase routing agree. The filter demonstrates indexing only; no passband, stopband, or alias-suppression specification is claimed.}

\lesson{TAPS / T16}{Carry the needed state across chunk boundaries}{T16_taps}
{\text{direct-form FIR state: }L-1\text{ previous input samples},
 \qquad y_{\mathrm{chunks}}[n]=y_{\mathrm{whole}}[n]}
{The dashed lines mark chunk boundaries, not restarts of the filter. The lower boxes hold the two input samples needed by the next chunk. The shaded tail is produced by an explicit two-zero flush after the final real sample.}
{For $(1,2,1)$, after twelve Fibonacci samples the saved state is $(55,89)$, oldest first. Feeding zero twice produces tail outputs $233$ and $89$. Resetting state at an interior chunk boundary would lose delayed contributions.}
{Boundary handling is part of the mathematical execution contract.}
{New direct-form FIR streaming example. Twenty signed-filter cases cover lengths $1$--$5$, four chunk partitions, and complete flushed tails. The $L-1$ rule describes this implementation, rather than a universal lower bound on every filter realization.}

\startvolume{Operator bridges}{Seeds, responses, and coefficient transforms}

\lesson{BRIDGE / 01}{Seeds choose which modes survive}{B01_surviving_modes}
{a_{n+2}=3a_{n+1}-2a_n,\qquad A(z)=\frac{a_0+(a_1-3a_0)z}{(1-z)(1-2z)}}
{Every row uses the same recurrence denominator. The seeds change its numerator, which can cancel a denominator factor. Filled marks are poles that remain after reducing the rational function; crosses mark factors that cancel. The two candidate modes are $1$ and $2^n$.}
{Seeds $(1,2)$ give numerator $1-z$, leaving $A(z)=1/(1-2z)$ and coefficients $2^n$. Seeds $(1,1)$ instead cancel $1-2z$, leaving the constant sequence. Seeds $(0,1)$ retain both modes and give $2^n-1$.}
{A recurrence denominator lists candidate modes. The numerator determines which are present in the specified sequence.}
{New elementary rational-series example, $n\ge0$. Numerator residuals and all three mode formulas are checked through index eighteen. Cancellations are exact; this is not a robustness claim for perturbed seeds or hidden states.}

\lesson{BRIDGE / 02}{Compile taps against an exact recurrence}{B02_recurrence_tap_reduction}
{\begin{gathered}D=1-z-z^2,\quad F=z/D,\quad h=(1+z)^2=(2+z)-D,\\ hF=(2+z)F-z\end{gathered}}
{The top panels compare the original three taps with a shorter list. The table shows why the shorter list needs an initial correction: its output differs at index one. Subtracting that impulse makes every coefficient agree, including the causal boundary. Here $F_n=0$ for $n<0$.}
{The reduced formula is $y_n=2F_n+F_{n-1}-\delta_{n,1}$. At $n=1$, it gives $2-1=1$; at $n=5$, it gives $2\cdot5+3=13$. If $DF=N$ and polynomial division gives $h=qD+r$, then $hF=rF+qN$.}
{The recurrence can simplify an operator while its numerator carries the remaining boundary contribution.}
{OEIS A000045, $F_0=0,F_1=1$, with new exact polynomial algebra. The coefficient check runs through index eighteen; finite-input tails are checked with $N=DF$, including terminal residuals. The reduction is for this recurrence input; it does not replace the original filter on arbitrary signals.}

\lesson{BRIDGE / 03}{A formal reciprocal need not be a stable filter}{B03_reciprocal_stability}
{H_r(z)=\frac1{1-rz}=\sum_{n\ge0}r^nz^n,\qquad |z|<1/|r|\quad(r\ne0)}
{The disks show where each causal series converges. The dashed unit circle lies inside the disk for $r=1/2$ and outside it for $r=2$. The right panel applies both responses to the same bounded input $x[n]=1$ and uses a logarithmic output axis.}
{For $r=1/2$, $y_n=2-2^{-n}$ tends to two, and every bounded input satisfies $|y_n|\le2\|x\|_\infty$. For $r=2$, $y_n=2^{n+1}-1$ grows without bound even though the input stays one.}
{A formal-series identity, an analytic convergence domain, and bounded-input stability have distinct contracts.}
{New causal geometric-response examples, $n\ge0$, with zero input at negative indices. Exact rational convolutions checked through index eighteen. The unstable rational expression can be evaluated on the unit circle, but its causal Fourier series there does not converge.}

\lesson{BRIDGE / 04}{Coefficient division becomes an area average}{B04_catalan_integration}
{B(z)=\frac1{\sqrt{1-4z}},\qquad C(z)=\frac1z\int_0^z B(t)\,dt=\frac{1-\sqrt{1-4z}}{2z}}
{The shaded area under the central-binomial generating function is divided by the interval length. The dashed line is that average, the Catalan generating function at $z=1/8$. The table shows the coefficient operation: integrating raises the exponent and divides by $n+1$; dividing by $z$ restores its address.}
{At index four, $70z^4$ integrates to $14z^5$, then becomes $14z^4$ after division by $z$. At $z=1/8$, the area is $(2-\sqrt2)/4$, so its average is $4-2\sqrt2$.}
{An index-dependent coefficient normalization can come from a familiar continuous operator.}
{OEIS A000984 and A000108; new elementary integral derivation. The square-root branch equals one at zero, $C(0)=1$ by removal, and the analytic identities hold on $|z|<1/4$. Coefficient identities also hold over formal rational series; exact checks cover indices zero through eighteen.}

\startvolume{Operator contracts}{Normalization, translation, and boundaries}

\lesson{OPERATOR / 01}{Transport the operator with the normalization}{O01_factorial_transport}
{\begin{gathered}(Na)_n=\frac{a_n}{n!},\qquad (Hb)_n=\sum_{k=0}^{\min(K,n)}h_kb_{n-k},\\
(N^{-1}HN a)_n=\sum_{k=0}^{\min(K,n)}h_k(n)_k a_{n-k},\qquad (n)_k=\frac{n!}{(n-k)!}\end{gathered}}
{The same stationary taps $(1,-1)$ are applied in two different coefficient coordinates. The top pair uses integer counts; the middle pair uses Taylor coefficients. Returning the latter output to count coordinates gives the bottom row. Its operator has index-dependent weights, rather than the original stationary taps.}
{Start with $a_n=1$ for all $n\ge0$, so $b_n=1/n!$ and the EGF is $e^z$. Filtering the counts gives $(1,0,0,\ldots)$ and EGF $1$. Filtering the Taylor coefficients gives $(1-z)e^z$. Its count coordinates are $1-n$, beginning $1,0,-1,-2,\ldots$.}
{A stationary Taylor FIR pulls back to counts as $N^{-1}HN$. Applying the same stationary weights in both coordinates is a different operation.}
{Exact formal coefficient identity over $\mathbb Q$, $\mathbb R$, or $\mathbb C$; $n,K\ge0$ are integers and $h_0,\ldots,h_K$ are finite taps. Counts transported in the opposite direction use $NHN^{-1}$ with weights $h_k/(n)_k$ for $n\ge k$. Prefixes are not finite input cutoffs.}

\lesson{OPERATOR / 02}{Bernoulli resolves a discrete antiderivative}{O02_bernoulli_antiderivative}
{\begin{gathered}B_4(x)=x^4-2x^3+x^2-\frac1{30},\qquad B_4(x+1)-B_4(x)=4x^3,\\
\frac{B_4(N)-B_4(0)}4=\sum_{k=0}^{N-1}k^3\qquad(N\ge0)\end{gathered}}
{Each bar is a forward translation gap divided by four. The taps compare the polynomial at arguments $k+1$ and $k$; they do not shift its order or its Taylor-coefficient index. Adding the five visible gaps cancels all interior endpoint values. The two remaining Bernoulli endpoints determine the sum.}
{For $N=5$, the gaps are $0,1,8,27,64$, totaling $100$. Also $B_4(5)=11999/30$ and $B_4(0)=-1/30$, so their difference is $400$. Division by four recovers exactly the same sum.}
{Translation taps can turn a special function into an exact discrete primitive. Its additive constant disappears under differencing.}
{Bernoulli-polynomial normalization from Paper I and sequence lesson S11, with $B_1=-1/2$. Expanding the stated polynomial proves its difference identity for every real or complex $x$; telescoping proves the integer sum. Exact sample checks cover $0\le N\le50$.}

\lesson{OPERATOR / 03}{The boundary changes what an inverse can recover}{O03_boundary_inversion}
{D_{\rm per}x[n]=x[n]-x[(n-1)\bmod4],\quad
\ker D_{\rm per}=\langle\boldsymbol1\rangle,\quad
\operatorname{im}D_{\rm per}=\left\{y:\sum_{n=0}^3y[n]=0\right\}}
{The periodic cycle places two possible amplitudes at every node. Both inputs give the same edge differences because adding a constant changes no periodic difference. The causal panel uses the same taps $(1,-1)$ with zero extension. Its initial output and flushed tail retain the boundary information.}
{The four-value inputs $(1,3,2,0)$ and $(6,8,7,5)$ both have periodic output $(1,2,-1,-2)$. Their causal outputs, including one zero flush, are $(1,2,-1,-2,0)$ and $(6,2,-1,-2,-5)$. In the periodic problem, fixing $x[0]=1$ reconstructs the first input uniquely.}
{A periodic inverse requires zero-sum output and one initial value. Causal summation instead recovers the zero-extended input from its full difference stream.}
{Exact four-periodic linear algebra over $\mathbb Q$ or $\mathbb R$. In the causal panel each four-value vector is zero outside indices $0$--$3$; adding five affects those stored values only. Equal periodic differences force a constant difference, while telescoping gives the zero-sum condition.}

\startvolume{Worked capstone}{One record through channels, residuals, and repair}

\lesson{CAPSTONE / 01}{A finite record leaves endpoint impulses}{C01_finite_residual}
{(1-z-z^2)P_N(z)=z-F_Nz^N-F_{N-1}z^{N+1},\qquad N\ge2}
{The record is $x=(0,1,1,2,3,5,8,13)$ at indices zero through seven, with zero elsewhere. Apply the causal taps $(1,-1,-1)$ and retain all ten output positions. The green strip marks the surviving seed. The warm strip marks the two impulses created by ending the record.}
{Write $P_8(z)=\sum_{n=0}^{7}F_nz^n$. Interior coefficients cancel by the Fibonacci recurrence. At index eight, the input has ended, so the output is $0-13-8=-21$. At index nine it is $0-0-13=-13$. Thus the complete output is $z-21z^8-13z^9$.}
{A finite cutoff has a boundary residual even when every interior recurrence step is exact.}
{Direct coefficient derivation, $P_N=\sum_{n=0}^{N-1}F_nz^n$, $F_0=0,F_1=1$. OEIS A000045 supplies the record. Exact finite checks cover $N=2$--$20$; the displayed general identity follows by collecting coefficients, not by extrapolating those checks.}

\lesson{CAPSTONE / 02}{Transfer the whole residual into three channels}{C02_three_phase_residual}
{A(C_3)=\begin{pmatrix}1&-u&-u\\-1&1&-u\\-1&-1&1\end{pmatrix},\qquad
 A(C_3)f=\begin{pmatrix}-13u^3\\1\\-21u^2\end{pmatrix}}
{Split the same eight input values into residues zero, one, and two. Each box holds an ordinary generating polynomial in the internal index. The matrix represents $A(z)=1-z-z^2$. Its factors of $u$ retain the coarse delay when a fine tap crosses a block boundary. Rows are output phases; columns are input phases.}
{The input bank is $f=(2u+8u^2,\,1+3u+13u^2,\,1+5u)^T$. In row zero, subtracting $u$ times the other two channels leaves $-13u^3$. Synthesis of the output bank gives $-13z^9+z-21z^8$, precisely the scalar output in C01.}
{Scalar filtering and full packet filtering agree, including the endpoint impulses.}
{Direct carry calculation, $A(C_3)=I-C_3-C_3^2$, with $u=z^3$ under synthesis. Exact integer polynomial arithmetic checks the complete three-channel result and all ten synthesized output coefficients. No channel is discarded.}

\lesson{CAPSTONE / 03}{An interior perturbation has its own signature}{C03_interior_defect}
{Q=P_8+z^5,\qquad A(z)(Q-P_8)=z^5-z^6-z^7}
{Change only the sample at index five from five to six. The top row shows that changed record. The bottom row shows the change in the filtered output, not the full output. A single input impulse meets three taps and therefore changes three consecutive residual positions. The seed and endpoint residuals stay unchanged.}
{At index five, the perturbed output is $6-3-2=1$. At index six it is $8-6-3=-1$, and at index seven it is $13-8-6=-1$. The full output becomes $z+z^5-z^6-z^7-21z^8-13z^9$.}
{Distinguish an interior model defect from the declared finite-record boundary.}
{Exact one-sample perturbation experiment over the integers. The identity follows by linearity of multiplication by $A(z)$. This example identifies the consequences of its specified change; an arbitrary residual alone does not decide whether data or an assumed recurrence model is wrong.}

\lesson{CAPSTONE / 04}{A discarded channel returns through the next operation}{C04_projection_obstruction}
{\mathcal R_{2,0}(Q-P_8)=0,\qquad
 \mathcal R_{2,0}\bigl(A(z)(Q-P_8)\bigr)=-u^3}
{Switch from C02's three-phase bank to the modulus-two observation retaining even indices only. The two records from C03 have identical observed inputs. After the same filter, their even outputs differ at index six. A contribution from the discarded odd sample has crossed into a retained output.}
{Both retained inputs are $u+3u^2+8u^3$. For $P_8$, the retained output is $-21u^4$. For $Q$, it is $-u^3-21u^4$. The coefficient of $u^3$ changes because the tap $-z$ carries the perturbation at index five to index six.}
{A projection must retain the information needed by the next operator.}
{Explicit finite counterexample, with $u=z^2$ under synthesis: no map $U$ on even input channels can satisfy $\mathcal R_{2,0}(A(z)P)=U(\mathcal R_{2,0}P)$ for every finite integer polynomial $P$. Equal arguments would force equal values. This does not rule out recovery under an additional model.}

\lesson{CAPSTONE / 05}{Repair with a declared recurrence model}{C05_model_repair}
{\begin{gathered}E_m=x[2m],\qquad a=E_0,\quad b=E_1-E_0,\\
 E_2=3E_1-E_0,\qquad E_3=3E_2-E_1\end{gathered}}
{Keep the modulus-two observation, and require the eight-value record to obey $x[n+2]=x[n+1]+x[n]$ for $n=0,\ldots,5$, with unknown seeds $(a,b)$. The first two even observations determine those seeds. The remaining two must pass the displayed consistency checks. The model then supplies the missing odd values.}
{For $E=(0,1,3,8)$, the inferred seeds are $(0,1)$. The checks give $3\cdot1-0=3$ and $3\cdot3-1=8$. Rebuilding from the seeds gives odd values $(1,2,5,13)$ and restores C01's record. The perturbed $Q$ has the same observations but lies outside this recurrence family.}
{A recovery hypothesis changes what the retained coordinates can identify.}
{Direct finite-state derivation: the even values are $(a,a+b,2a+3b,5a+8b)$. The two checks are necessary and sufficient for a unique completion inside this length-eight recurrence family. Exact tests cover 49 integer seed pairs and 196 inconsistent observation cases. No claim of recovery for arbitrary records or noisy observations is made.}

\lesson{CAPSTONE / 06}{Keep the history and the phase clock}{C06_streaming_clock}
{S_n=(x[n-2],x[n-1];\ n\bmod3),\qquad
 y[n]=x[n]-x[n-1]-x[n-2]}
{Return to the modulus-three output bank and execute the same filter in chunks of lengths two, three, and three. Before processing fine index $n$, the state holds the two preceding samples and the phase receiving $y[n]$. Dashed lines are chunk boundaries. The shaded tail comes from two explicit zero inputs after the last real sample.}
{After eight inputs, the history is $(8,13)$ and the next index is eight, phase two. The first zero gives $-21$ in phase two. The next gives $-13$ at index nine, phase zero. Continuing the clock yields output banks $(0,0,0,-13)$, $(1,0,0)$, and $(0,0,-21)$, exactly as in C02.}
{A chunk boundary changes storage, while the filter history and global address continue.}
{Exact direct-form streaming and output routing, with zero history initially. All 128 positive chunk compositions of this eight-value record are checked against whole-record convolution and routing, including the two-zero flush. The stated history describes this implementation; resetting the phase clock changes output addresses even if the arithmetic samples remain correct.}

\clearpage
\pdfbookmark[0]{Sequence, channel, and tap conventions}{addition-conventions}
{\LARGE\bfseries Sequence, channel, and tap conventions\par}
\vspace{0.15in}
The three atlases contain 48 illustrated lessons. They connect sequence definitions with exact coefficient channels and finite operations. Each sequence retains its own starting index and normalization.

\vspace{0.08in}
\begin{tabular}{>{\raggedright\arraybackslash}p{1.35in}>{\raggedright\arraybackslash}p{5in}}
\toprule
Convention & How to read the new lessons\\
\midrule
Ordinary series & $A(z)=\sum_{n\ge0}a_nz^n$. The coefficient is $a_n$.\\[0.09in]
Exponential series & $A(z)=\sum_{n\ge0}a_nz^n/n!$. The sequence value is $n!$ times the coefficient.\\[0.09in]
Residue channels & $A_{M,r}(u)=\sum_{m\ge0}a_{Mm+r}u^m$, with $0\le r<M$. The variable $u$ advances one coarse block.\\[0.09in]
Finite taps & $y_n=\sum_k h_k a_{n-k}$. Negative sample indices are zero in the causal tap examples.\\[0.09in]
Recurrence taps & Applying a finite annihilating kernel can leave an initial-condition residual. Its reciprocal response can have infinitely many terms.\\[0.09in]
Refinement masks & A mask routes a finer-scale representation. Its reconstruction contract is stated separately from an arbitrary convolution filter.\\
\bottomrule
\end{tabular}

\vspace{0.15in}
\textbf{OEIS identifications.} Thirty official OEIS records were read from the export dated 7 October 2026, pinned to commit \texttt{2ce625d68f85e35e26e15c0d9681f69280c095e7}. The source package retains names, offsets, first eight indexed terms, entry revisions, and extraction hashes. The dated records identify the sequences; the displayed transforms have their own derivations and checks.

\vspace{0.08in}
\textbf{Related project source.} The finite convolution and polyphase transfer examples also use \emph{Projector-Resolved Seed Calculus V}, Section 4, Theorem 4.2, Proposition 4.4, Corollary 4.5, and Example 4.6. That source writes coarse delays as $z^{-m}$; these new lessons use $u^m$. The underlying delay operation is the same.

\vspace{0.08in}
\textbf{Reading sequence.} Follow the sequence atlas for definitions and transforms, the polyphase atlas for exact routing, and the tap atlas for execution. These are applications of the coefficient language introduced in the foundations, with additional examples from combinatorics and signal processing.

\vspace{0.08in}
\textbf{Evidence.} The source records, author checks, independent exact computations, and final page inspection are separate records. Finite computations corroborate the implementations in their stated ranges. No new Lean certificate, universal constant-coefficient recurrence for all sequences, or historical novelty claim is made by this addition.

\textbf{Closing bridges and capstone.} Thirteen further lessons connect seeds to surviving modes, carry taps through coefficient normalization, use Bernoulli translation differences, and distinguish causal from periodic inversion. The capstone follows one eight-term Fibonacci record through its complete residual, three channels, a local defect, discarded-channel obstruction, model-constrained repair, and streaming phase state. Their elementary derivations and bounded exact checks are recorded separately.

\clearpage
\pdfbookmark[0]{Notation and source map}{notation-source-map}
{\LARGE\bfseries Notation and source map\par}
\vspace{0.15in}
\begin{tabular}{>{\raggedright\arraybackslash}p{1.3in}>{\raggedright\arraybackslash}p{5.05in}}
\toprule
Notation & Meaning in this guide\\
\midrule
$M,r,m$ & Coefficient modulus, residue address, and internal coefficient index.\\[0.08in]
$a,\beta,k$ & Arithmetic decay, positive scale exponent, and nonnegative integer depth. In the Bernoulli--Hurwitz sections, $a$ instead labels the declared translation argument.\\[0.08in]
$q$ in I--IV & Integer character modulus. Use the actual modulus for the character's zero set.\\[0.08in]
$\alpha,\gamma$ & Positive continuous shapes. Arithmetic shape is $\alpha=(2a-1)/\beta$.\\[0.08in]
$\kappa$ & Positive ladder label; the ordinary shape bridge is $\kappa=\alpha/2$.\\[0.08in]
$q$ in VIII & Real deformation parameter $0<q\le1$, with $q=1$ the ordinary limit.\\[0.08in]
$N$ & Total excitation in VII--VIII; a numerical cutoff elsewhere is labeled explicitly.\\[0.08in]
$\nu,E_k$ & Positive arithmetic moment measure and its moments. A signed companion is distinguished explicitly.\\[0.08in]
$h,\eta$ & Real translation and modulation coordinates in ambiguity formulas.\\[0.08in]
$\mathcal R,P$ & Signed orthogonal recoupling overlap and its entrywise squared measured channel.\\
\bottomrule
\end{tabular}\par
\vspace{0.1in}
In VIII.11, $\alpha,\beta,\gamma,\delta$ are local scalar q-Racah parameters. In VIII.12, $a,b,c,d$ are local scalar Askey--Wilson parameters. Their roles differ from the arithmetic and continuous-shape notation above.

\vspace{0.1in}
The source releases are the eight \emph{Projector Generated Special Functions} archives collected in the first consolidation. Source notes name the paper and claim or construction. The source package includes the guide scripts, computed figure files, per-chapter computation checks, source hashes, correction register, and the 173-record special-function claim index.

The original series contains 431 PDF pages and 212 distinct archived vector figures. This guide uses newly prepared lesson figures and explanatory examples. The eight original releases and the 38-page consolidated manuscript are included in this collection under \texttt{series/} and in its collected reader. The earlier registry is contextual evidence, not an additional numbered volume.

\textbf{Reproduction.} Run the figure scripts, then compile \texttt{main.tex} three times with XeLaTeX. The packaged README gives commands and dependency versions. The computation records identify exact finite checks, floating-point residuals, quadrature tolerances, and declared truncation bounds. Visual inspection and mathematical review are recorded separately.
\end{document}
